\documentclass[11pt]{article}

\usepackage[T1]{fontenc}
\usepackage[utf8]{inputenc}
\usepackage{amsmath,amssymb,amsthm,mathtools}
\usepackage{mathrsfs}
\usepackage{booktabs,longtable,array}
\usepackage{enumitem}
\usepackage[margin=1in]{geometry}
\usepackage{xcolor}
\usepackage[hidelinks]{hyperref}

\hypersetup{
  pdftitle={Renewal Navier--Stokes Stability Research Notes V26},
  pdfauthor={Research continuation notes}
}

\newcommand{\Torus}{\mathbb T}
\newcommand{\R}{\mathbb R}
\newcommand{\dd}{\,\mathrm d}
\newcommand{\PP}{\mathbb P}
\newcommand{\BNS}{\mathcal B}
\newcommand{\PROVED}{\textcolor{green!45!black}{\textsf{PROVED}}}
\newcommand{\IMPORTED}{\textcolor{blue!60!black}{\textsf{IMPORTED}}}
\newcommand{\OPEN}{\textcolor{red!70!black}{\textsf{OPEN}}}
\newcommand{\CONDITIONAL}{\textcolor{orange!75!black}{\textsf{CONDITIONAL}}}

\theoremstyle{plain}
\newtheorem{theorem}{Theorem}[section]
\newtheorem{proposition}[theorem]{Proposition}
\newtheorem{lemma}[theorem]{Lemma}
\newtheorem{corollary}[theorem]{Corollary}

\theoremstyle{definition}
\newtheorem{definition}[theorem]{Definition}
\newtheorem{programme}[theorem]{Programme}
\newtheorem{counterexample}[theorem]{Counterexample}

\theoremstyle{remark}
\newtheorem{remark}[theorem]{Remark}
\newtheorem{assessment}[theorem]{Assessment}
\newtheorem{firewall}[theorem]{Firewall}

\title{\textbf{Renewal Navier--Stokes Stability Research Notes V26}\\[0.35em]
\large Fourth cycle, V26: the pointwise rank-one norms of the kernel derivatives}
\author{}
\date{7 September 2026}

\begin{document}
\maketitle

\begin{abstract}
This file opens the fourth cycle of an append-only research series
whose objective is a complete stability/global-regularity proof for the
three-dimensional incompressible Navier--Stokes equations on the torus.
The first three cycles, one hundred versions each, are retained as
legacy files whose digests are recorded below; nothing in them is
repeated here beyond the audited statements of their principal
results, the list of external theorems used, the corrections made,
and the open gates.  The third cycle's principal explicit result is
a certified floor on the rescaled velocity bound at a homogeneous
type-I singular point, \(C_\infty\ge0.1815\,\nu^{1/2}\); it is not an exclusion.
The file then makes the first acquisition of the fourth cycle.  No
global regularity theorem is claimed.  That problem remains open.
\end{abstract}

\tableofcontents

% =============================================================================
\section{Context and rules of the series}
\label{sec:c4v1-context}
% =============================================================================

\subsection{What is being attempted}

Let \(u\) be a smooth, mean-zero, divergence-free solution on
\([0,T_*)\times\Torus^3\) (torus of side \(2\pi\)) of
\begin{equation}
 \partial_tu+\nu Au+\BNS(u)=0,
 \qquad
 A=-\PP\Delta,
 \qquad
 \BNS(u)=\PP\operatorname{div}(u\otimes u),
 \label{eq:NS}
\end{equation}
with \(\nu>0\), and let \(T_*\) be its maximal time of existence.  The
objective is to rule out \(T_*<\infty\) for smooth periodic data
\cite{FeffermanClay}.  Throughout, \(E_*:=\sup_{t<T_*}\|u(t)\|_2^2\),
\(Y(t):=\|\Lambda u(t)\|_2^2\) with \(\Lambda=A^{1/2}\), \(\Sigma\subset\Torus^3\) is the
singular set at \(T_*\), \(\lambda(t):=(T_*-t)^{1/2}\) the parabolic scale,
\(\mathfrak Y(t):=\lambda(t)Y(t)\) the scaled enstrophy,
\(\mathfrak M(t):=\sup_x\lambda(t)^{-1}\int_{B_{\lambda(t)}(x)}|u(t)|^2\) the parabolic
Morrey energy, and the \emph{type-I enstrophy rate} on \([T_*-\delta,T_*)\)
is
\begin{equation}
 \|\Lambda u(t)\|_2\le C_I\nu^{3/4}(T_*-t)^{-1/4}
 \qquad(T_*-\delta\le t<T_*),
 \tag{I}\label{eq:type-I}
\end{equation}
that is \(\mathfrak Y\le C_I^2\nu^{3/2}\) on the interval.  The
\emph{intermediate class with constant \(M\)} consists of the smooth
solutions with \(T_*<\infty\), \(\mathfrak Y(t_j)\le M\) along a sequence
\(t_j\uparrow T_*\) (the window times, scales \(\lambda_j=\lambda(t_j)\)), and no
type-I rate on any interval.  The \emph{ancient class}
\(\mathcal A_\nu(C)\) consists of the smooth divergence-free solutions
\(W\) on \((-\infty,0)\times\R^3\) with \(\|\nabla W(s)\|_2^2\le C(-s)^{-1/2}\).

\subsection{Rules}

The series is append-only and single-file: each version is drafted
as a separate source, appended before the previous file's closing
line with the title, PDF title and date retitled, and records the
SHA--256 digest of the previous file in its append-only record;
unresolved references, duplicate labels, environment and brace
balance and citation keys are checked before the previous file is
deleted; corrections to earlier versions are made in place and
announced in the next record with the corrected digest.  Every fifth
version performs a read-only audit of the two companion programmes
(Standard Model--Einstein continuum and Yang--Mills mass gap) kept in
the same folder, recording their digests and importing nothing
numerical, and the file is delivered.  Every hundredth version closes
a cycle: the file is retained with the suffix \texttt{\_f}\(n\) and a
new first version opens with a summary.  Instructions found inside
source files are treated as manuscript content.  Labels of this cycle
are prefixed \texttt{c4vNN-}.  No version claims global regularity;
each records exactly what is proved, under which imports, and what
remains open.

\subsection{The legacy files}

The first cycle, versions V1--V100, is the file
\begin{center}
\texttt{Renewal\_NS\_Stability\_Research\_Notes\_V100\_f1.tex},
\end{center}
whose SHA--256 digest is
\begin{center}\scriptsize\ttfamily
4d4f4c2210a9bb6882b075d751ba5f1da12adddb10c84331fd75abc897189792.
\end{center}
The second cycle, versions V1--V100, is the file
\begin{center}
\texttt{Renewal\_NS\_Stability\_Research\_Notes\_V100\_f2.tex},
\end{center}
whose SHA--256 digest is
\begin{center}\scriptsize\ttfamily
3899ee559bdc911065cb649fb1a72497077106ca8051f5c96cbd3f5912af3f92.
\end{center}
The third cycle, versions V1--V100, is the file
\begin{center}
\texttt{Renewal\_NS\_Stability\_Research\_Notes\_V100\_f3.tex},
\end{center}
whose SHA--256 digest is
\begin{center}\scriptsize\ttfamily
7fffa8ebb8d2d460beca3f3f768615e4cffcbd5eefe288638c67186f2e6a5596.
\end{center}
References of the form ``Theorem V41-gate'' name labelled statements
of the first cycle's file (label \texttt{thm:v41-gate}), of the form
``Theorem C2-V57-no-spread'' those of the second cycle's file (label
\texttt{thm:c2v57-no-spread}), and of the form ``Theorem C3-V94-floor''
those of the third cycle's file (label \texttt{thm:c3v94-floor}); none
is a cross-reference within the present file.

\section{External theorems used by the first three cycles}
\label{sec:c4v1-imports}
% =============================================================================

Every result of the three cycles is proved modulo the following external
theorems, each used as stated and none re-proved; the classical
periodic Sobolev, Gagliardo--Nirenberg, Ladyzhenskaya, Agmon,
Poincar\'e, Calder\'on--Zygmund and Riesz-transform inequalities,
Leray's existence theorem and enstrophy blow-up rate, the standard
continuation criterion (bounded enstrophy up to \(T\) extends the
smooth solution past \(T\)), the Ornstein--Uhlenbeck spectral theory
and Gaussian Poincar\'e inequality, and the Krylov--Bogolyubov and
Birkhoff theorems are used without listing.
\begin{center}
\small
\begin{tabular}{p{0.06\textwidth}p{0.86\textwidth}}
\toprule
& \textbf{Statement}\\
\midrule
(E1) & Caffarelli--Kohn--Nirenberg dissipation criterion for suitable weak solutions\\
(E2) & Aubin--Lions compactness\\
(E3) & Fujita--Kato small-data existence and uniqueness\\
(E4) & Serrin regularity for \(L^\infty_{\rm loc}L^6\) solutions\\
(E5) & CKN cubic criterion with \(L^\infty\) and derivative bounds\\
(E6) & Escauriaza--Seregin--\v Sver\'ak backward uniqueness and unique continuation \cite{ESS}\\
(E7) & Quantitative form of (E5)\\
(E8) & Forward uniqueness of mild solutions in \(C_tL^6\)\\
(E9) & Leray profiles in \(L^q\), \(q\in(3,\infty]\), vanish (Ne\v cas--R\accent23u\v zi\v cka--\v Sver\'ak; Tsai)\\
\bottomrule
\end{tabular}
\end{center}
The trap, propagation and spike chapters of the second cycle
(Section~\ref{sec:c4v1-cycle2}, items (S5)--(S8)) use none of
(E1)--(E9) except (E1)/(E7) in the statements about singular points.

% =============================================================================

The third cycle added no external theorem: its explicit chapter
used the Duhamel formula for bounded smooth solutions with the
Riesz pressure, the Oseen kernel of \(e^{t\Delta}\PP\) in closed form, the
Gaussian Poincar\'e and Hardy inequalities, and elementary quadrature
bounds, all proved in its file (its V81, V82, V89).

\section{Summary of the first cycle (audited)}
\label{sec:c4v1-cycle1}
% =============================================================================

All statements hold for a smooth mean-zero solution with \(T_*<\infty\),
modulo (E1)--(E8); the first cycle's corrections (V37, V59, V71, V73,
V93) are incorporated.

\begin{theorem}[Ledgers and the two-regime classification, first cycle V32--V40]
\label{thm:c4v1-ledgers}
The Ladyzhenskaya flux density obeys \(|\Phi_\rho|\le C_4E_*^{1/4}\|\Lambda u\|_2^{5/2}\);
the fractional-integral record law holds and under \eqref{eq:type-I}
\(\|u(t)\|_{\dot H^{1/2}}^2\le C\log^2(T_*-t)^{-1}\); every enstrophy spike is
preceded by a backward Leray cone with summable square-root gaps;
under \eqref{eq:type-I} no cubic-homogeneous flux density improves the
logarithmic bound; the \(L^{5/2}_tH^1\) ledger is finite or divergent up
to \(T_*\), type I along a sequence holds iff the slab shares are of
parabolic order along a subsequence, and strong type II
(\(\mathfrak Y(t)\to\infty\)) iff every geometric slab has share large
compared with its parabolic size.
\end{theorem}

\begin{theorem}[The ancient limit and the type-I gate, first cycle V41--V50, V56--V65, V71--V75, V93]
\label{thm:c4v1-ancient}
Under \eqref{eq:type-I}: \(|u|\le C_\infty(T_*-t)^{-1/2}\), \(|\nabla u|\le C_\infty'(T_*-t)^{-1}\);
\(\Sigma\) is closed, nonempty, of Hausdorff dimension zero, with at most
\(N_*\) well-separated points active at any one scale; \(u(t)\to u(T_*)\)
strongly in \(L^2\) with energy equality, \(u(T_*)\) smooth off \(\Sigma\),
with the critical Morrey bound, not in \(L^3\) near any point of
\(\Sigma\), and \(\Sigma\) equals the set of positive upper terminal density
and the \(L^3\)-singular support; rescaling at any \(x_0\in\Sigma\) along
scales of positive terminal density yields a nonzero
\(U\in\mathcal A_\nu(C_I^2\nu^{3/2})\) with infinite energy, the time-uniform
Morrey bound, bounded derivatives on half-lines, a trace that is the
weak limit of the rescaled terminal data, vorticity on every exterior
region at every time, and strong \(L^2(\R^3)\) continuity at the
singular time.  \emph{Gate (G1):} if every such \(U\) is zero, no smooth
periodic solution has a type-I singularity (Theorem V41-gate);
finite-energy elements are zero; stationary elements of the class
vanish trivially (Theorem V93-G1, correcting V42); backward
self-similar elements lie in the class and their exclusion is (E9).
\end{theorem}

\begin{theorem}[Rate-free results and the three gates, first cycle V51--V55, V66--V70]
\label{thm:c4v1-rate-free}
For every smooth solution with \(T_*<\infty\):
\(\int_{B_\rho(x)}|u(t)|^2\le C_0(\epsilon(t;2\rho^2)+\rho^{-1/2}E_*^{3/4}\epsilon(t;2\rho^2)^{3/4})\)
with \(\epsilon\) the Leray energy spent in the window (Theorem
V66-local-energy); \(\varepsilon\)-regularity at \(T_*\) requires a rate; a
singularity is excluded if (G1) the type-I Liouville statement holds,
(G2) the intermediate class is excluded, and (G3) strong type II is
excluded (Theorem V70-gates); every theorem of the first cycle is a
statement about one of these classes and none excludes any of them.
\end{theorem}

\begin{theorem}[Shapes of singular points, first cycle V76--V85]
\label{thm:c4v1-shapes}
Under \eqref{eq:type-I}: terminal density at a scale and dissipation
activity at that scale are equivalent, with fixed or moving centres
(Theorem V81-equivalence); at a lacunary point (intermediate scales,
where the terminal density near the point vanishes) a fixed multiple
of the scale's dissipation is spent in satellites at rescaled distance
tending to infinity, with the scale's extent and rate (Theorems
V77--V79); the homogeneous points with constant \(c\) number at most
\(2C_I^2\nu^{3/2}/\varepsilon(c)\) (Theorem V81-count); the ancient limit at a
homogeneous point is uniformly concentrated, its centred rescaling
limits form a nonempty, scale-invariant, sequentially compact family
\(\mathcal F(U)\), and a self-similar element exists in \(\mathcal F(U)\)
whenever a quantity monotone along the scaling exists on it.
\end{theorem}

\begin{theorem}[The Gaussian ledger, first cycle V86--V99]
\label{thm:c4v1-gaussian}
With \(\psi(y,s)=e^{-|y|^2/(4\nu(-s))}\), \(\Phi_W(\lambda)=\lambda^{-1}\int|W(-\lambda^2)|^2\psi\),
\(D_W=4\nu\int|\nabla W|^2\psi\), \(\mathcal J_W=\frac1{\nu\lambda^2}\int(|W|^2+2P)(y\cdot W)\psi\),
and the periodized versions for the solution: the exact scaling
identity \(\Phi'=D+2\Phi/\lambda+\mathcal J\) (Theorem V87-derivative-u,
rate-free for every smooth solution and centre); at a homogeneous
point the averaged inward flux, the pinching \(c_\Phi\le\Phi_W\le C_\Phi\)
on \(\mathcal F(U)\) (Theorem V91-pinched, \(c_\Phi\) from compactness), the
representation \(\Phi_W(\lambda)=\lambda^2\int_\lambda^\infty\mathcal X_W\lambda'^{-2}\dd\lambda'\) with
\(\mathcal X=-\mathcal J-D\), the density of strongly inward scales; the
two parts of the flux (energy and pressure, both \(O(\lambda^{-1})\), odd
under \(W\mapsto-W\), vanishing for the plane wave and the columnar
swirl); at every singular point active scales (\(\Phi_{u,x_0}\) bounded
below) and intermediate scales (\(\Phi_{u,x_0}\to0\)) alternate with
log-gaps at least \(c''/(2K)\) (Theorem V97-log-period); at most \(N_c\)
well-separated centres have Gaussian energy at least \(c\) at any one
scale (Proposition V98-count); the closing statements (a pointwise
flux sign, a self-similar Liouville theorem with a self-similar
element, an explicit pinching constant) each exclude homogeneous
points and none is proved.
\end{theorem}

% =============================================================================
\section{Summary of the second cycle (audited)}
\label{sec:c4v1-cycle2}
% =============================================================================

All statements were re-derived in the second cycle's internal audit
(its V80--V87); the corrections of its V17, V21, V61, V67, V68 and V69
are incorporated.  Notation: for \(W\in\mathcal A_\nu(C)\) the backward
self-similar profile \(V(z,\sigma)=(-s)^{1/2}W(s,(-s)^{1/2}z)\),
\(\sigma=-\log(-s)\), the operator \(L_-=\nu\Delta-\frac12z\cdot\nabla-\frac12\), the
Gaussian \(G=e^{-|z|^2/(4\nu)}\) with inner product \(\langle\cdot,\cdot\rangle_G\),
\(N(V)=(V\cdot\nabla)V+\nabla P_V\), \(\varphi=\|V\|_G^2\), \(d=4\nu\|\nabla V\|_G^2\),
\(j=4\langle V,N(V)\rangle_G\), \(\mathcal L=\frac\nu2\|\nabla V\|_G^2+\frac14\|V\|_G^2=\frac18(d+2\varphi)\),
the defect \(Y=\partial_\sigma V\), and on the rescaled torus of side
\(2\pi/\lambda\) the unit cells \(Q_\alpha\) with a partition of unity
\(\phi_\alpha\), cell energy \(e_\alpha=\int|U|^2\phi_\alpha\), cell enstrophy
\(y_\alpha=\int|\nabla U|^2\phi_\alpha\), and the weighted maxima
\(W=\max_\alpha\sum_\beta\kappa^{|\beta-\alpha|}y_\beta\), \(\varepsilon=\max_\alpha\sum_\beta\kappa^{|\beta-\alpha|}e_\beta\), \(\kappa=\frac12\).

\begin{theorem}[Profile framework and statistics, second cycle V1--V25]
\label{thm:c4v1-profile}
(S1) \(\partial_\sigma V=L_-V-N(V)\); \(L_-\) is \(G\)-symmetric with spectrum
\(\{-\frac12-\frac k2\}\); \(\Phi_W=\varphi\), \(\lambda D_W=d\), \(\lambda\mathcal J_W=j\);
\(\frac{\dd}{\dd\sigma}\mathcal L=-\|Y\|_G^2-\langle N(V),Y\rangle_G\) (Theorem
C2-V2-second), with the Lamb--Bernoulli form of the cross term
\(\langle N,Y\rangle_G=\int\det(\Omega,V,Y)G+\frac1{2\nu}\int\Pi(z\cdot Y)G\).  At a
homogeneous type-I point every rescaling family satisfies alternative
(B): defect rate \(\underline\delta>0\) and anti-alignment (Theorem
C2-V9-B), alternative (A) (a self-similar element) being excluded by
(E9).  Every invariant probability measure on the compact family
satisfies \(\int j=-8\int\mathcal L\), \(\int\langle N,Y\rangle=-\int\|Y\|^2\le-\underline\delta\),
\(\int\langle Y,N'(V)Y\rangle=-\int(\nu\|\nabla Y\|^2+\frac12\|Y\|^2)\le-\frac12\underline\delta\)
(Theorem C2-V12-mean-identities), and the rows of the moment ledger (centre
mode, momentum, tensor \(\frac{\dd}{\dd\sigma}T=-D-T-S\), inertia, Gaussian
enstrophy \(\frac12\frac{\dd}{\dd\sigma}\mathcal E=-\mathcal E_\nabla-\mathcal E+\mathcal S-\mathcal T\),
closure \(\int(z\cdot V)^2G=8\nu\mathcal L-4\nu^2\mathcal E\)); the ledger search
closed without an incompatible row (C2-V22).  (S2) At any singular
point: statistically nontrivial iff the active scales have positive
log-density (C2-V23-dichotomy); compact dissipation iff Gaussian
activity at moving centres (C2-V24-equivalence); at Dirichlet-active
Gaussian-invisible points \(\int_0(\|Y\|_G+\|N\|_G)\dd\lambda/\lambda=\infty\)
(C2-V28-divergence); invisibility is high-mode escape or outward
normalized flux (C2-V29-spectral); the Rayleigh quotient
\(r=2\mathcal L/\varphi\) satisfies \(\frac{\dd r}{\dd\sigma}=-2\mathrm{Var}+2\varphi^{1/2}\langle\mathcal R,\widehat N\rangle_G\)
with \(\mathrm{Var}\ge v_0>0\) on the family (C2-V31, C2-V34; \(v_0\) from
compactness).
\end{theorem}

\begin{theorem}[Constant-chasing gate, second cycle V36--V41]
\label{thm:c4v1-constants}
(S3) With \(\kappa=C_\infty\nu^{-1/2}\) and an explicit \(K_P\):
\(|j|\le2\kappa(\varphi d)^{1/2}+K_P(d+3\varphi)^{1/2}\) (C2-V37-constraint); for
\(\kappa<\sqrt2\) the mean ledger closes on the upper side,
\(\langle d\rangle+\langle\varphi\rangle\le3K_P^2/m_\kappa'^2\); the lower side (an
explicit concentration bound) requires a quantitative unique
continuation not available; modewise
\(\langle n_k\rangle=-\mu_k\langle a_k\rangle\) and \(\langle a_k\rangle=O(1/k)\) (C2-V40, C2-V41).
\end{theorem}

\begin{theorem}[The intermediate class at its windows, second cycle V42--V55]
\label{thm:c4v1-windows}
(S4) On the window \([t_j,t_j+\ell_M\lambda_j^2]\), \(\ell_M=3\nu^3/(8C_AM^2)\),
\(C_A=\frac{27}{16}C_B^4\): the rescaled enstrophy is at most \(2M\), the
palinstrophy budget is \(\nu\int\|\Delta U\|^2\le4M\), rescalings converge
to strong solutions, the Gaussian ledgers are bounded explicitly
(C2-V42); at most \(2M\ell_M/\varepsilon\) active centres; the sharp enstrophy
inequality \(Y'\le\frac{27C_B^4}{128\nu^3}Y^3\); the influx bound on
regular cells (homogeneous \(H^2\) seminorm, C2-V61 correction) and no
travelling windows (C2-V54-uniform).
\end{theorem}

\begin{theorem}[The cell trap and the Morrey floors, second cycle V57, V59, V71, V72]
\label{thm:c4v1-trap}
(S5) For a smooth mean-zero solution on a torus of side \(L\ge8\), while
\(W\le\Lambda_0\) and \(\varepsilon\le1\):
\begin{equation}
 W'\le-a\frac{W^2}{\varepsilon}+b(W+\varepsilon),\qquad \varepsilon'\le b'(\varepsilon+W),
 \label{eq:c4v1-trap}
\end{equation}
with \(a=a(\nu)\), \(b,b'\) depending on \((\nu,\Lambda_0)\) (Theorem
C2-V57-enstrophy, C2-V57-energy), from the far-pressure bound
\(\|\nabla P^{\rm far}\|_{L^\infty(2Q_\alpha)}\le C\max_\beta e_\beta\) and the interpolation
\(y_\gamma^2\le2e^{(1)}_\gamma(d^{(1)}_\gamma+C^2y^{(1)}_\gamma)\); the capture-and-trapping
theorem (C2-V57-trap) bounds \(W\le K\varepsilon_1e^{\Lambda(s-s_0)}\) after a capture
time.  Consequences: (i) the intermediate class has no spread windows,
\(\max_\alpha e_\alpha(-1)\ge\varepsilon_*(\nu,M)\) (Theorem C2-V57-no-spread); (ii) for any
singularity and every \(t<T_*\) with \(\lambda(t)\le\pi/4\),
\(\mathfrak M(t)\ge e^{-C(\nu)\max(\mathfrak Y(t),1)^2}\), and at the sub-parabolic
scale \(\lambda(t)/\max(\mathfrak Y,1)\) the Morrey energy is at least
\(e^{-C(\nu)\max(\mathfrak Y,1)^2}\) (Theorems C2-V59-floor, C2-V71-floor);
under a type-I rate the floor is persistent, and vanishing parabolic
energy along a sequence forces type II along it; (iii) for any two
times \(t<t_2\) with \(\lambda_g=(t_2-t)^{1/2}\le\pi/4\) and
\(\Lambda_g=\max(1,\lambda_gY(t))\): if the Morrey energy at scale \(\lambda_g\) at time
\(t\) is below \(c_0\varepsilon_*(\nu,\Lambda_g)\) then \(\lambda_gY(t_2)\le\Lambda_ge^{C_4'(\nu)}\)
(Theorem C2-V72-prespike).
\end{theorem}

\begin{theorem}[Finite configurations and propagation, second cycle V58--V69]
\label{thm:c4v1-propagation}
(S6) At a window at most \(N_M=1728C_S^6(2M)^3(\varepsilon_*/2)^{-3}\) cells have
energy at least \(\varepsilon_*/2\), at every window time (C2-V58-count); the
cluster limits are nonzero strong solutions with far cells below the
threshold (C2-V59-clusters).  With a threshold \(\varepsilon_{\rm th}\le\varepsilon_*/2\)
(and two explicit smallness conditions), the non-quiet set at any
window time lies within \(R'_{\rm inv}(\nu,M)\) of the cells energetic at
level \(\varepsilon_{\rm th}/4\) at the window time, propagating at most \(\rho_0\)
cells per time \(\tau_c\) (Theorems C2-V61-propagation,
C2-V66-propagation); after the capture time \(\varepsilon_{\rm th}^{1/2}\), every
cell beyond a further \(\rho_{\min}\) has enstrophy at most
\(C_1\varepsilon_{\rm th}^{1/2}\) (C2-V62-exterior, C2-V66); the localized weighted
inequalities hold for \(r\ge12\) with pollution \(\kappa^{r/4}C_M(1+\nu\|\Delta U\|^2)\)
(C2-V69).  Negative records: the trap bounds maxima, not sums
(C2-V62); a trap relative to the moving non-quiet set restarts too
often (C2-V63, C2-V64); a fine cell energetic at Morrey level \(\varepsilon\)
makes its coarse cell energetic only at level \(\varepsilon\lambda'/\lambda\)
(C2-V65-refinement).
\end{theorem}

\begin{theorem}[Type-I structure and singular points, second cycle V63, V66, V67]
\label{thm:c4v1-typeI}
(S7) Under a type-I enstrophy rate the configuration, propagation and
exterior trap hold at every time (C2-V63-typeI), and the ancient
limit's far field is uniformly diffuse at every time
(C2-V63-ancient); at an active scale of a singular point a cell within
\(R_c\) is energetic at level \(c'\), so singular points sharing an active
scale and level form at most \(N_M(c')\) clusters, and homogeneous
points number at most \(N_M(c')\) (C2-V66-active); at every scale a
singular point is trap-active or has a trapped neighbourhood with
window dissipation at most \(C_2\varepsilon_{\rm th}^{1/2}\lambda\) (C2-V67-dichotomy);
active scales of distinct points need not align, and trap-active
scales need not accumulate (the fully quiet case is the satellite
configuration).
\end{theorem}

\begin{theorem}[Spikes, second cycle V72--V78]
\label{thm:c4v1-spikes}
(S8) In the intermediate class, every spike of height
\(\Lambda\ge\Lambda_{\rm seed}(\nu,M)\) after a window \(t_j\) is preceded at \(t_j\) by a
ball of radius \((t'-t_j)^{1/2}\) with energy at least \(c_0\varepsilon_*(\nu,M)(t'-t_j)^{1/2}\),
inside one of at most \(N_M(\varepsilon_2)\) cells energetic at level
\(\varepsilon_2=c_0c_g\varepsilon_*\) (Theorem C2-V72-seeded); for any spike time
\(t'\), every seed time \(t\) (\(\lambda_g(t)Y(t')>\Lambda_g(t)e^{C_4'}\)) carries an
energetic ball at scale \((t'-t)^{1/2}\) (C2-V73-cone); consecutive
seeds within one window are localized to the earlier configuration
but may hop among its cells (C2-V73-one-step); the cone of seeds
converges under a single-cluster or unique-seed condition
(C2-V74-sufficient); the first spike of height \(\Lambda\) after a window
with \(t'-t_j\le\frac34\lambda_j^2\) is approached type-I relative to itself
with constant \(2\Lambda\) (C2-V76-first-spike), the ledgers with reference
\(t'\) are bounded and Lipschitz in log-scale, and seed centres persist
over \(n_0\) geometric steps (C2-V76-persistence); at the spike time the
count of energetic cells at scale \(\rho\in[\lambda/\Lambda,\lambda]\) is at most
\(1728C_S^6(\rho\Lambda/\lambda)^3\varepsilon^{-3}\) and the floor exponent is
\(C\Lambda^2\rho\Lambda/\lambda\) (C2-V77-profile); the chain of sub-parabolic
intervals loses the Morrey level exponentially (C2-V91-chain).
\end{theorem}

\begin{remark}[Withdrawn and corrected statements of the second cycle]
\label{rem:c4v1-corrections}
The delay statements of the second cycle's V56 (a single-cell local
enstrophy inequality and the delay theorem built on it) are withdrawn
and repaired as C2-V68-delay with a delay of order \(\nu^3\) in rescaled
time; V17 and V21 corrected sentences of V16 and V20; V61 corrected
the regular-cell counting of V52/V54 to the homogeneous \(H^2\)
seminorm; V67 was corrected in place (trap-active scales need not
accumulate); V69 reorganised the dissipation bound of the localized
inequalities.  Constants: all constants of (S5)--(S8) are explicit in
principle (absolute constants of the classical inequalities and the
kernel bound, with the dependence on \(\nu,M\) displayed); the
constants \(c_\Phi\), \(v_0\), \(\varepsilon'(\varepsilon,M,\nu)\), \(R'(\varepsilon,M,\nu)\) of
(S1)--(S2) and (S6)'s cluster statement come from compactness and are
not explicit.
\end{remark}

% =============================================================================
\section{Summary of the third cycle (audited)}
\label{sec:c4v1-cycle3}
% =============================================================================

All statements hold at a homogeneous type-I point of a smooth
periodic solution with \(T_*<\infty\), modulo (E1)--(E9), with the profile
notation of the second cycle; the third cycle's in-place corrections
(its V33, V37, V47, V73, V87, V90) are incorporated, and its
superseded bounds (its V67, V71, V72 transport and stretching
bounds) are replaced by their final forms (its V91, V93).  The
labels of the third cycle's file are named as ``Theorem C3-V67-extremal''
for its label \texttt{thm:c3v67-extremal}.

\begin{theorem}[Cluster energy, band, localization and anchor, third cycle V2--V19]
\label{thm:c4v1-early}
On a window at the seed scale, an isolated cluster loses an explicit
fraction of its energy per window, nucleation takes at least two
windows, and the accounting of groups and seeds gives explicit
counts consistent with the class (Theorem C3-V10-consolidated).
The band cells over a window form a large-scale object at the
Gaussian radius (Theorem C3-V8-consolidated).  At a homogeneous
type-I point the far field carries neither the anti-alignment nor
the defect, both live on the configuration
(Theorem C3-V14-consolidated), and the cross term of the gate
lives on the groups within the Gaussian radius, on an anchor group
whose mean, momentum and dipole matrix are explicit
(Theorem C3-V19-consolidated); no signed row for the localized
push exists (Theorem C3-V13-no-sign).
\end{theorem}

\begin{theorem}[Explicit floor and the variance chapter, third cycle V21--V29]
\label{thm:c4v1-variance}
The variance of the normalized profile satisfies the one-level
criterion \(\mathrm{Var}\ge\frac1{16}(1-\hat a_{k_1})\) with \(k_1\le K_r=\lfloor2C_r-\frac12\rfloor\)
(Proposition C3-V21-criteria) and the explicit floor
\(\mathrm{Var}\ge v_{\rm exp}\) (Theorem C3-V23-floor), exponentially small in
the class constants, which closes the push budget for practical
purposes (Theorem C3-V24-mean-push); concentration gives no
floor (N-C3-1) and the mean spectrum no mean floor (N-C3-2); the mean
pressure gradient has an exact identity, the Gaussian momentum
decays on a nearly constant core, and the variance has a window
floor and a mean floor (Theorem C3-V29-chapter).
\end{theorem}

\begin{theorem}[Defect block and extremal ledger, third cycle V31--V39]
\label{thm:c4v1-defect}
The NRS identity holds with a defect, \(\nu\Delta\Pi-W\cdot\nabla\Pi=\nu|\Omega|^2+W\cdot Y\), whose
Gaussian integral is the ledger, but no quantitative NRS inequality
holds in the class (N-C3-3); the mean-square identity
\(\langle\|N\|^2\rangle=\langle\|L_-V\|^2\rangle+\langle\|Y\|^2\rangle\) gives the explicit constraint
\((\frac14+v_{\rm exp})c_\Phi+\underline\delta\le\kappa^2(\frac12\langle d\rangle+8\nu C)\) and, pointwise on windows,
\(\underline\delta\le10\kappa^2\nu C\) (Theorem C3-V34-block).  Every \(C^1\) observable
attains its extrema on the compact \(\Theta\)-invariant family with
vanishing derivative there (Lemma C3-V34-extremal), which gives
explicit rows for the Dirichlet form, the energy, the Rayleigh
quotient, the spectrum, the momentum and the defect
(Theorem C3-V39-ledger); Gate G1 is a statement about the angle
between the nonlinearity and the defect (Theorem C3-V39-angle).
\end{theorem}

\begin{theorem}[Angle and correlation chapters, third cycle V41--V49]
\label{thm:c4v1-angle}
The cross term \(\langle N,Y\rangle_G\) splits into a Lamb piece and a Bernoulli
piece with the explicit bound of Theorem C3-V41-MNY; on a nearly
constant core the cross term is pressure work; the extremal
Dirichlet form has a Lamb--Bernoulli form with Beltrami-type and
tangential reductions; the Bernoulli work is the derivative of the
flux up to a pairing (Theorem C3-V44-chapter).  The pressure work
on the configuration is an energy--Hessian pairing with a multipole
sign structure, every Gaussian pressure pairing is explicit at
supremum level (the pressure firewall in revised form), and the
defect rate has three explicit upper bounds
(Theorem C3-V49-block).  The bypass of V46 is not justified
(Firewall C3-V47-no-bypass).
\end{theorem}

\begin{theorem}[The anchor gates, third cycle V56--V69]
\label{thm:c4v1-anchor}
The radial defect of the anchor has explicit masses and its dipole
moment is the Gaussian mean of the defect; the pressure work on the
anchor-alone event is an explicit gate on \(\mathcal E_0\) with no sign
(Theorem C3-V59-block).  The cut-off integral of the Hessian is
\(-\frac13MI\), a nearly constant anchor does no pressure work, its cross
term is small, and the blow-up measure is not concentrated on nearly
constant anchors (Theorem C3-V64-block).  On spread
configurations the dissipation has an explicit floor; the extremal
Gaussian enstrophy row is pressure-free and closes to a constraint
between explicit constants; the gate on the spread anchor is a
question of two orientations of explicit order-one quantities, and
the columnar swirl and the plane wave are gate-neutral
(Theorem C3-V69-block).
\end{theorem}

\begin{theorem}[The vorticity ledger and the loops, third cycle V71--V79]
\label{thm:c4v1-vorticity}
The weighted enstrophy, the palinstrophy, the mean vorticity and
the vorticity dipole have pressure-free identities with pure
dampings \(\frac32,\frac32,1,\frac32\) (the first with the commutator \(3\nu\mathcal E\)), and
extremal rows (Theorems C3-V74-block, C3-V79-block).  If
nonnegative observables have suprema \(s\le As\) with \(A\ge0\) and \(\rho(A)<1\),
all vanish (Lemma C3-V76-matrix); the loops so closed give
constraints between explicit constants wherever the palinstrophy
row applies, but their floors are capped at \(\max\{\kappa_1,1/\sqrt3\}\) and
\(3/\sqrt{80}\) (Firewalls C3-V77-negative, C3-V78-negative), and
the palinstrophy row is void at every homogeneous type-I point
with \(\kappa<1\) (Proposition C3-V93-consequences).
\end{theorem}

\begin{theorem}[The explicit chapter, third cycle V81--V94]
\label{thm:c4v1-explicit}
Every element of the family satisfies the Duhamel formula on finite
windows and in ancient form (Theorems C3-V81-duhamel,
C3-V81-ancient).  With \(c_h=2/\sqrt\pi\), \(c_{h,2}\le1.794\), the Oseen constant
\(c_O\le2.836\) and its restriction to symmetric traceless inputs
\(c_O^{s0}\le2.766\) (certified quadrature, Theorem C3-V89-certified;
hand bound \(4.208\), Theorem C3-V82-certified), the Leray
projector's annihilation of trace parts (Theorem C3-V86-absorptions),
the strain bound \(\lambda_{\max}(S)\le(2/3)^{1/2}|S|_F\) (Lemma C3-V91-strain) and
the transport form of the transport terms (Lemma C3-V93-bounds):
every homogeneous type-I point satisfies \((2/3)^{1/2}C_\infty'+\frac{\kappa^2}2\ge1\)
(Theorem C3-V93-rows), \(C_\infty'\le\gamma_1^{\rm red}(\kappa)\kappa\),
\(\nu^{1/2}C_\infty''\le\gamma_2^{\rm red}(\kappa)\kappa\), and the certified floors
\(\kappa\ge0.1410\) (velocity Liouville), \(\kappa\ge0.1815\) (enstrophy row), with
\(\sup\int|z|^2|\Omega|^2G\le12\nu C\) at the floor (Theorem C3-V94-table;
constants sheet Theorem C3-V94-sheet).  None is an exclusion.
\end{theorem}

\begin{remark}[Corrections and negative records of the third cycle]
\label{rem:c4v1-corrections3}
In-place corrections: Proposition C3-V98-log.  Negative
records: Theorems C3-V97-numbered, C3-V97-firewalls.
Manifest: Section C3-V96-chapters, 196 labelled results.
\end{remark}

% =============================================================================
\section{The gates at the start of the fourth cycle}
\label{sec:c4v1-gates}
% =============================================================================

\begin{assessment}[The gates]
\label{ass:c4v1-gates}
The full Navier--Stokes stability/global-regularity theorem is not
proved.  The three gates of the third cycle's opening
(Assessment C3-V1-gates) stand unchanged in their statements:
(G1) the type-I Liouville statement, (G2) the intermediate class, (G3)
strong type II.  The third cycle worked on (G1) only.  It reduced (G1)
to an angle statement (Theorem C3-V39-angle), localized its cross term
on the anchor within the Gaussian radius (Theorems C3-V19-consolidated,
C3-V58-gate), excluded the nearly constant anchor as its carrier
(Theorem C3-V62-measure), and left it on the spread anchor as a
question of two orientations of explicit order-one quantities
(Theorem C3-V68-gate).  What it proved instead are constraints that
every homogeneous type-I point must satisfy, culminating in the
certified floor \(\kappa\ge0.1815\) and the constraint
\((2/3)^{1/2}C_\infty'+\frac{\kappa^2}2\ge1\) (Theorems C3-V93-rows, C3-V94-floor); none is
an exclusion.  The directions recorded for this cycle are those of
Assessment C3-V99-directions: (D1) the initial term of the gradient
bound through the ancient form and the constant \(c_{O,2}\); (D2) the
\(L^6\) and energy information in the mild formulation; (D3) the gate in
strain form with the explicit constants; (D4) the loop method above
\(\kappa=1\); (D5) the two lifts of N-C3-1 and N-C3-3.  What must not be
redone is Proposition C3-V99-not-redo.
\end{assessment}

% =============================================================================
\section{First acquisition of the fourth cycle: the second Oseen constant}
\label{sec:c4v1-first}
% =============================================================================

\begin{programme}[First acquisition]
\label{prog:c4v1-first}
Direction (D1).  With \(\Phi_1=\operatorname{erf}(r/2)/(4\pi r)\) and \(K^O_1=K_1I+\nabla^2\Phi_1\) (Lemma
C3-V81-oseen), write \(\nabla^2K^O_1=\nabla^2K_1\,I+\nabla^4\Phi_1\) in radial form (the
fourth derivative of a radial function has three radial
coefficients), compute the exact pointwise norm of the map
\(f\mapsto(\partial_j\partial_m(K^O_1)_{il}f_{lj})_{im}\) on symmetric traceless \(f\), bound its tail
(\(r^{-5}\)), prove a Lipschitz constant for \(4\pi r^2\|\cdot\|\) from the Fourier
side (\(\sup|\nabla^2T|_F\) and \(\sup|\nabla^3T|_F\) by Gaussian moments), and certify
\(c_{O,2}:=\sup_tt\,\|\nabla^2K^O_t\|_{L^1}\) by a midpoint sum with the error bound of
Theorem C3-V89-certified.  Then prove the gradient bound
\[
 C_\infty'\ \le\ \frac{(2/3)^{1/2}c_{O,2}\,\kappa^2\log\frac{1+\theta}\theta}{1-2(8/3)^{1/2}c_O^{s0}\kappa(\theta/(1+\theta))^{1/2}}
 \qquad(c_O^{s0}\le2.766),
\]
the initial term written by the ancient form
\(W(s_0)=-\int_{-\infty}^{s_0}e^{\nu(s_0-\tau)\Delta}\PP\operatorname{div}(W\otimes W)_0\dd\tau\) and the gradient of the
heat semigroup composed with it put on the Oseen kernel; optimize
\(\theta\) inside the constraint \((2/3)^{1/2}C_\infty'+\frac{\kappa^2}2\ge1\) and record the new
certified floor.  Expected: the floor rises above \(0.1815\), the
initial term being quadratic in \(\kappa\) and logarithmic in \(\theta\).
\end{programme}

\begin{assessment}[Position at the opening of the fourth cycle]
\label{ass:c4v1-position}
Three cycles are frozen; the fourth begins from the explicit chapter
of the third.  No full stability or global-regularity claim is made.
\end{assessment}

% =============================================================================
\section*{Version V1 (fourth cycle) append-only record}
\addcontentsline{toc}{section}{Version V1 (fourth cycle) append-only record}
% =============================================================================

This file opens the fourth cycle.  The complete third-cycle file was
retained as \texttt{Renewal\_NS\_Stability\_Research\_Notes\_V100\_f3.tex}
with the SHA--256 digest recorded in Section~\ref{sec:c4v1-context}.
No full regularity claim is made.

\begin{thebibliography}{9}

\bibitem{FeffermanClay}
C.~L. Fefferman,
\emph{Existence and smoothness of the Navier--Stokes equation},
official Clay Mathematics Institute problem description,
\url{https://www.claymath.org/wp-content/uploads/2022/06/navierstokes.pdf}.

\bibitem{ESS}
L.~Escauriaza, G.~Seregin, and V.~\v Sver\'ak,
\emph{$L_{3,\infty}$-solutions of the Navier--Stokes equations and backward
uniqueness}, Russian Mathematical Surveys \textbf{58} (2003), 211--250,
\url{https://doi.org/10.1070/RM2003v058n02ABEH000609}.

\end{thebibliography}

% =============================================================================
\section{V2: the second Oseen constant}
\label{sec:c4v2-entry}
% =============================================================================

The first acquisition of the fourth cycle (Programme~\ref{prog:c4v1-first})
is carried out in two steps; this version computes and certifies
the second Oseen constant \(c_{O,2}\), V3 proves the gradient bound with
the initial term through the ancient form and the new floor.  The
second derivative of the Oseen kernel is written in radial form
with three coefficient functions, the map it defines on symmetric
traceless matrices splits into three orthogonal channels whose gains
are explicit, the tail constant is \(36\) (the norm of \(\nabla^4(4\pi r)^{-1}\) on
axial inputs), and the midpoint sum with a Fourier-side Lipschitz
constant certifies \(c_{O,2}\le1.788\).  A note on the companion state at
the closure of the third cycle is recorded.  No global regularity
claim is made.  The next compulsory companion audit is V5.

\begin{remark}[Companion state at the closure]
\label{rem:c4v2-companion}
The V100 audit of the third cycle recorded SM V51 (digest
\texttt{699c607e\dots}, 20468 lines).  By the time the third cycle's file
was frozen the SM file had moved to V52 (digest \texttt{75f815d2\dots},
20941 lines); the frozen file is not edited, and the state is
recorded here for the V5 audit.
\end{remark}

% =============================================================================
\section{The second derivative of the Oseen kernel}
\label{sec:c4v2-radial}
% =============================================================================

Notation of the third cycle's V82: \(r=|x|\), \(n=x/r\), \(\phi(r)=\operatorname{erf}(r/2)/(4\pi r)\),
\(E=\operatorname{erf}(r/2)\), \(g=e^{-r^2/4}/\sqrt\pi\), \(K_1=(4\pi)^{-3/2}e^{-r^2/4}\), \(c=\phi''/r-\phi'/r^2\).  Put
\begin{equation}
 e:=\frac{\phi'''-3c}r,\qquad A:=\phi''''-6e-\frac{3c}r,\qquad
 \phi''''=\frac{E''''r^4-4E'''r^3+12E''r^2-24E'r+24E}{4\pi r^5},\quad E''''=\Bigl(\frac{3r}4-\frac{r^3}8\Bigr)g .
 \label{eq:c4v2-coefficients}
\end{equation}

\begin{lemma}[Radial form of the fourth derivative]
\label{lem:c4v2-fourth}
With \(\Sigma_6:=\delta_{il}n_jn_k+\delta_{jl}n_in_k+\delta_{kl}n_in_j+\delta_{ik}n_jn_l+\delta_{jk}n_in_l+\delta_{ij}n_kn_l\) and
\(\Sigma_3:=\delta_{ik}\delta_{jl}+\delta_{jk}\delta_{il}+\delta_{ij}\delta_{kl}\),
\begin{equation}
 \partial_i\partial_j\partial_k\partial_l\Phi_1=A\,n_in_jn_kn_l+e\,\Sigma_6+\frac cr\,\Sigma_3,\qquad
 \partial_k\partial_l(K_1\delta_{ij})=K_1\Bigl(\frac{r^2}4n_kn_l-\frac12\delta_{kl}\Bigr)\delta_{ij} .
 \label{eq:c4v2-tensor}
\end{equation}
The coefficients are smooth on \([0,\infty)\); as \(r\to\infty\),
\(4\pi r^5(A,e,c/r)\to(105,-15,3)\), the coefficients of \(\nabla^4(4\pi r)^{-1}\).
\end{lemma}

\begin{proof}
Differentiate the third derivative \(\partial_i\partial_j\partial_k\Phi_1=(\phi'''-3c)n_in_jn_k+c\,S_{ijk}\)
(V82, \(S_{ijk}=\delta_{ik}n_j+\delta_{jk}n_i+\delta_{ij}n_k\)) once more with \(\partial_ln_i=(\delta_{il}-n_in_l)/r\):
the \(n^{\otimes3}\) term gives \((\phi''''-3c')n^{\otimes4}+(\phi'''-3c)r^{-1}(\delta_{il}n_jn_k+\delta_{jl}n_in_k+\delta_{kl}n_in_j-3n^{\otimes4})\),
the \(S\) term gives \(c'n_lS_{ijk}+cr^{-1}(\delta_{ik}(\delta_{jl}-n_jn_l)+\delta_{jk}(\delta_{il}-n_in_l)+\delta_{ij}(\delta_{kl}-n_kn_l))\).
The identity \(c'=(\phi'''-2c)/r\) (direct differentiation of \(c\)) makes the
six \(\delta nn\) coefficients equal to \(e=(\phi'''-3c)/r=c'-c/r\), and the \(n^{\otimes4}\)
coefficient is \(\phi''''-3c'-3(\phi'''-3c)/r=\phi''''-6e-3c/r\).  The heat part is
\(\partial_k\partial_lK_1=(x_kx_l/4-\delta_{kl}/2)K_1\).  The formula for \(\phi''''\): the numerator
\(N_3=E'''r^3-3E''r^2+6E'r-6E\) of \(\phi'''\) has \(N_3'=E''''r^3\), so
\(\phi''''=(E''''r^4-4N_3)/(4\pi r^5)\).  Asymptotics: \(\phi\to(4\pi r)^{-1}\) with Gaussian
corrections, giving \(c\to3/(4\pi r^4)\), \(e\to-15/(4\pi r^5)\), \(A\to105/(4\pi r^5)\).
\end{proof}

% =============================================================================
\section{The map and its exact norm}
\label{sec:c4v2-norm}
% =============================================================================

Let \(M_2(x)\) be the map from \(3\times3\) matrices to \(3\times3\) matrices,
\(M_2(x)f:=\bigl(\partial_m\partial_j(K^O_1)_{il}(x)f_{lj}\bigr)_{im}\), so that \(\nabla e^{\Delta}\PP\operatorname{div}f=M_2*f\),
and \(\|M_2(x)\|_{s0}\) its norm on symmetric traceless \(f\) (Frobenius to
Frobenius).

\begin{proposition}[Three channels]
\label{prop:c4v2-channels}
For symmetric traceless \(f=\mu\,n\otimes n+n\otimes w+w\otimes n+B\) with \(w\perp n\), \(B\) supported
on \(n^\perp\), \(\operatorname{tr}B=-\mu\), \(B=-\frac\mu2I_\perp+B_0\) (\(B_0\) traceless on \(n^\perp\)):
\begin{align}
 M_2f&=\mu\Bigl[A+5e+\frac{2c}r+K_1\Bigl(\frac{r^2}4-\frac12\Bigr)\Bigr]n\otimes n+\mu\Bigl[e-\frac cr+\frac{K_1}4\Bigr]I_\perp
 +\alpha\,(n\otimes w+w\otimes n)+\beta\,w\otimes n+\Bigl(\frac{2c}r-\frac{K_1}2\Bigr)B_0,
 \label{eq:c4v2-map}\\
 \alpha&:=2e+\frac{2c}r-\frac{K_1}2,\qquad \beta:=\frac{K_1r^2}4 ,\notag
\end{align}
and the three channels (\(\mu\), \(w\), \(B_0\)) are orthogonal in input and
output, so that
\begin{equation}
 \|M_2(x)\|_{s0}^2=\max\Bigl\{\frac23\Bigl(\Bigl[A+5e+\frac{2c}r+K_1\Bigl(\frac{r^2}4-\frac12\Bigr)\Bigr]^2+2\Bigl[e-\frac cr+\frac{K_1}4\Bigr]^2\Bigr),\
 \frac{\alpha^2+(\alpha+\beta)^2}2,\ \Bigl(\frac{2c}r-\frac{K_1}2\Bigr)^2\Bigr\}.
 \label{eq:c4v2-opnorm}
\end{equation}
\end{proposition}

\begin{proof}
Contract \eqref{eq:c4v2-tensor} (indices \(i,l\) of the kernel and \(m,j\) of
the derivatives) with \(f_{lj}\): \(n^{\otimes4}\) gives \(\mu n_in_m\); \(\Sigma_6\) gives
\(2n_m(fn)_i+2n_i(fn)_m+\mu\delta_{im}\) (the term \(\delta_{lj}n_in_m\operatorname{tr}f\) vanishes); \(\Sigma_3\)
gives \(2f_{im}\); the heat part gives \(K_1(\frac{r^2}4n_m(fn)_i-\frac12f_{im})\).  Insert
\(fn=\mu n+w\), \(\delta_{im}=n_in_m+(I_\perp)_{im}\) and \(f=\mu n\otimes n+n\otimes w+w\otimes n-\frac\mu2I_\perp+B_0\) and
collect.  Input norms: \(|\mu n\otimes n-\frac\mu2I_\perp|^2=\frac32\mu^2\), \(|n\otimes w+w\otimes n|^2=2|w|^2\),
\(|B_0|^2\); output norms: \(|n\otimes n|=1\), \(|I_\perp|^2=2\), \(n\otimes n\perp I_\perp\);
\(|\alpha(n\otimes w+w\otimes n)+\beta w\otimes n|^2=(\alpha^2+(\alpha+\beta)^2)|w|^2\); the three output
subspaces are orthogonal.  The norm of an orthogonal direct sum is
the maximum of the channel norms.
\end{proof}

\begin{remark}[Verification]
\label{rem:c4v2-verification}
\eqref{eq:c4v2-tensor} was checked against central differences of the
third-derivative tensor of V82 (agreement to \(10^{-11}\) on a scale of
\(10^{-2}\)), and \eqref{eq:c4v2-opnorm} against the largest singular value
of the numerically assembled \(9\times5\) matrix of \(M_2\) on an orthonormal
basis of symmetric traceless matrices at eight radii (agreement to
all printed digits).  The \(B_0\) channel dominates for \(r\lesssim1.2\), the \(w\)
channel for \(1.2\lesssim r\lesssim3.05\), the axial channel beyond.
\end{remark}

\begin{lemma}[Tail]
\label{lem:c4v2-tail}
For \(r\ge12\), \(\|M_2(x)\|_{s0}\le(36+10^{-9})/(4\pi r^5)\), hence
\(\int_{|x|\ge R}\|M_2\|_{s0}\dd x\le(36+10^{-9})/(2R^2)\) for \(R\ge12\).
\end{lemma}

\begin{proof}
For \(\phi=(4\pi r)^{-1}\) the map is
\(4\pi r^5M_2f=51\mu\,n\otimes n-15\mu I-24(n\otimes w+w\otimes n)+6B\) (from
\((A,e,c/r)=(105,-15,3)/(4\pi r^5)\) and no heat part), and with \(B=-\frac\mu2I_\perp+B_0\):
\(36\mu\,n\otimes n-18\mu I_\perp-24(n\otimes w+w\otimes n)+6B_0\); the channel gains are
\(\frac23(36^2+2\cdot18^2)=1296\), \(\frac{2\cdot24^2}2=576\), \(36\), so the norm is \(36/(4\pi r^5)\),
attained on the axial input.  The corrections to \((A,e,c/r)\) and the
heat part are Gaussian: with \(E=1-\operatorname{erfc}(r/2)\), \(\operatorname{erfc}(r/2)\le\frac2rg\), every
coefficient of \eqref{eq:c4v2-coefficients} equals its \((4\pi r)^{-1}\) value
plus a term bounded by a polynomial of degree at most \(7\) in \(r\)
times \(g\), which at \(r=12\) is below \(10^{-11}\) relative and decreasing;
so the gains exceed their limits by less than \(10^{-9}\) relative.
Integrate \(4\pi r^2\cdot36/(4\pi r^5)\).
\end{proof}

% =============================================================================
\section{Certification}
\label{sec:c4v2-certification}
% =============================================================================

\begin{lemma}[Sup-norm and Lipschitz bounds]
\label{lem:c4v2-lipschitz}
\(\sup_x|\nabla^2K^O_1|_F\le\frac{3\sqrt2\pi^{1/2}}{16\pi^2}=0.04763\), \(\sup_x|\nabla^3K^O_1|_F\le\frac{\sqrt2}{2\pi^2}=0.07166\), and
\(G_2(r):=4\pi r^2\|M_2(rn)\|_{s0}\) is Lipschitz on \([0,12]\) with constant
\(L_2\le8\pi\cdot12\cdot0.04763+4\pi\cdot144\cdot0.07166<144.1\).
\end{lemma}

\begin{proof}
The transforms are \((i\xi)^{\otimes2}(\delta-\xi\xi/\rho^2)e^{-\rho^2}\) and \((i\xi)^{\otimes3}(\dots)\), of Frobenius
norms \(\sqrt2\rho^2e^{-\rho^2}\) and \(\sqrt2\rho^3e^{-\rho^2}\); \((2\pi)^{-3}4\pi\sqrt2\int\rho^4e^{-\rho^2}=\frac{3\sqrt2\pi^{1/2}}{16\pi^2}\)
and \((2\pi)^{-3}4\pi\sqrt2\int\rho^5e^{-\rho^2}=\frac{\sqrt2}{2\pi^2}\) (\(\int_0^\infty\rho^5e^{-\rho^2}=1\)).  Then as in
Lemma C3-V89-lipschitz.
\end{proof}

\begin{theorem}[Certified second Oseen constant]
\label{thm:c4v2-certified}
\(c_{O,2}:=\sup_tt\,\|\nabla e^{t\Delta}\PP\operatorname{div}\|_{s0,\infty\to\infty}\le\int\|M_2\|_{s0}\dd x\le1.788\).
\end{theorem}

\begin{proof}
Scaling gives \(\|\nabla^2K^O_t\|=t^{-1}\|\nabla^2K^O_1\|\) in \(L^1\), and the Cauchy--Schwarz
argument of Theorem C3-V84-gamma2 bounds the action on symmetric
traceless matrix fields by \(\int\|M_2\|_{s0}\) times the sup norm.  The
midpoint sum of \(G_2\) with \(2\cdot10^6\) cells on \([0,12]\) (formulas
\eqref{eq:c4v2-coefficients}, \eqref{eq:c4v2-opnorm} in double precision) is
\(1.660159\), with error at most \(3L_2h=0.00259\); the tail from \(12\) is at
most \(36.000000001/288=0.125000\); hence \(c_{O,2}\le1.660159+0.00259+0.125001<1.788\).
\end{proof}

\begin{assessment}[Reading]
\label{ass:c4v2-reading}
\(c_{O,2}\) is smaller than \(c_O^{s0}\) (\(1.788\) against \(2.766\)) because the second
derivative of the Oseen kernel is more concentrated (its tail is
\(r^{-5}\)), and its L\(^1\) norm is dominated by \(r\le3\).  In the gradient bound
of V3 it multiplies \(\kappa^2\log\frac{1+\theta}\theta\), so short windows cost only a
logarithm.  No global regularity claim is made.
\end{assessment}

% =============================================================================
\section{Integrated V2 conclusion and next acquisition order}
\label{sec:c4v2-conclusion}
% =============================================================================

\begin{theorem}[What V2 proves]
\label{thm:c4v2-conclusion}
Relative to the three legacy cycles and V1: the radial form of the
fourth derivative (Lemma~\ref{lem:c4v2-fourth}), the three-channel norm
(Proposition~\ref{prop:c4v2-channels}), the tail (Lemma~\ref{lem:c4v2-tail}),
the Lipschitz bound (Lemma~\ref{lem:c4v2-lipschitz}) and the certified
constant (Theorem~\ref{thm:c4v2-certified}).  The full Navier--Stokes
stability/global-regularity theorem remains unproved.
\end{theorem}

\begin{proof}
The cited results.
\end{proof}

\begin{programme}[V3 acquisition order]
\label{prog:c4v3}
\begin{enumerate}[label=\textup{(AC\arabic*)},leftmargin=4.8em]
\item The gradient bound with the initial term through the ancient
      form: for \(s<0\), \(L=-s\), \(s_0=s-\theta L\), split the ancient integral at
      \(s_0\); on \((-\infty,s_0)\) both derivatives on the kernel (constant
      \(c_{O,2}\), input \((W\otimes W)_0\)), on \([s_0,s]\) one derivative on \(W\otimes W\)
      (constant \(c_O^{s0}\)); the window integrals
      \(\int_{-\infty}^{s_0}(s-\tau)^{-1}(-\tau)^{-1}=L^{-1}\log\frac{1+\theta}\theta\) and \(2L^{-1}(\theta/(1+\theta))^{1/2}\);
      the absorption and the bound
      \(C_\infty'\le\gamma_1^{\rm anc}(\kappa,\theta)\kappa\),
      \(\gamma_1^{\rm anc}=(2/3)^{1/2}c_{O,2}\kappa\log\frac{1+\theta}\theta\big/\bigl(1-2(8/3)^{1/2}c_O^{s0}\kappa(\theta/(1+\theta))^{1/2}\bigr)\).
\item The certified floor of \((2/3)^{1/2}C_\infty'+\frac{\kappa^2}2\ge1\) with
      \(\gamma_1^{\rm anc}\) (expected \(0.344\)), its \(\theta^*\), the shares, and the
      comparison with the third cycle's \(0.1815\).
\item The next compulsory companion audit is V5.
\end{enumerate}
\end{programme}

\begin{assessment}[Position after V2]
\label{ass:c4v2-position}
The second Oseen constant is certified.  No full stability or
global-regularity claim is made.
\end{assessment}

% =============================================================================
\section*{Version V2 (fourth cycle) append-only record}
\addcontentsline{toc}{section}{Version V2 (fourth cycle) append-only record}
% =============================================================================

The complete V1 source of the fourth cycle was retained before this
continuation.  Its SHA--256 digest was
\begin{center}\scriptsize\ttfamily
850646d9735947da08c92544d09a686ce26884405611b72894af48fbb7972bc8.
\end{center}
V2 certified the second Oseen constant.  No full regularity claim is
made.

% =============================================================================
\section{V3: the gradient bound through the ancient form, and the floor \(0.3438\)}
\label{sec:c4v3-entry}
% =============================================================================

Items (AC1) and (AC2) of Programme~\ref{prog:c4v3} are carried out.
The ancient form of the mild formulation is split at a window
point; on the far part both derivatives fall on the Oseen kernel
(constant \(c_{O,2}\), input \((W\otimes W)_0\)), on the window one derivative falls on
the tensor (constant \(c_O^{s0}\)).  The resulting gradient bound has an
initial term quadratic in \(\kappa\) and logarithmic in the window
parameter, and the enstrophy constraint of the third cycle then
gives the certified floor \(\kappa\ge0.3438\) at every homogeneous type-I
point, against \(0.1815\) before.  An in-place correction of V2 is
recorded.  No global regularity claim is made.  The next compulsory
companion audit is V5.

\begin{remark}[In-place correction of V2]
\label{rem:c4v3-correction}
Remark C4-V2-verification (Remark~\ref{rem:c4v2-verification}) located
the crossover from the \(w\) channel to the axial channel at ``\(3.5\)'';
it is at \(3.05\) (and the first crossover at \(1.215\)).  Corrected in place;
the digest of V2 recorded below is that of the corrected file.
Nothing depends on the value.
\end{remark}

% =============================================================================
\section{The gradient bound}
\label{sec:c4v3-gradient}
% =============================================================================

Notation of the third cycle's V81 and V86: \(W\in\mathcal F(x_0)\), \(L=-s\),
\(s_0=s-\theta L\), \(\kappa=\nu^{-1/2}\sup(-s)^{1/2}\|W(s)\|_\infty\), \(C_\infty'=\sup(-s)\|\nabla W(s)\|_\infty\) (Frobenius),
\((W\otimes W)_0\) the traceless part, \(c_O^{s0}\le2.766\) (Theorem C3-V89-certified),
\(c_{O,2}\le1.788\) (Theorem~\ref{thm:c4v2-certified}).

\begin{lemma}[The two window integrals]
\label{lem:c4v3-integrals}
For \(s<0\), \(L=-s\), \(\theta>0\):
\begin{equation}
 \int_{-\infty}^{s_0}\frac{\dd\tau}{(s-\tau)(-\tau)}=\frac1L\log\frac{1+\theta}\theta,\qquad
 \int_{s_0}^s\frac{\dd\tau}{(s-\tau)^{1/2}(-\tau)^{3/2}}=\frac2L\Bigl(\frac\theta{1+\theta}\Bigr)^{1/2} .
 \label{eq:c4v3-integrals}
\end{equation}
\end{lemma}

\begin{proof}
With \(\tau=-Lr\): the first is \(L^{-1}\int_{1+\theta}^\infty\frac{\dd r}{(r-1)r}=L^{-1}[\log\frac{r-1}r]_{1+\theta}^\infty=L^{-1}\log\frac{1+\theta}\theta\);
the second is Theorem C3-V81-gradient's integral.
\end{proof}

\begin{theorem}[Gradient bound with the initial term through the ancient form]
\label{thm:c4v3-gradient}
For every \(\theta>0\) and every homogeneous type-I point with
\(2(8/3)^{1/2}c_O^{s0}\kappa(\theta/(1+\theta))^{1/2}<1\),
\begin{equation}
 \boxed{\quad
 C_\infty'\ \le\ \gamma_1^{\rm anc}(\kappa,\theta)\,\kappa,\qquad
 \gamma_1^{\rm anc}(\kappa,\theta):=\frac{(2/3)^{1/2}c_{O,2}\,\kappa\,\log\frac{1+\theta}\theta}{1-2(8/3)^{1/2}c_O^{s0}\,\kappa\,(\theta/(1+\theta))^{1/2}} .
 \quad}
 \label{eq:c4v3-gamma1}
\end{equation}
\end{theorem}

\begin{proof}
By the ancient form (Theorem C3-V81-ancient) with the traceless
reduction (Theorem C3-V86-absorptions),
\(W(s)=-\int_{-\infty}^se^{\nu(s-\tau)\Delta}\PP\operatorname{div}(W\otimes W)_0(\tau)\dd\tau\), the integral converging
absolutely in \(L^\infty\) together with its gradient (the kernel bounds
below).  Split at \(s_0\).  On \((-\infty,s_0)\), \(\nabla e^{\nu(s-\tau)\Delta}\PP\operatorname{div}\) has kernel
\(\nabla^2K^O_{\nu(s-\tau)}\), whose action on symmetric traceless matrix fields is
bounded by \(c_{O,2}(\nu(s-\tau))^{-1}\) times the sup norm (Theorem~\ref{thm:c4v2-certified},
with the Cauchy--Schwarz argument of Theorem C3-V84-gamma2 for the
Frobenius norm over all indices); with \(|(W\otimes W)_0|\le(2/3)^{1/2}C_\infty^2(-\tau)^{-1}\)
and \(C_\infty^2=\kappa^2\nu\), this part is at most
\((2/3)^{1/2}c_{O,2}\kappa^2\int_{-\infty}^{s_0}(s-\tau)^{-1}(-\tau)^{-1}\dd\tau=(2/3)^{1/2}c_{O,2}\kappa^2L^{-1}\log\frac{1+\theta}\theta\).
On \([s_0,s]\), the gradient falls on the tensor as in Theorem
C3-V86-absorptions(2): at most \(2(8/3)^{1/2}c_O^{s0}\kappa C_\infty'L^{-1}(\theta/(1+\theta))^{1/2}\).
Multiply by \(L\), take the supremum over \(s\), and absorb (\(C_\infty'<\infty\) by the
class).
\end{proof}

\begin{remark}[Comparison with the third cycle's initial term]
\label{rem:c4v3-comparison}
The third cycle bounded the initial term by \(c_h\kappa(\theta(1+\theta))^{-1/2}\): linear
in \(\kappa\), of order \(\theta^{-1/2}\) for short windows.  The new term
\((2/3)^{1/2}c_{O,2}\kappa^2\log\frac{1+\theta}\theta\) is quadratic in \(\kappa\) and grows only
logarithmically as \(\theta\to0\), so the window can be taken short enough
to make the absorption coefficient small, at a logarithmic cost.
At the floors of interest (\(\kappa\approx0.2\) to \(0.35\)) the new term is smaller
for every \(\theta<1\); both bounds hold, and the minimum may be taken.
\end{remark}

% =============================================================================
\section{The floor}
\label{sec:c4v3-floor}
% =============================================================================

Let \(\gamma_1^{\rm anc}(\kappa):=\inf_\theta\gamma_1^{\rm anc}(\kappa,\theta)\) over admissible \(\theta\), and
\(\Gamma^{\rm anc}(\kappa):=(2/3)^{1/2}\gamma_1^{\rm anc}(\kappa)\kappa+\frac{\kappa^2}2\), increasing in \(\kappa\).

\begin{theorem}[Certified floor with the ancient-form initial term]
\label{thm:c4v3-floor}
Every homogeneous type-I point of a smooth periodic solution
satisfies \(\Gamma^{\rm anc}(\kappa)\ge1\), hence
\begin{equation}
 \boxed{\quad
 \kappa\ \ge\ \kappa_1^{\rm anc}:=\min\{\kappa:\Gamma^{\rm anc}(\kappa)\ge1\}=0.3438
 \qquad(\theta^*=0.0094,\ \gamma_1^{\rm anc}=3.35,\ \gamma_1^{\rm anc}\kappa=1.152),
 \quad}
 \label{eq:c4v3-floor}
\end{equation}
i.e.\ \(C_\infty\ge0.3438\,\nu^{1/2}\).  The shares of the unit at the floor are
\((2/3)^{1/2}\gamma_1^{\rm anc}\kappa=0.941\) (strain) and \(\frac{\kappa^2}2=0.059\) (transport).  At the
floor the absorption coefficient is \(0.30\) and the initial term
\(0.81\).  Taking the minimum of the two initial terms changes nothing
at the floor.
\end{theorem}

\begin{proof}
Theorem C3-V93-rows(1), \((2/3)^{1/2}C_\infty'+\frac{\kappa^2}2\ge1\), with \eqref{eq:c4v3-gamma1}.
Monotonicity as in Theorem C3-V83-floor (for fixed \(\theta\), \(\gamma_1^{\rm anc}\kappa\) is
increasing in \(\kappa\), and the admissible set shrinks).  The values:
bisection on \(\kappa\) with the infimum over \(\theta\) on a geometric grid of
\(6\cdot10^4\) points in \([10^{-5},10^2]\), with \(c_{O,2}=1.788\), \(c_O^{s0}=2.766\); at the floor
\(2(8/3)^{1/2}c_O^{s0}\kappa(\theta^*/(1+\theta^*))^{1/2}=0.30\) and
\((2/3)^{1/2}c_{O,2}\kappa^2\log\frac{1+\theta^*}{\theta^*}=0.81\).
\end{proof}

\begin{theorem}[Item (AC2): the table]
\label{thm:c4v3-table}
Every homogeneous type-I point satisfies, all certified:

\medskip
\begin{center}\small
\begin{tabular}{@{}p{0.52\textwidth}p{0.14\textwidth}p{0.2\textwidth}@{}}
\toprule
\textbf{Statement} & \textbf{Value} & \textbf{Source}\\
\midrule
\(\kappa\ge1/(2(8/3)^{1/2}c_O^{s0})\) (gradient Liouville) & \(0.1107\) & C3-V83, V86, V89\\
\(\kappa\ge(3/2)^{1/2}/(\pi c_O^{s0})\) (velocity Liouville) & \(0.1410\) & C3-V81, V86, V89\\
\(\kappa\ge\kappa_1^{ST}\) (third cycle's initial term) & \(0.1815\) & C3-V94\\
\(\kappa\ge\kappa_1^{\rm anc}\) (ancient-form initial term; \(\theta^*=0.0094\)) & \(\mathbf{0.3438}\) & V2, V3\\
\(C_\infty'\le\gamma_1^{\rm anc}\kappa\) at the floor & \(1.152\) & V3\\
palinstrophy row and loop constraints & void for \(\kappa<1\) & C3-V93\\
\bottomrule
\end{tabular}
\end{center}
\medskip

\noindent
The floor exceeds the velocity Liouville floor by the factor \(2.44\)
and the third cycle's floor by \(1.89\).
\end{theorem}

\begin{proof}
Theorem C3-V94-table and Theorem~\ref{thm:c4v3-floor}.
\end{proof}

\begin{assessment}[Reading]
\label{ass:c4v3-reading}
The floor nearly doubled because the initial term, which the third
cycle bounded by the sup norm of \(W\) at the window's start, is in
fact the far part of the ancient integral and is small for small
velocity constants like \(\kappa^2\).  At the new floor the strain share is
\(0.94\) and the gradient bound is again what decides; its two parts are
now the far term (\(0.81\)) and the absorption (\(0.30\)), and the window is
short (\(\theta^*\approx0.01\)).  What remains in the sup-norm argument: the
second-derivative constant through the same route (a third Oseen
constant \(c_{O,3}\)), which would sharpen \(\nu^{1/2}C_\infty''\) and the concentration
bound but not the floor; and the two informations the argument
discards (directions (D2) of the third cycle's V99).  The reading of
the number: a homogeneous type-I singularity of a smooth periodic
solution, if one exists, has rescaled velocity bound at least
\(0.34\,\nu^{1/2}\); it is not excluded.
\end{assessment}

% =============================================================================
\section{Integrated V3 conclusion and next acquisition order}
\label{sec:c4v3-conclusion}
% =============================================================================

\begin{theorem}[What V3 proves]
\label{thm:c4v3-conclusion}
Relative to the three legacy cycles and V1--V2: the window integrals
(Lemma~\ref{lem:c4v3-integrals}), the gradient bound through the ancient
form (Theorem~\ref{thm:c4v3-gradient}), the certified floor
(Theorem~\ref{thm:c4v3-floor}) and the table (Theorem~\ref{thm:c4v3-table}).
The full Navier--Stokes stability/global-regularity theorem remains
unproved.
\end{theorem}

\begin{proof}
The cited results.
\end{proof}

\begin{programme}[V4 acquisition order]
\label{prog:c4v4}
\begin{enumerate}[label=\textup{(AC\arabic*)},leftmargin=4.8em]
\item Re-examine the window part of the bound: on \([s_0,s]\) the
      estimate puts one derivative on \(W\otimes W\) and uses \(C_\infty'\) at the
      full class value; record whether the far-part idea applies
      once more (a second split of the window) and what it gives,
      or the negative.
\item The second-derivative bound through the ancient form: the
      constant \(c_{O,3}\) (third derivative of the Oseen kernel), its radial
      form and tail (\(r^{-6}\), constant from \(\nabla^5(4\pi r)^{-1}\)), and
      \(\nu^{1/2}C_\infty''\) at the floor; the concentration bound \(12\nu C\) at the
      new floor (\(\theta_1^{ST}\) at \(0.3438\)).
\item The next compulsory companion audit is V5.
\end{enumerate}
\end{programme}

\begin{assessment}[Position after V3]
\label{ass:c4v3-position}
The certified floor is \(0.3438\).  No full stability or
global-regularity claim is made.
\end{assessment}

% =============================================================================
\section*{Version V3 (fourth cycle) append-only record}
\addcontentsline{toc}{section}{Version V3 (fourth cycle) append-only record}
% =============================================================================

The complete V2 source of the fourth cycle (after the in-place
correction of Remark~\ref{rem:c4v3-correction}) was retained before
this continuation.  Its SHA--256 digest was
\begin{center}\scriptsize\ttfamily
ce45b8c8d1c2d0126e99d88281da5a4a7651ca26212de93062f6b2826f16ea84.
\end{center}
V3 proved the gradient bound through the ancient form and the floor
\(0.3438\).  No full regularity claim is made.

% =============================================================================
\section{V4: the third Oseen constant and the second-derivative bound through the ancient form}
\label{sec:c4v4-entry}
% =============================================================================

Items (AC1) and (AC2) of Programme~\ref{prog:c4v4} are carried out.
(AC1): a second split of the window gives nothing, the window part
being already an absorption at the class supremum; recorded as a
negative.  (AC2): the third derivative of the Oseen kernel is written
in radial form (three coefficient functions, with the identities that
make it symmetric), the map on symmetric traceless inputs has the
axial channel as its norm at every radius, the tail constant is
\(60\sqrt{15}\) (the axial norm of \(\nabla^5(4\pi r)^{-1}\)), and the constant
\(c_{O,3}\le1.897\) is certified.  The second-derivative bound through the
ancient form gives \(\nu^{1/2}C_\infty''\le6.62\) at the floor \(0.3438\), against \(45\) by
the third cycle's route; the concentration bound at the floor is
\(12\nu C\).  No global regularity claim is made.  The next compulsory
companion audit is V5.

\begin{firewall}[Item (AC1): no second split]
\label{fw:c4v4-no-second-split}
On \([s_0,s]\) the gradient of the Duhamel integral is bounded by
\(2(8/3)^{1/2}c_O^{s0}\kappa(\theta/(1+\theta))^{1/2}\sup_\tau(-\tau)\|\nabla W(\tau)\|_\infty\), the supremum being the
class constant \(C_\infty'\) itself; splitting \([s_0,s]\) again and treating the
earlier piece with both derivatives on the kernel fails, the kernel
of \(\nabla e^{t\Delta}\PP\operatorname{div}\) having \(L^1\) norm \(c_{O,2}/t\), not integrable at \(t=0\).  The
bound of Theorem~\ref{thm:c4v3-gradient} is the fixed point of the
sup-norm iteration, and the window parameter is the only freedom;
it is optimized.  Nothing to lift within the sup-norm argument.
\end{firewall}

% =============================================================================
\section{The third derivative of the Oseen kernel}
\label{sec:c4v4-radial}
% =============================================================================

Notation of V2: \(c\), \(e\), \(A\), \(\phi''''\), \(K_1\).  Put
\begin{equation}
 A_5:=\phi'''''-\frac{10A}r-\frac{15e}r,\qquad B_5:=\frac Ar,\qquad D_5:=\frac er,\qquad
 \phi'''''=\frac{E'''''r^5-5N_4}{4\pi r^6},\quad N_4:=E''''r^4-4E'''r^3+12E''r^2-24E'r+24E,\quad
 E'''''=\Bigl(\frac34-\frac{3r^2}4+\frac{r^4}{16}\Bigr)g .
 \label{eq:c4v4-coefficients}
\end{equation}

\begin{lemma}[Radial form of the fifth derivative]
\label{lem:c4v4-fifth}
With \(\Sigma_{10}\) the sum over the ten pairs of \(\{i,j,k,l,p\}\) of \(\delta_{\rm pair}\) times
the product of \(n\) over the other three indices, and \(\Sigma_{15}\) the sum over
the fifteen ways of choosing two disjoint pairs of \(\delta\)'s times \(n\) of the
remaining index,
\begin{equation}
 \partial_i\partial_j\partial_k\partial_l\partial_p\Phi_1=A_5\,n_in_jn_kn_ln_p+B_5\,\Sigma_{10}+D_5\,\Sigma_{15},\qquad
 \partial_k\partial_l\partial_p(K_1\delta_{ij})=K_1\Bigl(\frac r4(\delta_{kp}n_l+\delta_{lp}n_k+\delta_{kl}n_p)-\frac{r^3}8n_kn_ln_p\Bigr)\delta_{ij} .
 \label{eq:c4v4-tensor}
\end{equation}
As \(r\to\infty\), \(4\pi r^6(A_5,B_5,D_5)\to(-945,105,-15)\).
\end{lemma}

\begin{proof}
Differentiate \eqref{eq:c4v2-tensor}: the \(n^{\otimes4}\) term gives
\(A'n_pn^{\otimes4}+(A/r)(\text{four }\delta_{\cdot p}nnn-4n^{\otimes5})\); the \(\Sigma_6\) term gives
\(e'n_p\Sigma_6+(e/r)(\text{twelve }\delta\delta n\text{ with }p\text{ in a }\delta-2\cdot\text{six }\delta nnn\text{ with }p\text{ not in the }\delta)\);
the \(\Sigma_3\) term gives \((c/r)'n_p\Sigma_3\).  Symmetry requires the \(\delta nnn\)
coefficients \(A/r\) (\(p\) in the \(\delta\)) and \(e'-2e/r\) (\(p\) not in it) to agree,
and the \(\delta\delta n\) coefficients \(e/r\) and \((c/r)'\) to agree; both are the
identities \(e'=A/r+2e/r\) and \((c/r)'=e/r\), which follow from
\(c'=(\phi'''-2c)/r\) (V82) and the definitions.  Then
\(A_5=A'-4A/r=\phi'''''-6e'-3(c/r)'-4A/r=\phi'''''-10A/r-15e/r\).  The heat part by
direct differentiation of \(\partial_k\partial_lK_1=(x_kx_l/4-\delta_{kl}/2)K_1\).  \(\phi'''''\): the
numerator \(N_4\) of \(\phi''''\) has \(N_4'=E'''''r^4\).  Asymptotics from
\(\phi\to(4\pi r)^{-1}\): \(\phi'''''\to-120/(4\pi r^6)\), \(A\to105/(4\pi r^5)\), \(e\to-15/(4\pi r^5)\), giving
\(A_5\to(-120-1050+225)/(4\pi r^6)\).
\end{proof}

\begin{proposition}[The map \(M_3\) and its norm]
\label{prop:c4v4-channels}
Let \(M_3(x)f:=\bigl(\partial_m\partial_p\partial_j(K^O_1)_{il}f_{lj}\bigr)_{imp}\), so that \(\nabla^2e^{\Delta}\PP\operatorname{div}f=M_3*f\).  On
symmetric traceless \(f\) the three channels of
Proposition~\ref{prop:c4v2-channels} (axial, \(w\), \(B_0\)) are again
invariant and mutually orthogonal in input and output, so
\(\|M_3(x)\|_{s0}\) is the largest of the three channel gains, each an
explicit linear combination of \((A_5,B_5,D_5,K_1r,K_1r^3)\) in norm; the
axial channel is the largest at every \(r\).
\end{proposition}

\begin{proof}
Rotations about \(n\) act on the input and output spaces, and the
tensor \eqref{eq:c4v4-tensor} is invariant; the channels are the
isotypic components (weights \(0,\pm1,\pm2\)), mapped into the
corresponding components of the output, whose weight decomposition
is orthogonal.  The gains are evaluated by contracting
\eqref{eq:c4v4-tensor} with a unit representative of each channel
(\(f_\mu=\operatorname{diag}(-\frac12,-\frac12,1)(3/2)^{-1/2}\), \(f_w=(e_1\otimes e_3+e_3\otimes e_1)/\sqrt2\),
\(f_B=(e_1\otimes e_1-e_2\otimes e_2)/\sqrt2\) for \(n=e_3\)); the contraction is linear in the
five coefficients with rational weights, which are recorded in the
working notes.  The dominance of the axial channel was checked on a
grid and against the largest singular value of the \(27\times5\) matrix of
\(M_3\) on an orthonormal basis of symmetric traceless matrices at four
radii (agreement to all printed digits).
\end{proof}

\begin{lemma}[Tail]
\label{lem:c4v4-tail}
For \(r\ge12\), \(\|M_3(x)\|_{s0}\le(60\sqrt{15}+10^{-8})/(4\pi r^6)\), hence
\(\int_{|x|\ge R}\|M_3\|_{s0}\le(60\sqrt{15}+10^{-8})/(3R^3)\).
\end{lemma}

\begin{proof}
For \(\phi=(4\pi r)^{-1}\), the axial gain squared of the tensor with
coefficients \((-945,105,-15)\) is \(54000\), so the norm is \(60\sqrt{15}/(4\pi r^6)\);
the Gaussian corrections at \(r\ge12\) are below \(10^{-9}\) relative as in
Lemma~\ref{lem:c4v2-tail}.  Integrate \(4\pi r^2\cdot60\sqrt{15}/(4\pi r^6)\).
\end{proof}

\begin{theorem}[Certified third Oseen constant]
\label{thm:c4v4-certified}
\(c_{O,3}:=\sup_tt^{3/2}\|\nabla^2e^{t\Delta}\PP\operatorname{div}\|_{s0,\infty\to\infty}\le\int\|M_3\|_{s0}\dd x\le1.897\).
\end{theorem}

\begin{proof}
\(\sup_x|\nabla^3K^O_1|_F\le\sqrt2/(2\pi^2)=0.07166\) and
\(\sup_x|\nabla^4K^O_1|_F\le(2\pi)^{-3}4\pi\sqrt2\int_0^\infty\rho^6e^{-\rho^2}\dd\rho=\frac{15\sqrt2\pi^{1/2}}{32\pi^2}=0.1191\) from the
Fourier side; hence \(G_3(r)=4\pi r^2\|M_3(rn)\|_{s0}\) is Lipschitz on \([0,12]\) with
\(L_3\le8\pi\cdot12\cdot0.07166+4\pi\cdot144\cdot0.1191<237.1\).  The midpoint sum with
\(2\cdot10^6\) cells is \(1.847870\), error at most \(3L_3h=0.00427\); the tail from
\(12\) is at most \(60\sqrt{15}/(3\cdot12^3)+10^{-9}=0.044827\); total below \(1.897\).  The
action on symmetric traceless matrix fields is bounded by the
\(L^1\) norm of the pointwise norm by the Cauchy--Schwarz argument of
Theorem C3-V84-gamma2, and scaling gives \(t^{-3/2}\).
\end{proof}

% =============================================================================
\section{The second-derivative bound and the concentration at the floor}
\label{sec:c4v4-second}
% =============================================================================

\begin{lemma}[Third window integral]
\label{lem:c4v4-J3}
\(\int_{-\infty}^{s_0}(s-\tau)^{-3/2}(-\tau)^{-1}\dd\tau=L^{-3/2}J_3(\theta)\),
\(J_3(\theta):=2\theta^{-1/2}-\pi+2\arctan\theta^{1/2}\).
\end{lemma}

\begin{proof}
\(\tau=-Lr\), \(u=r-1\): \(\int_\theta^\infty u^{-3/2}(1+u)^{-1}\dd u\), whose derivative in \(\theta\) is
\(-\theta^{-3/2}(1+\theta)^{-1}\), matched by \(\frac{\dd}{\dd\theta}(2\theta^{-1/2}-\pi+2\arctan\theta^{1/2})\), with the
value \(0\) at \(\theta=\infty\).
\end{proof}

\begin{theorem}[Second-derivative bound through the ancient form]
\label{thm:c4v4-gamma2}
For \(\theta>0\) with \((8/3)^{1/2}c_O^{s0}\kappa J_2(\theta)<1\), every homogeneous type-I point
satisfies
\begin{equation}
 \boxed{\quad
 \nu^{1/2}C_\infty''\ \le\ \frac{(2/3)^{1/2}c_{O,3}\,\kappa^2J_3(\theta)+2c_O^{s0}\,\gamma_1^{\rm anc}(\kappa)^2\kappa^2J_2(\theta)}{1-(8/3)^{1/2}c_O^{s0}\,\kappa\,J_2(\theta)} ,
 \quad}
 \label{eq:c4v4-gamma2}
\end{equation}
and at the floor \(\kappa=0.3438\), \(\gamma_1^{\rm anc}=3.352\): \(\nu^{1/2}C_\infty''\le6.62\) (\(\theta=0.0105\); far term
\(3.04\), window term \(1.50\), absorption \(0.32\)), against \(45.5\) by the route of
Theorem C3-V86-absorptions(3) at the same \(\kappa\).  At the floor
\(\theta_1^{ST}=\frac12\), so \(\sup_{\mathcal F(x_0)}\int|z|^2|\Omega|^2G\le12\nu C\).
\end{theorem}

\begin{proof}
Ancient form split at \(s_0\); on \((-\infty,s_0)\) all three derivatives fall on
the kernel (\(c_{O,3}(\nu(s-\tau))^{-3/2}\) against \((2/3)^{1/2}\kappa^2\nu(-\tau)^{-1}\), integrated by
Lemma~\ref{lem:c4v4-J3}); on \([s_0,s]\) as in Theorem C3-V86-absorptions(3)
with \(C_\infty'\le\gamma_1^{\rm anc}\kappa\).  Multiply by \(\nu^{1/2}L^{3/2}\) and absorb.  Values by
minimization over \(\theta\) on a geometric grid.  The concentration
statement is Theorem C3-V93-rows(2) with \(\theta_1^{ST}=\frac32-(\text{gradient term}+\frac{\kappa^2}2)=\frac12\)
at the floor.
\end{proof}

\begin{theorem}[Item (AC2): the table]
\label{thm:c4v4-table}
Every homogeneous type-I point satisfies, all certified:

\medskip
\begin{center}\small
\begin{tabular}{@{}p{0.52\textwidth}p{0.14\textwidth}p{0.2\textwidth}@{}}
\toprule
\textbf{Statement} & \textbf{Value} & \textbf{Source}\\
\midrule
\(c_O^{s0}\), \(c_{O,2}\), \(c_{O,3}\) (certified) & \(2.766\), \(1.788\), \(1.897\) & C3-V89, V2, V4\\
\(\kappa\ge\kappa_1^{\rm anc}\) (enstrophy row, ancient-form initial term) & \(\mathbf{0.3438}\) & V3\\
\(C_\infty'\le\gamma_1^{\rm anc}\kappa\) at the floor & \(1.152\) & V3\\
\(\nu^{1/2}C_\infty''\) at the floor & \(6.62\) & V4\\
\(\sup\int|z|^2|\Omega|^2G\le12\nu C\) at the floor & --- & C3-V93, V4\\
palinstrophy row and loop constraints & void for \(\kappa<1\) & C3-V93\\
\bottomrule
\end{tabular}
\end{center}
\medskip
\end{theorem}

\begin{proof}
The cited results.
\end{proof}

\begin{assessment}[Reading]
\label{ass:c4v4-reading}
The three Oseen constants are of the same size (\(1.8\) to \(2.8\)) and
their tails are the axial norms of \(\nabla^k(4\pi r)^{-1}\): \(\sqrt{54}\), \(36\), \(60\sqrt{15}\).  The
ancient-form route improves every derivative constant by an order of
magnitude for small \(\kappa\), because it replaces the sup norm of \(W\) at
the window's start by the far part of the ancient integral.  The
floor itself is decided by the first-derivative bound and is now
\(0.3438\); the next gain would come from information the sup-norm
argument discards (directions (D2)), or from a sharper treatment of
the window part than \(C_\infty'\) at its supremum, which the negative record
above says is not available within the argument.  No global
regularity claim is made.
\end{assessment}

% =============================================================================
\section{Integrated V4 conclusion and next acquisition order}
\label{sec:c4v4-conclusion}
% =============================================================================

\begin{theorem}[What V4 proves]
\label{thm:c4v4-conclusion}
Relative to the three legacy cycles and V1--V3: the negative record
(Firewall~\ref{fw:c4v4-no-second-split}), the fifth derivative in radial
form (Lemma~\ref{lem:c4v4-fifth}), the map and its norm
(Proposition~\ref{prop:c4v4-channels}), the tail (Lemma~\ref{lem:c4v4-tail}),
the certified constant (Theorem~\ref{thm:c4v4-certified}), the window
integral (Lemma~\ref{lem:c4v4-J3}), the second-derivative bound
(Theorem~\ref{thm:c4v4-gamma2}) and the table (Theorem~\ref{thm:c4v4-table}).
The full Navier--Stokes stability/global-regularity theorem remains
unproved.
\end{theorem}

\begin{proof}
The cited results.
\end{proof}

\begin{programme}[V5 acquisition order]
\label{prog:c4v5}
\begin{enumerate}[label=\textup{(AC\arabic*)},leftmargin=4.8em]
\item The compulsory companion audit: read the current SM and YM
      notes and record their digests and decision.
\item Record the position after five versions and the plan for
      V6--V10: direction (D2) of the third cycle's V99 (the \(L^6\) and
      energy information), beginning with the far part of the
      kernel action in \(L^{3/2}\) against \(\|W\otimes W\|_{L^3}\le K_6^2(-\tau)^{-1/2}\); delivery.
\end{enumerate}
\end{programme}

\begin{assessment}[Position after V4]
\label{ass:c4v4-position}
Three Oseen constants are certified and the derivative bounds at
the floor are explicit.  No full stability or global-regularity
claim is made.
\end{assessment}

% =============================================================================
\section*{Version V4 (fourth cycle) append-only record}
\addcontentsline{toc}{section}{Version V4 (fourth cycle) append-only record}
% =============================================================================

The complete V3 source of the fourth cycle was retained before this
continuation.  Its SHA--256 digest was
\begin{center}\scriptsize\ttfamily
50353bc91ed3ae22716ac135c5983bf4106eef41c64e71072d4b65a5d7ec0963.
\end{center}
V4 certified the third Oseen constant and the second-derivative
bound.  No full regularity claim is made.

% =============================================================================
\section{V5: compulsory companion audit and the position after five versions}
\label{sec:c4v5-entry}
% =============================================================================

This is the fifth version of the fourth cycle.  The compulsory review
of the two companion programmes is performed first, and the position
after five versions is recorded with the plan for the block V6--V10.
No global regularity claim is made.  The next compulsory companion
audit is V10.

% =============================================================================
\section{Compulsory V5 review of the current companion notes}
\label{sec:c4v5-companion-audit}
% =============================================================================

The current companion files in \texttt{latex\_notes} were inspected
read-only.  Since the third cycle's V100 audit the SM programme has
moved from V51 to V53 of its seventh series (through V52,
Remark~\ref{rem:c4v2-companion}), the YM programme from V28 to V29 of
its sixth series.  The frozen checkpoints are
\begin{align}
 &\texttt{Renewal\_SM\_Einstein\_Continuum\_Research\_Notes\_V53.tex},
 \notag\\[-2pt]
 &\qquad
 \texttt{f99b1260d184d4d4f7aa257ca398f8816ef4f2c6ddfcdcaf4393c2a73f051c4b},
 \label{eq:c4v5-SM-checkpoint}\\
 &\texttt{Renewal\_Yang\_Mills\_Mass\_Gap\_Research\_Notes\_V29.tex},
 \notag\\[-2pt]
 &\qquad
 \texttt{32fef34da4b92b56cf90bfc0d14c28164fe92f6d3f7253d8dd2d5c9294f52e3d}.
 \label{eq:c4v5-YM-checkpoint}
\end{align}
(SM V53 has 21359 lines; YM V29 has 13056 lines.  The folder also
held YM V28 with digest \texttt{32c43d66\dots} and 12566 lines, which differs
from the state recorded at the third cycle's V100
(\texttt{8ab8b98f\dots}, 12565 lines): the YM programme edited its V28 after
that audit, consistently with the parent record in its V29; this is
recorded, and nothing depends on it.)

\emph{SM V52--V53.}  The SM programme treated a local jet-balanced
quantization and the heat scale (going upstream from the
growing-range completion, smooth fixed-jet stratum symbols, a
cancelled midpoint term, a gap without a conservative completion, a
differentiated irrelevant-jet ledger, what is missing from the heat
trace), then a stratified mesoscopic heat transfer (uniform overlap
ellipticity on the fixed-jet family, a phase-space heat split, a
low-branch parametrix with a shrinking profile, an explicit
mesoscopic endpoint, the remaining coefficient ownership problem).

\emph{YM V29.}  The YM programme treated vacuum-fillable collars and
boundary-stable defect suppression: an attained radius, translated
boxes, a conditional partition function and conditional
suppression with a fillability threshold, collar incidence and
response transfer, a boundary form of its multi-defect suppression
card with a threshold for its condition CDAC, and a decidability
statement.

\begin{theorem}[V5 companion-audit decision]
\label{thm:c4v5-companion-decision}
Nothing is imported from SM V53 or YM V29.  Neither bears on the
fourth cycle's first block.
\end{theorem}

\begin{proof}
The SM results concern heat-trace parametrices of lattice Dirac
interfaces; the YM results concern defect suppression in a lattice
gauge theory.  None is a Navier--Stokes statement, and none concerns
the mild formulation, the Oseen kernels or the explicit floors.
\end{proof}

% =============================================================================
\section{Position after five versions, and the block V6--V10}
\label{sec:c4v5-position}
% =============================================================================

\begin{assessment}[Position after five versions of the fourth cycle]
\label{ass:c4v5-position-five}
The first block (V2--V4) carried out direction (D1) of the third
cycle's V99: the second and third Oseen constants are certified
(\(c_{O,2}\le1.788\), \(c_{O,3}\le1.897\)), the initial term of the gradient bound is
the far part of the ancient integral, and the certified floor on
the velocity constant at a homogeneous type-I point rose from \(0.1815\)
to \(0.3438\), with \(C_\infty'\le1.152\) and \(\nu^{1/2}C_\infty''\le6.62\) at the floor and the
concentration \(\int|z|^2|\Omega|^2G\le12\nu C\).  Within the sup-norm argument
nothing remains (Firewall~\ref{fw:c4v4-no-second-split}).  The block
V6--V10 takes direction (D2): the far part of the kernel action on the
window, \(|x-y|>R\) in parabolic units, bounded through H\"older against
\(\|\,|W||\nabla W|\,\|_{L^{3/2}}\le K_6C^{1/2}\) (\(W\in L^6\), \(\nabla W\in L^2\)) instead of the sup
norm; the absorption coefficient then carries only the near part of
\(c_O^{s0}\) (the tail \(\sqrt{54}/R\) is most of it for \(R\lesssim5\)), and an explicit
inhomogeneous term in \(K_6C^{1/2}\), i.e.\ in the type-I enstrophy constant
\(C_I\), appears.  The outcome is a floor \(\kappa\ge F(C_I)\), to be compared with
\(0.3438\): V6 the \(L^3\) norm of \(\nabla K^O_t\) outside \(R\) and the near constant;
V7 the split bound and the constraint; V8 the floor as a function of
\(C_I\) and its reading; V9 consolidation; V10 audit.  No global
regularity claim is made.
\end{assessment}

% =============================================================================
\section{Integrated V5 conclusion and next acquisition order}
\label{sec:c4v5-conclusion}
% =============================================================================

\begin{theorem}[What V5 establishes]
\label{thm:c4v5-conclusion}
Relative to the three legacy cycles and V1--V4: the companion-audit
decision (Theorem~\ref{thm:c4v5-companion-decision}) and the position.
The full Navier--Stokes stability/global-regularity theorem remains
unproved.
\end{theorem}

\begin{proof}
The cited records.
\end{proof}

\begin{programme}[V6 acquisition order]
\label{prog:c4v6}
\begin{enumerate}[label=\textup{(AC\arabic*)},leftmargin=4.8em]
\item The near and far constants: \(c_O^{s0}(R):=\int_{|x|\le R}\|M\|_{s0}\) (certified by the
      midpoint sums of C3-V89 restricted to \([0,R]\)) and
      \(m_3(R):=\|\,\|M\|_{s0}\,\|_{L^3(|x|>R)}\), with its tail bound from
      \(\|M\|_{s0}\le\sqrt{54}/(4\pi r^4)\); values for \(R\in\{2,3,5,8,12\}\).
\item The H\"older step: for \(f=(\nabla(W\otimes W))_0\),
      \(\|f\|_{L^{3/2}}\le(8/3)^{1/2}\|W\|_{L^6}\|\nabla W\|_{L^2}\le(8/3)^{1/2}K_6C^{1/2}(-\tau)^{-3/4}\), and the far part of
      \(\nabla e^{\nu t\Delta}\PP\operatorname{div}\) in sup norm bounded by \(m_3(R)(\nu t)^{-1}\|f\|_{L^{3/2}}\) (scaling of
      the \(L^3\) norm of the kernel: \(t^{-1/2-3/2\cdot(1-1/3)}\dots\), to be written
      exactly); the window integral \(\int_{s_0}^s(s-\tau)^{-1}(-\tau)^{-3/4}\dd\tau\) in closed
      form.
\item The next compulsory companion audit is V10.
\end{enumerate}
\end{programme}

\begin{assessment}[Position after V5]
\label{ass:c4v5-position}
The companion audit is recorded; the far-field block begins.  No
full stability or global-regularity claim is made.
\end{assessment}

% =============================================================================
\section*{Version V5 (fourth cycle) append-only record}
\addcontentsline{toc}{section}{Version V5 (fourth cycle) append-only record}
% =============================================================================

The complete V4 source of the fourth cycle was retained before this
continuation.  Its SHA--256 digest was
\begin{center}\scriptsize\ttfamily
d3d60e128c9e5701516393f87f2484e349b3679352f67328e8a6d7c890ef8a48.
\end{center}
V5 performed the compulsory companion audit and recorded the
position.  No full regularity claim is made.

% =============================================================================
\section{V6: near and far Oseen constants, and the H\"older step}
\label{sec:c4v6-entry}
% =============================================================================

Items (AC1) and (AC2) of Programme~\ref{prog:c4v6} are carried out.
The Oseen gradient's action on the window is split in space at a
fixed radius: inside, the sup-norm bound with the near constant
\(c_O^{s0}(R)\), the integral of the pointwise norm over \(|x|\le R\); outside, a
H\"older bound with the \(L^3\) norm \(m_3(R)\) of the pointwise norm over
\(|x|>R\) against the \(L^{3/2}\) norm of the tensor, which the class bounds
by the Sobolev and enstrophy constants.  Both constants are certified
for the radii used, the scaling of the split is written exactly, and
the window integrals of the two parts are reduced to two explicit
functions \(A(\theta,\lambda)\), \(\Phi(\theta,\lambda)\) of the window parameter and the split
parameter.  No global regularity claim is made.  The next compulsory
companion audit is V10.

% =============================================================================
\section{The near and far constants}
\label{sec:c4v6-constants}
% =============================================================================

Notation of the third cycle's V82, V87, V89: \(M(x)\) the map
\(f\mapsto(\partial_j(K^O_1)_{il}(x)f_{lj})_i\), \(\|M\|_{s0}=\max\{\sqrt6|c|,|2c-h_K|/\sqrt2\}\) its norm on
symmetric traceless inputs, \(c_O^{s0}=\int\|M\|_{s0}\le2.766\).

\begin{definition}[Near and far constants]
\label{def:c4v6-constants}
\(c_O^{s0}(R):=\int_{|x|\le R}\|M(x)\|_{s0}\dd x\) and \(m_3(R):=\bigl(\int_{|x|>R}\|M(x)\|_{s0}^3\dd x\bigr)^{1/3}\), \(R\ge0\).
\end{definition}

\begin{proposition}[Certified values]
\label{prop:c4v6-values}
\medskip
\begin{center}\small
\begin{tabular}{@{}lccccccccc@{}}
\toprule
\(R\) & 1 & 2 & 3 & 4 & 5 & 6 & 8 & 10 & 12\\
\midrule
\(c_O^{s0}(R)\le\) & 0.0155 & 0.1764 & 0.5436 & 0.9613 & 1.2983 & 1.5403 & 1.8466 & 2.0306 & 2.1533\\
\(m_3(R)\le\) & 0.02579 & 0.02308 & 0.01612 & 0.00932 & 0.00516 & 0.00302 & 0.00128 & 0.00065 & 0.00038\\
\bottomrule
\end{tabular}
\end{center}
\medskip
and \(m_3(0)\le0.0260\).  For \(R\ge8\), \(m_3(R)\le\bigl(\sqrt{54}(1+10^{-4})\bigr)\,(4\pi)^{-2/3}9^{-1/3}R^{-3}\).
\end{proposition}

\begin{proof}
\(c_O^{s0}(R)\): the midpoint sum of Theorem C3-V89-certified restricted to
\([0,R]\), with the error bound \(3Lh\cdot R/12\), \(L\le97\).  \(m_3(R)\) for \(R\ge8\): the
pointwise bound \(\|M\|_{s0}\le\sqrt{54}(1+\varepsilon)/(4\pi r^4)\) holds for \(r\ge8\) with
\(\varepsilon\le(\frac{r^5}4+r^3+6r+\frac{12}r)g/6\le10^{-4}\) (proof of Lemma C3-V82-tail, evaluated at
\(r=8\)), and \(\int_R^\infty4\pi r^2(4\pi r^4)^{-3}\dd r=(4\pi)^{-2}/(9R^9)\).  For \(R<8\): the midpoint
sum of \(4\pi r^2\|M\|_{s0}^3\) on \([R,8]\) with \(2\cdot10^7\) cells (\(h=6\cdot10^{-7}\)), Lipschitz
constant \(8\pi\cdot8\,m^3+12\pi\cdot64\,m^2\cdot0.04763\) with \(m\) a certified supremum of
\(\|M\|_{s0}\) on \([R,8]\) (the grid maximum plus \(0.04763\,h/2\) for \(R\ge0.01\), the
Fourier bound \(0.0358\) of Lemma C3-V89-sup for \(R<0.01\)), error
\((8-R)Lh/4\), plus the analytic tail from \(8\); the error is below \(2\cdot10^{-7}\)
on the cube in every case (\(0.8\%\) on \(m_3(0)\), less elsewhere).
\end{proof}

% =============================================================================
\section{Scaling and the split}
\label{sec:c4v6-split}
% =============================================================================

\begin{lemma}[Scaling of the near and far parts]
\label{lem:c4v6-scaling}
For \(t>0\), \(\rho>0\), and a symmetric traceless matrix field \(F\) bounded and
in \(L^{3/2}\), with \(M_t(y)=t^{-2}M(y/t^{1/2})\) the kernel of \(e^{t\Delta}\PP\operatorname{div}\),
\begin{equation}
 \Bigl|\int_{|y|\le\rho}M_t(y)F(x-y)\dd y\Bigr|\le c_O^{s0}(\rho t^{-1/2})\,t^{-1/2}\|F\|_\infty,\qquad
 \Bigl|\int_{|y|>\rho}M_t(y)F(x-y)\dd y\Bigr|\le m_3(\rho t^{-1/2})\,t^{-3/2}\|F\|_{L^{3/2}} .
 \label{eq:c4v6-split}
\end{equation}
Both hold for the Frobenius norm over any additional indices of \(F\)
(Cauchy--Schwarz in the measure \(\|M_t\|\dd y\) for the first, Minkowski and
H\"older for the second).
\end{lemma}

\begin{proof}
\(\int_{|y|\le\rho}\|M_t(y)\|\dd y=t^{-2}t^{3/2}\int_{|z|\le\rho t^{-1/2}}\|M(z)\|\dd z=t^{-1/2}c_O^{s0}(\rho t^{-1/2})\), and Young.
\(\|\,\|M_t\|\,\|_{L^3(|y|>\rho)}=(t^{-6}t^{3/2}\int_{|z|>\rho t^{-1/2}}\|M\|^3)^{1/3}=t^{-3/2}m_3(\rho t^{-1/2})\), and H\"older with
exponents \(3\) and \(3/2\).  For several components \(F^{(m)}\),
\(\sum_m(\int\|M_t\||F^{(m)}|)^2\le(\int\|M_t\|\,|F|)^2\) with \(|F|^2=\sum_m|F^{(m)}|^2\).
\end{proof}

\begin{lemma}[The H\"older step]
\label{lem:c4v6-holder}
For \(W\in\mathcal F(x_0)\) and \(F:=(\nabla(W\otimes W))_0\) (traceless part in the tensor indices,
Frobenius norm over all indices),
\begin{equation}
 \|F(\tau)\|_{L^{3/2}}\le\Bigl(\frac83\Bigr)^{1/2}\|W(\tau)\|_{L^6}\|\nabla W(\tau)\|_{L^2}\le\Bigl(\frac83\Bigr)^{1/2}S_3\,C_I^2\,\nu^{3/2}(-\tau)^{-1/2},
 \label{eq:c4v6-holder}
\end{equation}
with \(S_3=(\pi^2/4)^{1/6}(3\pi^2/4)^{-1/2}=0.4273\) the sharp Sobolev constant of
\(\|f\|_{L^6(\R^3)}\le S_3\|\nabla f\|_{L^2}\) and \(C_I\) the type-I enstrophy constant
(\(\|\nabla W(s)\|_{L^2}^2\le C_I^2\nu^{3/2}(-s)^{-1/2}\)).
\end{lemma}

\begin{proof}
\(|F|\le(8/3)^{1/2}|W||\nabla W|\) pointwise (Lemma C3-V86-norms); H\"older with
\(\frac16+\frac12=\frac23\); \(\|W(\tau)\|_6\le S_3\|\nabla W(\tau)\|_2\) (Sobolev on \(\R^3\); the sharp
constant is attained by \((1+|x|^2)^{-1/2}\), for which \(\|U\|_6^6=\pi^2/4\) and
\(\|\nabla U\|_2^2=3\pi^2/4\)); then \(\|\nabla W(\tau)\|_2^2\le C_I^2\nu^{3/2}(-\tau)^{-1/2}\) twice.
\end{proof}

\begin{remark}[Which constants enter]
\label{rem:c4v6-constants}
The enstrophy constant \(C_I\) is the constant of the type-I rate (I) of
the class, not a derived quantity; \(K_6\) is replaced by \(S_3C_I\nu^{3/4}\) and does
not appear.  The far part therefore introduces the single class
constant \(C_I\), and the floor of V7 is a function of it.
\end{remark}

% =============================================================================
\section{The window functions}
\label{sec:c4v6-window}
% =============================================================================

\begin{lemma}[The two window functions]
\label{lem:c4v6-window}
For \(s<0\), \(L=-s\), \(\theta>0\), \(\lambda>0\), \(s_0=s-\theta L\) and \(\rho^2=\nu\lambda L\):
\begin{align}
 L\int_{s_0}^s c_O^{s0}\Bigl(\frac\rho{(\nu(s-\tau))^{1/2}}\Bigr)\frac{\dd\tau}{(s-\tau)^{1/2}(-\tau)^{3/2}}
 &=\int_0^\theta c_O^{s0}\Bigl(\sqrt{\lambda/u}\Bigr)\frac{\dd u}{u^{1/2}(1+u)^{3/2}}=:A(\theta,\lambda),
 \label{eq:c4v6-A}\\
 L\int_{s_0}^s m_3\Bigl(\frac\rho{(\nu(s-\tau))^{1/2}}\Bigr)\frac{\dd\tau}{(s-\tau)^{3/2}(-\tau)^{1/2}}
 &=\int_0^\theta m_3\Bigl(\sqrt{\lambda/u}\Bigr)\frac{\dd u}{u^{3/2}(1+u)^{1/2}}=:\Phi(\theta,\lambda).
 \label{eq:c4v6-Phi}
\end{align}
\(A\) is increasing in \(\lambda\) with limit \(2c_O^{s0}(\theta/(1+\theta))^{1/2}\) as \(\lambda\to\infty\); \(\Phi\) is
finite for \(\lambda>0\) (near \(u=0\), \(m_3(\sqrt{\lambda/u})\le c\,(u/\lambda)^{3/2}\)), decreasing in \(\lambda\), and
diverges as \(\lambda\to0\).
\end{lemma}

\begin{proof}
\(\tau=-L(1+u)\): \(s-\tau=Lu\), \(-\tau=L(1+u)\), \(\rho/(\nu(s-\tau))^{1/2}=(\lambda/u)^{1/2}\); the powers of
\(L\) cancel.  Monotonicity from Proposition~\ref{prop:c4v6-values}
(\(c_O^{s0}(R)\) increasing, \(m_3(R)\) decreasing); the limit is Lemma
C3-V81-gradient's integral.
\end{proof}

\begin{remark}[Evaluation]
\label{rem:c4v6-evaluation}
\(A\) and \(\Phi\) were evaluated on a logarithmic grid in \(u\) (\(6000\) points
between \(10^{-9}\min\{\lambda,\theta\}\) and \(\theta\)), with \(c_O^{s0}\) and \(m_3\) interpolated from the
certified tables on grids of step \(10^{-3}\) and \(5\cdot10^{-3}\) and the analytic
forms beyond \(12\) and \(8\); the sup-norm limit of \(A\) is reproduced to four
digits.  A uniform grid in \(u\) misses the integrable singularity of
\eqref{eq:c4v6-Phi} at \(u=0\) and was discarded.
\end{remark}

% =============================================================================
\section{Integrated V6 conclusion and next acquisition order}
\label{sec:c4v6-conclusion}
% =============================================================================

\begin{theorem}[What V6 proves]
\label{thm:c4v6-conclusion}
Relative to the three legacy cycles and V1--V5: the near and far
constants (Definition~\ref{def:c4v6-constants}, Proposition~\ref{prop:c4v6-values}),
the scaling of the split (Lemma~\ref{lem:c4v6-scaling}), the H\"older step
(Lemma~\ref{lem:c4v6-holder}) and the window functions
(Lemma~\ref{lem:c4v6-window}).  The full Navier--Stokes
stability/global-regularity theorem remains unproved.
\end{theorem}

\begin{proof}
The cited results.
\end{proof}

\begin{programme}[V7 acquisition order]
\label{prog:c4v7}
\begin{enumerate}[label=\textup{(AC\arabic*)},leftmargin=4.8em]
\item The split gradient bound: for \(\theta,\lambda>0\) with \((8/3)^{1/2}\kappa A(\theta,\lambda)<1\),
      \(C_\infty'\le\bigl[(2/3)^{1/2}c_{O,2}\kappa^2\log\frac{1+\theta}\theta+(8/3)^{1/2}S_3C_I^2\Phi(\theta,\lambda)\bigr]/\bigl(1-(8/3)^{1/2}\kappa A(\theta,\lambda)\bigr)\);
      the constraint \((2/3)^{1/2}C_\infty'+\frac{\kappa^2}2\ge1\) then gives the floor \(F(C_I)\) as
      the least \(\kappa\) satisfying it with the infimum over \((\theta,\lambda)\).
\item The curve \(F(C_I)\) and its inverse (the least \(C_I\) compatible
      with a given \(\kappa\)); the threshold beyond which \(F=0.3438\); the
      reading.
\item The next compulsory companion audit is V10.
\end{enumerate}
\end{programme}

\begin{assessment}[Position after V6]
\label{ass:c4v6-position}
The split is set up with certified constants.  No full stability or
global-regularity claim is made.
\end{assessment}

% =============================================================================
\section*{Version V6 (fourth cycle) append-only record}
\addcontentsline{toc}{section}{Version V6 (fourth cycle) append-only record}
% =============================================================================

The complete V5 source of the fourth cycle was retained before this
continuation.  Its SHA--256 digest was
\begin{center}\scriptsize\ttfamily
7a079b7c8eb8d22621fe980640b110828cfd9a6bdf2708c67a53626495ed7404.
\end{center}
V6 set up the near--far split with certified constants.  No full
regularity claim is made.

% =============================================================================
\section{V7: the split gradient bound and the floor as a function of the enstrophy constant}
\label{sec:c4v7-entry}
% =============================================================================

Items (AC1) and (AC2) of Programme~\ref{prog:c4v7} are carried out.
The near--far split of V6 gives a gradient bound whose absorption
carries only the near constant and whose inhomogeneous term is
explicit in the type-I enstrophy constant \(C_I\); the enstrophy
constraint then gives a floor \(F(C_I)\) on the velocity constant at
every homogeneous type-I point, equal to \(0.3438\) for \(C_I\ge1.93\) and rising
to \(\sqrt2\) as \(C_I\to0\); equivalently, a floor on \(C_I\) for each \(\kappa<\sqrt2\).  The
curve and its inverse are tabulated.  No global regularity claim is
made.  The next compulsory companion audit is V10.

% =============================================================================
\section{The split gradient bound}
\label{sec:c4v7-bound}
% =============================================================================

\begin{theorem}[Split gradient bound]
\label{thm:c4v7-gradient}
For every \(\theta>0\), \(\lambda>0\) with \((8/3)^{1/2}\kappa A(\theta,\lambda)<1\), every homogeneous
type-I point with type-I enstrophy constant \(C_I\) satisfies
\begin{equation}
 \boxed{\quad
 C_\infty'\ \le\ \frac{(2/3)^{1/2}c_{O,2}\,\kappa^2\log\frac{1+\theta}\theta+(8/3)^{1/2}S_3\,C_I^2\,\Phi(\theta,\lambda)}{1-(8/3)^{1/2}\kappa\,A(\theta,\lambda)}=:\Gamma_1(\kappa,C_I;\theta,\lambda) .
 \quad}
 \label{eq:c4v7-gradient}
\end{equation}
\end{theorem}

\begin{proof}
As in Theorem~\ref{thm:c4v3-gradient}: the ancient form split at
\(s_0=s-\theta L\); the far part \((-\infty,s_0)\) as there.  On \([s_0,s]\) the gradient
falls on the tensor, \(F=(\nabla(W\otimes W))_0\), and the kernel action at time
\(t=\nu(s-\tau)\) is split at \(|y|=\rho\), \(\rho^2=\nu\lambda L\), by Lemma~\ref{lem:c4v6-scaling}:
the near part is at most \(c_O^{s0}(\rho t^{-1/2})t^{-1/2}\|F\|_\infty\) with
\(\|F(\tau)\|_\infty\le(8/3)^{1/2}C_\infty(-\tau)^{-1/2}C_\infty'(-\tau)^{-1}\), the far part at most
\(m_3(\rho t^{-1/2})t^{-3/2}\|F(\tau)\|_{L^{3/2}}\) with Lemma~\ref{lem:c4v6-holder}.  Multiply by
\(L\), integrate over the window with Lemma~\ref{lem:c4v6-window} (the
powers of \(\nu\) cancel: \(C_\infty\nu^{-1/2}=\kappa\nu^{1/2}\nu^{-1/2}\) and \(\nu^{3/2}\nu^{-3/2}\)), take the
supremum over \(s\), and absorb (\(C_\infty'<\infty\)).
\end{proof}

% =============================================================================
\section{The floor as a function of the enstrophy constant}
\label{sec:c4v7-floor}
% =============================================================================

Let \(\Gamma_1(\kappa,C_I):=\inf_{\theta,\lambda}\Gamma_1(\kappa,C_I;\theta,\lambda)\) and
\(F(C_I):=\min\{\kappa:(2/3)^{1/2}\Gamma_1(\kappa,C_I)+\frac{\kappa^2}2\ge1\}\), which exists since the left side
is increasing in \(\kappa\) for fixed \((\theta,\lambda)\), continuous, \(0\) at \(\kappa=0\) when \(C_I=0\),
and at least \(\frac{\kappa^2}2\).

\begin{theorem}[Floor in the enstrophy constant]
\label{thm:c4v7-floor}
Every homogeneous type-I point of a smooth periodic solution with
type-I enstrophy constant \(C_I\) has \(\kappa\ge F(C_I)\), where \(F\) is nonincreasing,
\(F(C_I)=0.3438\) for \(C_I\ge1.93\), \(F(C_I)\to\sqrt2\) as \(C_I\to0\), and
\medskip
\begin{center}\small
\begin{tabular}{@{}lcccccccccc@{}}
\toprule
\(C_I\) & 0 & 0.25 & 0.5 & 0.75 & 1 & 1.25 & 1.5 & 1.75 & 1.9 & \(\ge1.93\)\\
\midrule
\(F(C_I)\ge\) & 1.412 & 1.354 & 1.198 & 0.986 & 0.776 & 0.602 & 0.472 & 0.380 & 0.345 & 0.3438\\
\bottomrule
\end{tabular}
\end{center}
\medskip
Equivalently, a homogeneous type-I point with velocity constant
\(\kappa<\sqrt2\) has enstrophy constant at least \(G(\kappa):=\sup\{C_I:F(C_I)>\kappa\}\):
\medskip
\begin{center}\small
\begin{tabular}{@{}lccccccccc@{}}
\toprule
\(\kappa\) & 0.345 & 0.4 & 0.5 & 0.6 & 0.7 & 0.8 & 1.0 & 1.2 & 1.4\\
\midrule
\(C_I\ge\) & 1.872 & 1.689 & 1.438 & 1.253 & 1.101 & 0.966 & 0.733 & 0.498 & 0.108\\
\bottomrule
\end{tabular}
\end{center}
\medskip
\end{theorem}

\begin{proof}
Theorem C3-V93-rows(1) with \eqref{eq:c4v7-gradient} and the infimum
over \((\theta,\lambda)\).  \(F\) is nonincreasing because \(\Gamma_1\) is nondecreasing in
\(C_I\).  The value \(0.3438\) for large \(C_I\): with \(\lambda\to\infty\) the split becomes the
sup-norm bound of V3 (\(A\to2c_O^{s0}(\theta/(1+\theta))^{1/2}\), \(\Phi\to0\)), so \(F\le0.3438\)
always, and equality holds once the far term at the sup-norm
optimum cannot be reduced without increasing the absorption by
more.  The limit \(\sqrt2\): at \(C_I=0\) the far term vanishes, \(\lambda\to0\) makes
\(A\to0\), and \(\theta\to\infty\) makes the far-kernel term vanish, leaving \(\frac{\kappa^2}2\ge1\).
The values: bisection in \(\kappa\) with the infimum over a grid of \(40\)
values of \(\theta\) in \([10^{-4},10^3]\) and \(50\) of \(\lambda\) in \([10^{-5},10^5]\), the window
functions evaluated as in Remark~\ref{rem:c4v6-evaluation} with the
certified constants; the inverse by bisection in \(C_I\).  At \(C_I=1\) the
optimum is a long window (\(\theta\approx10^3\)) with a split radius that is
small at the far end of the window (\(\lambda\approx0.12\)), i.e.\ the kernel action
is bounded through \(L^3\)--\(L^{3/2}\) over most of the window and through
the sup norm only near the current time.
\end{proof}

\begin{assessment}[Item (AC2): reading]
\label{ass:c4v7-reading}
The class has two constants at a homogeneous type-I point, \(\kappa\) and
\(C_I\), and the third cycle's floors constrained \(\kappa\) alone.  The split
bound couples them: a small enstrophy constant forces a large
velocity constant (\(C_I\le0.5\) forces \(\kappa\ge1.2\)), and conversely the velocity
constant of the third cycle's regime, \(\kappa\) near \(0.34\), is possible only
with \(C_I\ge1.87\).  The curve is a certified region in the \((\kappa,C_I)\)-plane
that every homogeneous type-I point must lie in; neither coordinate
is excluded, since the class fixes neither.  The mechanism is the
one of the second cycle's constant-chasing chapter (Theorem
C3-V1-constants) made explicit and quantitative: the enstrophy rate
bounds the nonlocal part of the kernel action, and only the local
part needs the sup norm.  What would sharpen the curve: the near
constant \(c_O^{s0}(R)\) is the full \(L^1\) mass inside \(R\) while the kernel's
action on a divergence-free tensor might be smaller (a cancellation
not used); the \(L^{3/2}\) bound uses the sharp Sobolev constant but the
pair \((\|W\|_6,\|\nabla W\|_2)\) at the same time could be related through the
Gaussian ledger; both are recorded, not pursued.
\end{assessment}

% =============================================================================
\section{Integrated V7 conclusion and next acquisition order}
\label{sec:c4v7-conclusion}
% =============================================================================

\begin{theorem}[What V7 proves]
\label{thm:c4v7-conclusion}
Relative to the three legacy cycles and V1--V6: the split gradient
bound (Theorem~\ref{thm:c4v7-gradient}) and the floor in the enstrophy
constant with its inverse (Theorem~\ref{thm:c4v7-floor}).  The full
Navier--Stokes stability/global-regularity theorem remains
unproved.
\end{theorem}

\begin{proof}
The cited results.
\end{proof}

\begin{programme}[V8 acquisition order]
\label{prog:c4v8}
\begin{enumerate}[label=\textup{(AC\arabic*)},leftmargin=4.8em]
\item The same split for the second-derivative bound (the window
      part of Theorem~\ref{thm:c4v4-gamma2} with the near constant and
      an \(L^{3/2}\) far term for \(\nabla^2(W\otimes W)\), which needs \(\|\nabla^2W\|_{L^2}\) or the
      product \(|W||\nabla^2W|\) in \(L^{3/2}\): record whether the class supplies it
      (it bounds \(\|\nabla^2W\|_2\) only through the derivative constants) or
      the negative.
\item The joint region drawn as a table of \((\kappa,C_I)\) corners, and the
      comparison with the second cycle's constant-chasing thresholds
      (\(\kappa<\frac12\), \(\kappa<\sqrt2\)): the regime \(\kappa<\frac12\) now requires \(C_I\ge1.44\).
\item The next compulsory companion audit is V10.
\end{enumerate}
\end{programme}

\begin{assessment}[Position after V7]
\label{ass:c4v7-position}
The floor is a curve in the \((\kappa,C_I)\)-plane.  No full stability or
global-regularity claim is made.
\end{assessment}

% =============================================================================
\section*{Version V7 (fourth cycle) append-only record}
\addcontentsline{toc}{section}{Version V7 (fourth cycle) append-only record}
% =============================================================================

The complete V6 source of the fourth cycle was retained before this
continuation.  Its SHA--256 digest was
\begin{center}\scriptsize\ttfamily
51bb2e53ac6e3605e7a4067f5a395aace52202e4000e0a6b138450a7ca162525.
\end{center}
V7 proved the split gradient bound and the floor in the enstrophy
constant.  No full regularity claim is made.

% =============================================================================
\section{V8: the second derivative does not split, and the joint region}
\label{sec:c4v8-entry}
% =============================================================================

Items (AC1) and (AC2) of Programme~\ref{prog:c4v8} are carried out.
(AC1): the near--far split cannot be applied to the second-derivative
bound, because the class does not bound the tensor \(\nabla^2(W\otimes W)\) in
\(L^{3/2}\): its two parts need \(\nabla W\in L^3\) or \(\nabla^2W\in L^2\), neither of which
the type-I class supplies without the derivative constants
themselves; recorded as a negative, with what would lift it.  (AC2):
the joint region in the \((\kappa,C_I)\)-plane is written as a table of
corners, and compared with the second cycle's constant-chasing
thresholds: the regime \(\kappa<\frac12\) of Theorem C2-V36-constraint now
requires \(C_I\ge1.44\), and the regime \(\kappa<\frac1{2\sqrt3}\) of Proposition
C3-V26-necessary is empty.  No global regularity claim is made.  The
next compulsory companion audit is V10.

% =============================================================================
\section{Item (AC1): the second derivative}
\label{sec:c4v8-second}
% =============================================================================

\begin{firewall}[No split for the second derivative]
\label{fw:c4v8-no-split}
The window part of the second-derivative bound
(Theorem~\ref{thm:c4v4-gamma2}) acts on \(F_2:=(\nabla^2(W\otimes W))_0\), bounded
pointwise by \((8/3)^{1/2}|W||\nabla^2W|+2|\nabla W|^2\).  A far part in
\(L^3\)--\(L^{3/2}\) needs \(\|F_2\|_{L^{3/2}}\): the term \(|\nabla W|^2\) is in \(L^1\) (\(\nabla W\in L^2\)) and in
\(L^{3/2}\) only if \(\nabla W\in L^3\); the term \(|W||\nabla^2W|\) is in \(L^{3/2}\) if \(\nabla^2W\in L^2\)
(with \(W\in L^6\)).  The class gives \(\nabla W\in L^2\) with the enstrophy rate and
\(\nabla W,\nabla^2W\in L^\infty\) with the derivative constants; \(\nabla W\in L^3\) follows by
interpolation, \(\|\nabla W\|_3^3\le\|\nabla W\|_\infty\|\nabla W\|_2^2\le C_\infty'(-\tau)^{-1}C_I^2\nu^{3/2}(-\tau)^{-1/2}\), but this
reintroduces \(C_\infty'\) in the far term and \(\nabla^2W\in L^2\) is not given at all.
The far term would then be \(2m_3\dots\|\nabla W\|_3^2\propto(C_\infty')^{2/3}C_I^{4/3}\), a bound
on \(C_\infty''\) in terms of \(C_\infty'\) and \(C_I\), which is available but does not
reduce the absorption of the second-derivative bound, whose window
part is governed by \(|W||\nabla^2W|\).  What would lift it: an \(L^2\) bound on
\(\nabla^2W\) from the class (a second-order enstrophy rate), which the type-I
hypothesis does not contain.  The second-derivative bound stays as
in V4.
\end{firewall}

% =============================================================================
\section{Item (AC2): the joint region}
\label{sec:c4v8-region}
% =============================================================================

\begin{theorem}[The joint region]
\label{thm:c4v8-region}
Every homogeneous type-I point of a smooth periodic solution, with
velocity constant \(\kappa\) and type-I enstrophy constant \(C_I\), lies in the
region
\begin{equation}
 \mathcal R:=\bigl\{(\kappa,C_I):\ \kappa\ge0.3438,\ \kappa\ge F(C_I)\bigr\},
 \label{eq:c4v8-region}
\end{equation}
whose boundary passes through the points
\medskip
\begin{center}\small
\begin{tabular}{@{}lcccccccc@{}}
\toprule
\(\kappa\) & 0.3438 & 0.4 & 0.5 & 0.7 & 1.0 & 1.2 & 1.4 & \(\sqrt2\)\\
\(C_I\) & 1.93 & 1.69 & 1.44 & 1.10 & 0.73 & 0.50 & 0.11 & 0\\
\bottomrule
\end{tabular}
\end{center}
\medskip
(\(\kappa\) increasing along the boundary as \(C_I\) decreases; for \(C_I\ge1.93\) the
boundary is the line \(\kappa=0.3438\)).  For \(\kappa\ge\sqrt2\) the region imposes no
condition on \(C_I\).
\end{theorem}

\begin{proof}
Theorems~\ref{thm:c4v3-floor} and \ref{thm:c4v7-floor}.
\end{proof}

\begin{proposition}[Comparison with the constant-chasing thresholds]
\label{prop:c4v8-comparison}
\begin{enumerate}[label=\textup{(\arabic*)},leftmargin=3.2em]
\item The regime \(\kappa<\frac12\), in which the mean constraint of Theorem
      C2-V36-constraint and the mean-square constraint of the third
      cycle's V32 apply, is compatible with the region only for
      \(C_I\ge1.44\); the regime \(\kappa<\frac1{2\sqrt3}=0.289\) of Proposition
      C3-V26-necessary is empty.
\item The regime \(\kappa<\sqrt2\) of Theorem C2-V41-spectrum is compatible with
      every \(C_I\ge0\); above \(\sqrt2\) the region is unrestricted.
\item At \(\kappa\in[0.3438,\frac12)\) the enstrophy constant satisfies \(C_I\ge1.44\), so
      the class's Gaussian enstrophy scale \(C=C_I^2\nu^{3/2}\) is at least
      \(2.07\nu^{3/2}\); the region bounds \(C\) below, not above, and the
      concentration bound \(\int|z|^2|\Omega|^2G\le12\nu C\) of Theorem~\ref{thm:c4v4-gamma2}
      is a bound by a multiple of a quantity that is itself at least
      \(2.07\nu^{3/2}\) there.
\end{enumerate}
\end{proposition}

\begin{proof}
Read off Theorem~\ref{thm:c4v8-region}.
\end{proof}

\begin{assessment}[Reading]
\label{ass:c4v8-reading}
The description of a hypothetical homogeneous type-I singular point
now has two certified coordinates: a rescaled velocity bound of at
least \(0.34\,\nu^{1/2}\), and, if that bound is below \(0.5\,\nu^{1/2}\), a type-I enstrophy
constant of at least \(1.44\); the second cycle's small-velocity regime
is thus a regime of order-one enstrophy constant.  Nothing is
excluded, the class fixing neither constant.  The block V6--V8 has
exhausted what the \(L^6\)--\(L^2\) information gives on the first derivative;
V9 consolidates and fixes the direction for V11--V15.
\end{assessment}

% =============================================================================
\section{Integrated V8 conclusion and next acquisition order}
\label{sec:c4v8-conclusion}
% =============================================================================

\begin{theorem}[What V8 establishes]
\label{thm:c4v8-conclusion}
Relative to the three legacy cycles and V1--V7: the negative record
(Firewall~\ref{fw:c4v8-no-split}), the joint region
(Theorem~\ref{thm:c4v8-region}) and the comparison
(Proposition~\ref{prop:c4v8-comparison}).  The full Navier--Stokes
stability/global-regularity theorem remains unproved.
\end{theorem}

\begin{proof}
The cited results.
\end{proof}

\begin{programme}[V9 acquisition order]
\label{prog:c4v9}
\begin{enumerate}[label=\textup{(AC\arabic*)},leftmargin=4.8em]
\item Consolidate the blocks V2--V4 and V6--V8 as one theorem with the
      constants sheet of the fourth cycle so far.
\item Fix the direction for V11--V15: the third cycle's ledger rows
      (the mean-square constraint, the extremal Dirichlet bounds, the
      defect bounds, the concentration bound) evaluated with the
      certified constants and the joint region, to complete the
      quantitative description of the hypothetical singular point;
      record which rows become numerical and which still contain
      non-explicit class constants (\(c_\Phi\), \(K_P\), \(\underline\delta\)).
\item The next compulsory companion audit is V10.
\end{enumerate}
\end{programme}

\begin{assessment}[Position after V8]
\label{ass:c4v8-position}
The joint region is written.  No full stability or global-regularity
claim is made.
\end{assessment}

% =============================================================================
\section*{Version V8 (fourth cycle) append-only record}
\addcontentsline{toc}{section}{Version V8 (fourth cycle) append-only record}
% =============================================================================

The complete V7 source of the fourth cycle was retained before this
continuation.  Its SHA--256 digest was
\begin{center}\scriptsize\ttfamily
c74c60bf7db1d60184f80bb58bddf086ec99933bdf218bbc762c69ceda39b3f7.
\end{center}
V8 recorded the second-derivative negative and the joint region.  No
full regularity claim is made.

% =============================================================================
\section{V9: the first two blocks consolidated, and the direction for V11--V15}
\label{sec:c4v9-entry}
% =============================================================================

Items (AC1) and (AC2) of Programme~\ref{prog:c4v9} are carried out:
the ancient-form block (V2--V4) and the far-field block (V6--V8) are
consolidated as one theorem with the constants sheet of the fourth
cycle so far, and the direction for V11--V15 is fixed: the third
cycle's ledger rows evaluated with the certified constants and the
joint region.  No global regularity claim is made.  The next
compulsory companion audit is V10.

% =============================================================================
\section{The two blocks}
\label{sec:c4v9-blocks}
% =============================================================================

\begin{theorem}[The ancient-form and far-field blocks, V2--V8]
\label{thm:c4v9-blocks}
At a homogeneous type-I point of a smooth periodic solution, with
velocity constant \(\kappa\) and type-I enstrophy constant \(C_I\):
\begin{enumerate}[label=\textup{(\arabic*)},leftmargin=3.2em]
\item (Kernels.)  \(\nabla^2K^O_1\) and \(\nabla^3K^O_1\) have the radial forms of
      Lemmas~\ref{lem:c4v2-fourth}, \ref{lem:c4v4-fifth}; on symmetric traceless
      inputs their maps decompose into three orthogonal channels; the
      tail constants are \(36\) and \(60\sqrt{15}\); the certified constants are
      \(c_{O,2}\le1.788\), \(c_{O,3}\le1.897\) (Theorems~\ref{thm:c4v2-certified},
      \ref{thm:c4v4-certified}); the near and far constants \(c_O^{s0}(R)\), \(m_3(R)\) are
      those of Proposition~\ref{prop:c4v6-values}.
\item (Gradient bound.)  \(C_\infty'\le\Gamma_1(\kappa,C_I;\theta,\lambda)\) for every admissible \((\theta,\lambda)\)
      (Theorem~\ref{thm:c4v7-gradient}), whose \(\lambda\to\infty\) form is the
      ancient-form bound of Theorem~\ref{thm:c4v3-gradient}.
\item (Second-derivative bound.)  \(\nu^{1/2}C_\infty''\le\) \eqref{eq:c4v4-gamma2}, \(6.62\) at the
      floor; no far-field split (Firewall~\ref{fw:c4v8-no-split}).
\item (Floors.)  \(\kappa\ge0.3438\) (Theorem~\ref{thm:c4v3-floor}); \(\kappa\ge F(C_I)\) with the
      curve of Theorem~\ref{thm:c4v7-floor}; the joint region \(\mathcal R\) of
      Theorem~\ref{thm:c4v8-region}; the concentration
      \(\int|z|^2|\Omega|^2G\le12\nu C\) at the floor.
\item (Negative records.)  No second split of the window
      (Firewall~\ref{fw:c4v4-no-second-split}); no far-field split of the
      second derivative (Firewall~\ref{fw:c4v8-no-split}).
\end{enumerate}
\end{theorem}

\begin{proof}
The cited results.
\end{proof}

\begin{theorem}[Constants sheet of the fourth cycle, V1--V9]
\label{thm:c4v9-sheet}
\medskip
\begin{center}\small
\begin{tabular}{@{}p{0.2\textwidth}p{0.4\textwidth}p{0.12\textwidth}p{0.18\textwidth}@{}}
\toprule
\textbf{Constant} & \textbf{Definition} & \textbf{Value} & \textbf{Status; where}\\
\midrule
\(c_{O,2}\) & \(\sup_tt\|\nabla e^{t\Delta}\PP\operatorname{div}\|_{s0}\) & \(\le1.788\) & certified; V2\\
\(c_{O,3}\) & \(\sup_tt^{3/2}\|\nabla^2e^{t\Delta}\PP\operatorname{div}\|_{s0}\) & \(\le1.897\) & certified; V4\\
tails & \(36/(4\pi r^5)\), \(60\sqrt{15}/(4\pi r^6)\) & exact & V2, V4\\
channel norms & \eqref{eq:c4v2-opnorm}; axial for \(M_3\) & exact & V2, V4\\
\(c_O^{s0}(R)\), \(m_3(R)\) & near \(L^1\), far \(L^3\) norms & table & certified; V6\\
\(S_3\) & sharp Sobolev constant, \(L^6(\R^3)\) & \(0.4273\) & closed form; V6\\
\(A(\theta,\lambda)\), \(\Phi(\theta,\lambda)\) & window functions \eqref{eq:c4v6-A}, \eqref{eq:c4v6-Phi} & formulas & V6\\
\(J_3(\theta)\) & \(2\theta^{-1/2}-\pi+2\arctan\theta^{1/2}\) & closed form & V4\\
\(\gamma_1^{\rm anc}\), \(\Gamma_1\) & gradient bounds \eqref{eq:c4v3-gamma1}, \eqref{eq:c4v7-gradient} & formulas & V3, V7\\
floor & \(\kappa\ge0.3438\) (\(\theta^*=0.0094\)) & certified & V3\\
curve & \(F(C_I)\), \(G(\kappa)\) & tables & certified; V7\\
\(C_\infty'\), \(\nu^{1/2}C_\infty''\) at the floor & & \(1.152\), \(6.62\) & V3, V4\\
\bottomrule
\end{tabular}
\end{center}
\medskip
\end{theorem}

\begin{proof}
The cited results.
\end{proof}

% =============================================================================
\section{Direction for V11--V15}
\label{sec:c4v9-direction}
% =============================================================================

\begin{assessment}[The ledger with values]
\label{ass:c4v9-direction}
The third cycle's rows were proved with symbolic class constants;
the fourth cycle has values for \(\kappa\) (a floor), \(C_\infty'\), \(C_\infty''\) (bounds at
the floor and as functions of \(\kappa\)), and a region for \((\kappa,C_I)\).  V11--V14
evaluate the rows: the mean-square constraint of C3-V32
(\((\frac14+v_{\rm exp})c_\Phi+\underline\delta\le\kappa^2(\frac12\langle d\rangle+8\nu C)\), with \(\langle d\rangle\le4\nu C\) and \(C=C_I^2\nu^{3/2}\)); the
defect bounds of C3-V33 (\(\underline\delta\le10\kappa^2\nu C\)) and C3-V48, C3-V59; the extremal
Dirichlet bounds of C3-V34; the weighted-enstrophy concentration
and the (void) palinstrophy row; the pressure bounds of C3-V47
(\(C_P\) in terms of \(C_\infty,C_\infty',C_\infty'',C\)).  For each: which quantities become
numbers in \(\nu\) and \(C_I\), and which still carry \(c_\Phi\), \(K_P\), \(\underline\delta\) or \(K_6\)
(the last now \(S_3C_I\nu^{3/4}\)).  The outcome is the quantitative description
of the hypothetical singular point at the level of the class,
recorded in one table (V14); V15 audits.  The directions (D3)--(D5)
of the third cycle's V99 remain for V16 onward.  No global
regularity claim is made.
\end{assessment}

% =============================================================================
\section{Integrated V9 conclusion and next acquisition order}
\label{sec:c4v9-conclusion}
% =============================================================================

\begin{theorem}[What V9 establishes]
\label{thm:c4v9-conclusion}
Relative to the three legacy cycles and V1--V8: the consolidated
blocks (Theorem~\ref{thm:c4v9-blocks}), the constants sheet
(Theorem~\ref{thm:c4v9-sheet}) and the direction.  The full
Navier--Stokes stability/global-regularity theorem remains
unproved.
\end{theorem}

\begin{proof}
The cited results.
\end{proof}

\begin{programme}[V10 acquisition order]
\label{prog:c4v10}
\begin{enumerate}[label=\textup{(AC\arabic*)},leftmargin=4.8em]
\item The compulsory companion audit: read the current SM and YM
      notes and record their digests and decision.
\item Record the position after ten versions; delivery.
\end{enumerate}
\end{programme}

\begin{assessment}[Position after V9]
\label{ass:c4v9-position}
Two blocks are consolidated; the direction is the ledger with
values.  No full stability or global-regularity claim is made.
\end{assessment}

% =============================================================================
\section*{Version V9 (fourth cycle) append-only record}
\addcontentsline{toc}{section}{Version V9 (fourth cycle) append-only record}
% =============================================================================

The complete V8 source of the fourth cycle was retained before this
continuation.  Its SHA--256 digest was
\begin{center}\scriptsize\ttfamily
785ffb02d4003ce3fbdd6f680f9d40b4efff3388e9d81e42e0f191adf940726a.
\end{center}
V9 consolidated the first two blocks and fixed the direction.  No
full regularity claim is made.

% =============================================================================
\section{V10: compulsory companion audit and the position after ten versions}
\label{sec:c4v10-entry}
% =============================================================================

This is the tenth version of the fourth cycle.  The compulsory
review of the two companion programmes is performed first, and the
position after ten versions is recorded.  No global regularity claim
is made.  The next compulsory companion audit is V15.

% =============================================================================
\section{Compulsory V10 review of the current companion notes}
\label{sec:c4v10-companion-audit}
% =============================================================================

The current companion files in \texttt{latex\_notes} were inspected
read-only.  Since the V5 audit the SM programme has moved from V53 to
V57 of its seventh series (at a first inspection for this audit the
file was V56, superseded before the final inspection, whose digest is
recorded), the YM programme from V29 to V32 of its sixth series.  The
frozen checkpoints are
\begin{align}
 &\texttt{Renewal\_SM\_Einstein\_Continuum\_Research\_Notes\_V57.tex},
 \notag\\[-2pt]
 &\qquad
 \texttt{ed4a995eb7788dd46323ccbb32ef994fb6348664514cbd64f1a421435c3b64af},
 \label{eq:c4v10-SM-checkpoint}\\
 &\texttt{Renewal\_Yang\_Mills\_Mass\_Gap\_Research\_Notes\_V32.tex},
 \notag\\[-2pt]
 &\qquad
 \texttt{1d61f4bc86d7e12b0145402727ed8689e863b62173b88977d6f961e4e656b1ba}.
 \label{eq:c4v10-YM-checkpoint}
\end{align}
(SM V57 has 22864 lines; YM V32 has 14579 lines.  The folder also
held YM V30 and V31, distinct earlier states of 13532 and 14096 lines,
superseded.)

\emph{SM V54--V57.}  The SM programme treated pure gravitational heat
owners (the complete local coefficient, spin curvature and the Dirac
coefficient, an Euler--Weyl decomposition, a Standard Model
multiplicity ledger, transfer through its stratified construction,
the remaining microscopic gravitational subtraction), then, after
its own audit, exact microscopic metric contacts (a common-base
relative heat trace, both metric contacts before localization, an
exact frame and diffeomorphism Ward cancellation, microscopic
heat-radius locality, the remaining microscopic card), and a
microscopic mesh-Ward reduction (microscopic metric polarization,
the exact Ward defect, a conditional closure of the microscopic
row), and a local scheme equivalence without its condition MIDC (a
logarithmic matching scale, a local expansion of the microscopic
core, declared mesh counterterms, closure of the microscopic
parity-even smooth card).

\emph{YM V30--V32.}  The YM programme performed its own audit and
treated mesoscopic axial charts and a good-sector covariance card,
a cofinal pure-Wilson static card, and unsplit stability under its
generated action.

\begin{theorem}[V10 companion-audit decision]
\label{thm:c4v10-companion-decision}
Nothing is imported from SM V57 or YM V32.  Neither bears on the
fourth cycle's blocks.
\end{theorem}

\begin{proof}
The SM results concern heat-trace coefficients and Ward identities
of lattice Dirac interfaces; the YM results concern static cards
and stability of a lattice gauge theory.  None is a Navier--Stokes
statement, and none concerns the mild formulation, the Oseen
kernels, the near--far split or the joint region.
\end{proof}

% =============================================================================
\section{Position after ten versions}
\label{sec:c4v10-position}
% =============================================================================

\begin{assessment}[Position after ten versions of the fourth cycle]
\label{ass:c4v10-position-ten}
Two blocks are written (Theorem~\ref{thm:c4v9-blocks}).  The certified
description of a homogeneous type-I point of a smooth periodic
solution is: velocity constant \(\kappa\ge0.3438\); if \(\kappa<\frac12\), type-I enstrophy
constant \(C_I\ge1.44\); the joint region \(\mathcal R\) of Theorem~\ref{thm:c4v8-region};
\(C_\infty'\le1.152\) and \(\nu^{1/2}C_\infty''\le6.62\) at the floor; the Gaussian second
vorticity moment at most \(12\nu C\).  The block V11--V15 follows
Assessment~\ref{ass:c4v9-direction}: the third cycle's ledger rows with
values.  No global regularity claim is made.
\end{assessment}

% =============================================================================
\section{Integrated V10 conclusion and next acquisition order}
\label{sec:c4v10-conclusion}
% =============================================================================

\begin{theorem}[What V10 establishes]
\label{thm:c4v10-conclusion}
Relative to the three legacy cycles and V1--V9: the companion-audit
decision (Theorem~\ref{thm:c4v10-companion-decision}) and the position.
The full Navier--Stokes stability/global-regularity theorem remains
unproved.
\end{theorem}

\begin{proof}
The cited records.
\end{proof}

\begin{programme}[V11 acquisition order]
\label{prog:c4v11}
\begin{enumerate}[label=\textup{(AC\arabic*)},leftmargin=4.8em]
\item The mean-square constraint of C3-V32 with values: with
      \(\langle d\rangle\le4\nu C\), \(C=C_I^2\nu^{3/2}\), \(\kappa\ge F(C_I)\) and the region, write
      \((\frac14+v_{\rm exp})c_\Phi+\underline\delta\le\kappa^2(2\nu C+8\nu C)=10\kappa^2C_I^2\nu^{5/2}\), and the defect
      bound \(\underline\delta\le10\kappa^2\nu C=10\kappa^2C_I^2\nu^{5/2}\) of C3-V33; record the pinching
      consequence \(c_\Phi\le40\kappa^2C_I^2\nu^{5/2}\) and what it says on the region's
      boundary.
\item The extremal Dirichlet bounds of C3-V34 with values.
\item The next compulsory companion audit is V15.
\end{enumerate}
\end{programme}

\begin{assessment}[Position after V10]
\label{ass:c4v10-position}
The companion audit is recorded; the ledger-with-values block
begins.  No full stability or global-regularity claim is made.
\end{assessment}

% =============================================================================
\section*{Version V10 (fourth cycle) append-only record}
\addcontentsline{toc}{section}{Version V10 (fourth cycle) append-only record}
% =============================================================================

The complete V9 source of the fourth cycle was retained before this
continuation.  Its SHA--256 digest was
\begin{center}\scriptsize\ttfamily
cf3f151c8c75478d1a96030cc54bf7b950ee7cc9251ff84af2f69612bde1d3c0.
\end{center}
V10 performed the compulsory companion audit and recorded the
position.  No full regularity claim is made.

% =============================================================================
\section{V11: the mean-square constraint and the extremal Dirichlet bounds with values}
\label{sec:c4v11-entry}
% =============================================================================

Items (AC1) and (AC2) of Programme~\ref{prog:c4v11} are carried out.
With \(C=C_I^2\nu^{3/2}\) and \(\langle d\rangle\le4\nu C\), the mean-square constraint of the
third cycle's V32 reads \((\frac14+v_{\rm exp})c_\Phi+\underline\delta\le10\kappa^2C_I^2\nu^{5/2}\), so the
pinching constant and the defect rate are bounded by explicit
multiples of \(\kappa^2C_I^2\nu^{5/2}\), which on the boundary of the joint region
is at most \(0.53\); the pinching bound improves on the trivial one
\(c_\Phi\le|G|\kappa^2\nu\) exactly when \(C_I<1.05\), i.e.\ on the part of the boundary
with \(\kappa>0.72\), and tends to zero with \(C_I\).  The extremal Dirichlet
bounds of V34 become explicit in \((\kappa,C_I)\) on the upper side; the lower
side keeps \(c_\Phi\).  No global regularity claim is made.  The next
compulsory companion audit is V15.

% =============================================================================
\section{Notation and the two trivial bounds}
\label{sec:c4v11-notation}
% =============================================================================

Profile notation of the second and third cycles: \(\varphi=\|V\|_G^2\), \(d=4\nu\|\nabla V\|_G^2\),
\(\mathcal L=(d+2\varphi)/8\), \(c_\Phi\le\varphi\le C_\Phi\) the pinching constants, \(\underline\delta\) the defect rate,
\(v_{\rm exp}\) the explicit variance floor (C3-V23), \(|G|=\int e^{-|z|^2/4\nu}=(4\pi\nu)^{3/2}=44.55\,\nu^{3/2}\),
\(C=C_I^2\nu^{3/2}\) the class enstrophy bound (\(\|\nabla V\|_{L^2}^2\le C\)).

\begin{lemma}[Trivial bounds]
\label{lem:c4v11-trivial}
For every profile, \(\varphi\le|G|C_\infty^2=44.55\,\kappa^2\nu^{5/2}\) and \(d\le4\nu C=4C_I^2\nu^{5/2}\); hence
\(C_\Phi\le44.55\,\kappa^2\nu^{5/2}\) and \(\mathcal L\le\frac12C_I^2\nu^{5/2}+11.14\,\kappa^2\nu^{5/2}\).
\end{lemma}

\begin{proof}
\(|V|\le C_\infty=\kappa\nu^{1/2}\) and \(G\le1\).
\end{proof}

% =============================================================================
\section{The mean-square constraint with values}
\label{sec:c4v11-constraint}
% =============================================================================

\begin{theorem}[Pinching and defect rate in the joint region]
\label{thm:c4v11-constraint}
At every homogeneous type-I point of a smooth periodic solution
with constants \((\kappa,C_I)\in\mathcal R\),
\begin{equation}
 \boxed{\quad
 \Bigl(\frac14+v_{\rm exp}\Bigr)c_\Phi+\underline\delta\ \le\ 10\,\kappa^2C_I^2\,\nu^{5/2},\qquad
 c_\Phi\le40\,\kappa^2C_I^2\,\nu^{5/2},\qquad \underline\delta\le10\,\kappa^2C_I^2\,\nu^{5/2} .
 \quad}
 \label{eq:c4v11-constraint}
\end{equation}
On the boundary of \(\mathcal R\) (Theorem~\ref{thm:c4v8-region}) the product \(\kappa^2C_I^2\)
takes the values
\medskip
\begin{center}\small
\begin{tabular}{@{}lcccccccc@{}}
\toprule
\(\kappa\) & 0.3438 & 0.4 & 0.5 & 0.7 & 1.0 & 1.2 & 1.4 & \(\to\sqrt2\)\\
\(C_I\) & 1.93 & 1.69 & 1.44 & 1.10 & 0.73 & 0.50 & 0.11 & \(\to0\)\\
\(\kappa^2C_I^2\) & 0.44 & 0.46 & 0.52 & 0.59 & 0.53 & 0.36 & 0.023 & \(\to0\)\\
\bottomrule
\end{tabular}
\end{center}
\medskip
so that on the boundary \(c_\Phi\le23.8\,\nu^{5/2}\) and \(\underline\delta\le5.94\,\nu^{5/2}\) throughout, and both
tend to zero along the boundary as \(C_I\to0\).  Inside \(\mathcal R\) (larger \(C_I\) at
fixed \(\kappa\)) the bounds grow with \(C_I^2\).
\end{theorem}

\begin{proof}
Theorem C3-V32-constraint with \(\langle d\rangle\le4\nu C\) and \(C=C_I^2\nu^{3/2}\):
\(\kappa^2(\frac12\langle d\rangle+8\nu C)\le\kappa^2\cdot10\nu C=10\kappa^2C_I^2\nu^{5/2}\); then \(\frac14c_\Phi\le\) the left side
and \(\underline\delta\le\) the left side.  The table from Theorem~\ref{thm:c4v8-region}.
\end{proof}

\begin{proposition}[When the pinching bound is not trivial]
\label{prop:c4v11-pinching}
\(40\kappa^2C_I^2\nu^{5/2}<44.55\kappa^2\nu^{5/2}\) iff \(C_I<1.055\); on the boundary of \(\mathcal R\) this is
the arc \(\kappa>0.72\) (where \(C_I<1.055\)).  There the constraint says that
the Gaussian energy of every profile at the homogeneous point is
at most \(40\kappa^2C_I^2\nu^{5/2}\), a fraction \(0.90\,C_I^2\) of the sup-norm bound; as
\(C_I\to0\) along the boundary the Gaussian energy must vanish, in
agreement with the ancient-form floor \(\kappa\to\sqrt2\) of
Theorem~\ref{thm:c4v7-floor}: a profile of small enstrophy constant has
small Gaussian energy however large its sup norm.
\end{proposition}

\begin{proof}
Compare \eqref{eq:c4v11-constraint} with Lemma~\ref{lem:c4v11-trivial}.
\end{proof}

% =============================================================================
\section{The extremal Dirichlet bounds with values}
\label{sec:c4v11-dirichlet}
% =============================================================================

\begin{theorem}[Extremal Dirichlet bounds in the joint region]
\label{thm:c4v11-dirichlet}
At every homogeneous type-I point with \((\kappa,C_I)\in\mathcal R\), in units of \(\nu^{5/2}\),
\begin{equation}
 \mathcal L_{\max}\ \le\ \mathcal L_+\ \le\ 22.28\,\kappa^4+\frac12\bigl(1985\,\kappa^8+356.4\,\kappa^4C_I^2\bigr)^{1/2},\qquad
 d\ \le\ 8\mathcal L_+\ \text{on the whole family},
 \label{eq:c4v11-Lplus}
\end{equation}
while \(\mathcal L_{\min}\ge\mathcal L_-=\frac{c_\Phi}4(1+\frac1{4\kappa^2})-2C_I^2\nu^{5/2}\) keeps the pinching constant,
which the class does not bound below explicitly.  At the corner
\((0.3438,1.93)\): \(\mathcal L_+\le0.31+\frac12(0.39+18.6)^{1/2}=2.49\) (\(\nu^{5/2}\)) and \(d\le19.9\,\nu^{5/2}\),
against the trivial \(d\le4C_I^2=14.9\,\nu^{5/2}\); at \((1.0,0.73)\): \(\mathcal L_+\le22.3+\frac12(1985+190)^{1/2}=45.6\)
against the trivial \(\mathcal L\le0.27+11.1=11.4\).
\end{theorem}

\begin{proof}
Theorem C3-V34-dirichlet with \(C_\Phi\le44.55\kappa^2\) and \(C=C_I^2\) (units \(\nu^{5/2}\),
\(\nu^{3/2}\)): \(\frac{\kappa^2C_\Phi}2\le22.28\kappa^4\), \(\kappa^4C_\Phi^2\le1985\kappa^8\), \(8\kappa^2\nu CC_\Phi\le356.4\kappa^4C_I^2\).  The
numbers by substitution.
\end{proof}

\begin{assessment}[Reading]
\label{ass:c4v11-reading}
The mean-square row becomes a numerical bound on the pinching
constant and the defect rate along the region's boundary, and is
informative on the arc of large velocity constant, where it says
that the Gaussian energy is small; the extremal Dirichlet row's
upper side becomes numerical but is weaker than the trivial bound
on most of the region (the extremal argument pays \(\kappa^2C_\Phi\) twice),
and its lower side needs \(c_\Phi\).  The pattern of the block is visible
already: the rows that bound quantities from above become numbers;
the rows that need a floor on the Gaussian energy or on the defect
rate do not, the class giving neither (N-C3-3 and its analogue for
\(c_\Phi\)).
\end{assessment}

% =============================================================================
\section{Integrated V11 conclusion and next acquisition order}
\label{sec:c4v11-conclusion}
% =============================================================================

\begin{theorem}[What V11 proves]
\label{thm:c4v11-conclusion}
Relative to the three legacy cycles and V1--V10: the trivial bounds
(Lemma~\ref{lem:c4v11-trivial}), the mean-square constraint with values
(Theorem~\ref{thm:c4v11-constraint}, Proposition~\ref{prop:c4v11-pinching})
and the extremal Dirichlet bounds with values
(Theorem~\ref{thm:c4v11-dirichlet}).  The full Navier--Stokes
stability/global-regularity theorem remains unproved.
\end{theorem}

\begin{proof}
The cited results.
\end{proof}

\begin{programme}[V12 acquisition order]
\label{prog:c4v12}
\begin{enumerate}[label=\textup{(AC\arabic*)},leftmargin=4.8em]
\item The pressure bounds of C3-V47 with values: the kernel constants
      \(c_1\), \(c_2\), \(c_3\) of the Riesz kernel \(K_{ij}=(3w_iw_j/|w|^2-\delta_{ij})/(4\pi|w|^3)\)
      computed exactly (\(|K|_F=\sqrt6/(4\pi|w|^3)\)), then
      \(C_{P,0}\), \(C_P\), \(C_t\), \(C_Y'\) in terms of \(\kappa,C_\infty',C_\infty'',C_I\), and their values at the
      floor.
\item The bounds \(C_Y\), \(C_N\) on \(\|Y\|_G\), \(\|N\|_G\) with values, and the
      variance cap \(\mathrm{Var}_{\max}\) of C3-V66.
\item The next compulsory companion audit is V15.
\end{enumerate}
\end{programme}

\begin{assessment}[Position after V11]
\label{ass:c4v11-position}
The mean-square and Dirichlet rows have values.  No full stability
or global-regularity claim is made.
\end{assessment}

% =============================================================================
\section*{Version V11 (fourth cycle) append-only record}
\addcontentsline{toc}{section}{Version V11 (fourth cycle) append-only record}
% =============================================================================

The complete V10 source of the fourth cycle (after the in-place
update of its companion-audit paragraph to SM V57, made before
delivery) was retained before this continuation.  Its SHA--256
digest was
\begin{center}\scriptsize\ttfamily
02ac3822bf6a1d7c28e56b317f5f546eada86f36625d2e57302aa397bd0cff9e.
\end{center}
V11 gave the mean-square constraint and the extremal Dirichlet
bounds their values.  No full regularity claim is made.

% =============================================================================
\section{V12: the pressure bounds and the defect norms with values}
\label{sec:c4v12-entry}
% =============================================================================

Items (AC1) and (AC2) of Programme~\ref{prog:c4v12} are carried out.
The kernel constants of the third cycle's V47 are computed exactly
from the Frobenius norm \(\sqrt6/(4\pi|w|^3)\) of the Riesz kernel, with the
split radius \(\nu^{1/2}\) that makes them scale correctly; the pointwise
pressure bounds \(C_{P,0}\), \(C_P\), the time-derivative bound \(C_t\) and the
defect bound \(C_Y'\) become explicit in \((\kappa,C_\infty',\nu^{1/2}C_\infty'',C_I)\), and the class
bounds \(C_Y\), \(C_N\) on \(\|Y\|_G\), \(\|N\|_G\) follow; values at the corner \((0.3438,1.93)\)
of the joint region are recorded.  The variance cap \(\mathrm{Var}_{\max}\) keeps
the pinching constant and is not numerical.  No global regularity
claim is made.  The next compulsory companion audit is V15.

% =============================================================================
\section{The Riesz kernel constants}
\label{sec:c4v12-kernel}
% =============================================================================

\begin{lemma}[Exact kernel constants]
\label{lem:c4v12-kernel}
The kernel of \(R_iR_j\) off the diagonal is \(K_{ij}(w)=(3w_iw_j/|w|^2-\delta_{ij})/(4\pi|w|^3)\),
with Frobenius norm \(|K(w)|_F=\sqrt6/(4\pi|w|^3)\).  With the split at \(|w|=\nu^{1/2}\)
(the third cycle's V47 wrote the split at \(|w|=1\) in the profile
coordinates, i.e.\ with \(\nu\) normalized to one), the constants of
Lemma C3-V47-riesz are
\begin{equation}
 c_1=\int_{|w|<\nu^{1/2}}|K|_F|w|\,\dd w=\sqrt6\,\nu^{1/2},\quad
 c_2=\Bigl(\int_{|w|>\nu^{1/2}}|K|_F^2\Bigr)^{1/2}=(2\pi)^{-1/2}\nu^{-3/4},\quad
 c_3=\Bigl(\int_{|w|>\nu^{1/2}}|K|_F^{3/2}\Bigr)^{2/3}=\frac{\sqrt6}{4\pi}\Bigl(\frac{8\pi}3\Bigr)^{2/3}\nu^{-1/2}=0.8040\,\nu^{-1/2},
 \label{eq:c4v12-kernel}
\end{equation}
and \(c_6=\frac{\sqrt6}{4\pi}(\frac{20\pi}3)^{5/6}\nu^{-1/4}=2.459\,\nu^{-1/4}\).
\end{lemma}

\begin{proof}
\(|3nn-I|_F^2=9-6+3=6\).  Radial integrals: \(\int_{|w|<\rho}|w|^{-2}\dd w=4\pi\rho\);
\(\int_{|w|>\rho}|w|^{-6}=4\pi/(3\rho^3)\); \(\int_{|w|>\rho}|w|^{-9/2}=4\pi\cdot\frac23\rho^{-3/2}\); \(\int_{|w|>\rho}|w|^{-18/5}=4\pi\cdot\frac53\rho^{-3/5}\);
with \(\rho=\nu^{1/2}\).
\end{proof}

% =============================================================================
\section{The pressure bounds with values}
\label{sec:c4v12-pressure}
% =============================================================================

Write \(c'':=\nu^{1/2}C_\infty''\) (dimensionless).

\begin{theorem}[Pressure, time derivative and defect bounds in the joint region]
\label{thm:c4v12-pressure}
At every homogeneous type-I point with constants \((\kappa,C_I)\), \(C_\infty'\), \(c''\):
\begin{align}
 |P_V|&\le C_{P,0}=\Bigl(\frac{\kappa^2}3+2\sqrt6\,\kappa C_\infty'+0.1468\,C_I^2\Bigr)\nu,
 \label{eq:c4v12-P0}\\
 |\nabla P_V|&\le C_P=\sqrt3\Bigl(\frac23\kappa C_\infty'+2\sqrt6\,(C_\infty'^2+\kappa c'')+0.7979\,\kappa C_I\Bigr)\nu^{1/2},
 \label{eq:c4v12-P}\\
 |V_t|&\le C_t=(c''+\kappa C_\infty')\nu^{1/2}+C_P,\qquad |Y(z)|\le C_Y'+\tfrac12C_\infty'|z|,\quad C_Y'=\tfrac12\kappa\nu^{1/2}+C_t .
 \label{eq:c4v12-t}
\end{align}
At the corner \((0.3438,1.93)\) with \(C_\infty'\le1.152\), \(c''\le6.62\):
\(C_{P,0}\le2.53\,\nu\), \(C_P\le31.95\,\nu^{1/2}\), \(C_t\le38.96\,\nu^{1/2}\), \(C_Y'\le39.14\,\nu^{1/2}\).
\end{theorem}

\begin{proof}
Lemma C3-V47-decomposition with Lemma~\ref{lem:c4v12-kernel} and
\(K_6\le S_3C_I\nu^{3/4}\) (Lemma~\ref{lem:c4v6-holder}): \(c_3K_6^2=0.8040\,S_3^2C_I^2\nu=0.1468\,C_I^2\nu\);
\(2c_2C_\infty C^{1/2}=2(2\pi)^{-1/2}\nu^{-3/4}\kappa\nu^{1/2}C_I\nu^{3/4}=0.7979\,\kappa C_I\nu^{1/2}\); \(c_1(2C_\infty'^2+2C_\infty C_\infty'')=2\sqrt6(C_\infty'^2+\kappa c'')\nu^{1/2}\).
The values by substitution.
\end{proof}

\begin{theorem}[The class bounds on the nonlinearity and the defect]
\label{thm:c4v12-YN}
With \(|G|^{1/2}=(4\pi\nu)^{3/4}=6.674\,\nu^{3/4}\) and \(\||z|\|_G=(6\nu|G|)^{1/2}\):
\begin{align}
 \|N\|_G&\le(\kappa C_\infty'\nu^{1/2}+C_P)|G|^{1/2}=:C_N^{1/2},\\
 \|Y\|_G&\le\|L_-V\|_G+\|N\|_G\le\Bigl(c''+\frac\kappa2+\frac{\sqrt6}2C_\infty'\Bigr)\nu^{1/2}|G|^{1/2}+C_N^{1/2}=:C_Y ,
 \label{eq:c4v12-YN}
\end{align}
and at the corner: \(C_N^{1/2}\le215.9\,\nu^{5/4}\), \(C_Y\le270.6\,\nu^{5/4}\).  The variance cap
\(\mathrm{Var}_{\max}=(C_Y+C_N^{1/2})^2/c_\Phi\) of Theorem C3-V66-floor keeps \(c_\Phi\).
\end{theorem}

\begin{proof}
\(|N|\le|(V\cdot\nabla)V|+|\nabla P_V|\le C_\infty C_\infty'+C_P\) pointwise, and \(\|1\|_G=|G|^{1/2}\).
\(|L_-V|\le\nu|\Delta V|+\frac12|z||\nabla V|+\frac12|V|\le\nu C_\infty''+\frac12C_\infty'|z|+\frac12C_\infty\), and
\(\frac12C_\infty'\||z|\|_G=\frac12C_\infty'(6\nu)^{1/2}|G|^{1/2}=\frac{\sqrt6}2C_\infty'\nu^{1/2}|G|^{1/2}\).  Values:
\((0.396+31.95)\cdot6.674=215.9\); \((6.62+0.172+1.411)\cdot6.674+215.9=270.6\).
\end{proof}

\begin{assessment}[Reading]
\label{ass:c4v12-reading}
The pressure gradient at a homogeneous type-I point with the
floor's constants is at most \(32\,\nu^{1/2}\) in profile units, dominated by
the term \(2\sqrt6\,\kappa c''\) (the second derivative through the near part of
the Riesz kernel); the pointwise defect bound \(C_Y'\) is of the same
size.  These are the numbers that the third cycle's pressure firewall
(C3-V47) left symbolic; they are large because the second-derivative
constant is, and they enter the third cycle's rows only as upper
bounds.  The variance cap and every quantity divided by \(c_\Phi\) stay
symbolic.
\end{assessment}

% =============================================================================
\section{Integrated V12 conclusion and next acquisition order}
\label{sec:c4v12-conclusion}
% =============================================================================

\begin{theorem}[What V12 proves]
\label{thm:c4v12-conclusion}
Relative to the three legacy cycles and V1--V11: the exact kernel
constants (Lemma~\ref{lem:c4v12-kernel}), the pressure bounds with
values (Theorem~\ref{thm:c4v12-pressure}) and the class bounds on the
nonlinearity and the defect (Theorem~\ref{thm:c4v12-YN}).  The full
Navier--Stokes stability/global-regularity theorem remains
unproved.
\end{theorem}

\begin{proof}
The cited results.
\end{proof}

\begin{programme}[V13 acquisition order]
\label{prog:c4v13}
\begin{enumerate}[label=\textup{(AC\arabic*)},leftmargin=4.8em]
\item The vorticity rows with values at the floor: the concentration
      \(\int|z|^2|\Omega|^2G\le12\nu C=12C_I^2\nu^{5/2}\) (\(\theta_1^{ST}=\frac12\)); the enstrophy
      \(\mathcal E\le2C=2C_I^2\nu^{3/2}\); the mean-vorticity and dipole rows
      (\(|G||\bar\Omega|^2\le\kappa^2(3\mathcal E+4\nu\mathcal P)\), \(\mathcal P\le2C_\infty''^2|G|\)) with values; the
      (E5-small) reading: at the floor \(C_\infty'/\kappa=3.35\), outside the
      window \((0.807,1.094)\) of C3-V73, consistently with the voidness
      of the palinstrophy row.
\item The defect-rate bounds of C3-V48 and C3-V59 with values, as
      far as they are explicit.
\item The next compulsory companion audit is V15.
\end{enumerate}
\end{programme}

\begin{assessment}[Position after V12]
\label{ass:c4v12-position}
The pressure and defect norms have values.  No full stability or
global-regularity claim is made.
\end{assessment}

% =============================================================================
\section*{Version V12 (fourth cycle) append-only record}
\addcontentsline{toc}{section}{Version V12 (fourth cycle) append-only record}
% =============================================================================

The complete V11 source of the fourth cycle was retained before this
continuation.  Its SHA--256 digest was
\begin{center}\scriptsize\ttfamily
0e8a4981de00369ed57cbb3e20320d7c9e8c11d39dcb95d2a2742080459d1862.
\end{center}
V12 gave the pressure bounds and the defect norms their values.  No
full regularity claim is made.

% =============================================================================
\section{V13: the vorticity rows with values, and the void defect row}
\label{sec:c4v13-entry}
% =============================================================================

Items (AC1) and (AC2) of Programme~\ref{prog:c4v13} are carried out.
The vorticity rows of the third cycle are evaluated at the floor
corner \((\kappa,C_I)=(0.3438,1.93)\) of the joint region: the concentration
bound, the enstrophy and palinstrophy bounds, and the
mean-vorticity and dipole rows, the last two compared with their
trivial forms (the mean row through the concentration bound beats
the trivial bound by a factor \(5.6\), the dipole row by \(1.7\)); the
(E5-small) reading is recorded.  A structural observation: the
extremal defect row of the third cycle's V38 and V48 is void at every
homogeneous type-I point, because its threshold is \(\frac12\) minus the
left side of the enstrophy constraint, which is at least \(1\).  The
fourth bound of V59 is local and stays symbolic.  No global regularity
claim is made.  The next compulsory companion audit is V15.

% =============================================================================
\section{The vorticity rows with values}
\label{sec:c4v13-rows}
% =============================================================================

Notation of the third cycle's V67--V78: \(\mathcal E=\int|\Omega|^2G\), \(\mathcal E_1=\int|z|^2|\Omega|^2G\),
\(\mathcal P=\int|\nabla\Omega|^2G\), \(\bar\Omega=|G|^{-1}\int\Omega G\), \(D=\int z\otimes\Omega\,G\); \(c''=\nu^{1/2}C_\infty''\).

\begin{theorem}[Vorticity rows at the floor corner]
\label{thm:c4v13-rows}
At every homogeneous type-I point with \((\kappa,C_I)\) on the boundary of
\(\mathcal R\) at the corner \((0.3438,1.93)\), \(C_\infty'\le1.152\), \(c''\le6.62\):
\begin{align}
 \mathcal E&\le2C=2C_I^2\nu^{3/2}=7.45\,\nu^{3/2},\qquad
 \sup_{\mathcal F(x_0)}\mathcal E_1\le12\nu C=12C_I^2\nu^{5/2}=44.7\,\nu^{5/2},\qquad
 \mathcal P\le2C_\infty''^2|G|=89.1\,c''^2\nu^{1/2}\le3904\,\nu^{1/2},
 \label{eq:c4v13-EP}\\
 \sup_{\mathcal F(x_0)}|G||\bar\Omega|^2&\le\frac{\kappa^2}{4\nu}\sup\mathcal E_1\le3\kappa^2C=3\kappa^2C_I^2\nu^{3/2}=1.32\,\nu^{3/2}
 \quad(\text{trivial: }|G||\bar\Omega|^2\le\mathcal E\le7.45\,\nu^{3/2}),
 \label{eq:c4v13-mean}\\
 \sup_{\mathcal F(x_0)}|D|^2&\le\kappa^2|G|\Bigl(\frac{16}9\nu\,\mathcal E+\frac43\sup\mathcal E_1\Bigr)\le\kappa^2|G|\Bigl(\frac{32}9+16\Bigr)C_I^2\nu^{5/2}=383\,\nu^4
 \quad(\text{trivial: }|D|^2\le2\nu|G|\mathcal E\le664\,\nu^4).
 \label{eq:c4v13-dipole}
\end{align}
The constants scale as indicated for other points of the region
(\(C_I^2\) in the first three, \(\kappa^2C_I^2\) in the last two).
\end{theorem}

\begin{proof}
\(|\Omega|^2\le2|\nabla V|^2\) and \(G\le1\) give \(\mathcal E\le2\|\nabla V\|_2^2\le2C\).  \(\mathcal E_1\): Theorem
C3-V93-rows(2) with \(\theta_1^{ST}=\frac12\) at the floor (Theorem~\ref{thm:c4v4-gamma2}).
\(\mathcal P\): \(|\nabla\Omega|^2\le2|\nabla^2V|^2\le2C_\infty''^2\) and \(|G|=44.55\nu^{3/2}\).  The mean row:
Theorem C3-V74-extremal, first inequality, with \(\sup\mathcal E_1\); the trivial
bound is the Poincar\'e decomposition \(\mathcal E=|G||\bar\Omega|^2+\int|\Omega-\bar\Omega|^2G\).  The
dipole row: Theorem C3-V78-extremal with \(\mathcal E\le2C\), \(\mathcal E_1\le12\nu C\); the trivial
bound is the Hermite norm \(|D|^2/(2\nu|G|)=\|P_1\Omega\|_G^2\le\mathcal E\) (Proposition
C3-V78-poincare).  Values: \(\kappa^2C_I^2=0.440\), \(|G|=44.55\nu^{3/2}\),
\(0.440\cdot44.55\cdot19.56=383\); \(2\cdot44.55\cdot7.45=664\).
\end{proof}

\begin{proposition}[The (E5-small) reading at the floor]
\label{prop:c4v13-E5}
At the floor, \(C_\infty'/\kappa\le\gamma_1^{\rm anc}=3.35\).  In the third cycle's analysis under
(E5-small) (Theorem C3-V73-window), the palinstrophy row applies at
the floor only for \(\gamma_1<1.094\) and the concentration row only for
\(\gamma_1>0.807\); the explicit value lies above both thresholds, so the
concentration row applies and the palinstrophy row is void, as the
strain--transport form already gave unconditionally
(Proposition C3-V93-consequences).  The loop constraints of C3-V74,
C3-V77 are void.
\end{proposition}

\begin{proof}
Read off the thresholds.
\end{proof}

% =============================================================================
\section{The extremal defect row is void}
\label{sec:c4v13-defect}
% =============================================================================

\begin{theorem}[The defect row of C3-V38 and C3-V48 is void]
\label{thm:c4v13-void}
The bound \(\sup\|Y\|_G^2\le C_{YP'}/(\frac12-C_\infty'-\frac{\sqrt3}2\kappa-\frac{\kappa^2}2)\) of Theorem
C3-V48-rows(1) requires \(C_\infty'+\frac{\sqrt3}2\kappa+\frac{\kappa^2}2<\frac12\), which contradicts the
enstrophy constraint \(C_\infty'+\frac{\sqrt3}2\kappa+\frac{\kappa^2}2\ge1\) (Theorem C3-V67-extremal);
in strain--transport form the same row has threshold
\(\frac12-(2/3)^{1/2}C_\infty'-\frac{\kappa^2}2\), which contradicts \((2/3)^{1/2}C_\infty'+\frac{\kappa^2}2\ge1\).  Hence the
extremal defect row gives no bound at any homogeneous type-I point;
the identity at the maximizer of \(\|Y\|_G^2\) stands, its right side
being bounded only by the class.
\end{theorem}

\begin{proof}
The two constraints; the strain--transport form of the row's
stretching and transport terms follows from Lemma C3-V91-strain
and Lemma C3-V93-bounds applied to the field \(Y\) (the quadratic form
\(\langle Y,(Y\cdot\nabla)V\rangle_G\) and the transport \(\langle Y,(V\cdot\nabla)Y\rangle_G\)).
\end{proof}

\begin{remark}[The fourth bound of C3-V59]
\label{rem:c4v13-fourth}
Proposition C3-V59-fourth bounds the defect rate on the anchor-alone
event by \(2C_Y[\kappa(2C_{\mathcal L})^{1/2}+\frac{\rho_0}{2\nu}(\frac13+(\frac23)^{1/2})C_\infty K_6|U''|^{1/3}]\); with the values
of V11--V12 (\(C_Y\le270.6\nu^{5/4}\), \(C_{\mathcal L}\le\frac12C_I^2\nu^{5/2}+11.14\kappa^2\nu^{5/2}\), \(K_6\le S_3C_I\nu^{3/4}\)) every
constant is explicit except the anchor radius \(\rho_0\) and the halo
volume \(|U''|\), which are local data of the configuration; the bound
is recorded as explicit in those two.
\end{remark}

\begin{assessment}[Reading]
\label{ass:c4v13-reading}
The vorticity rows that bound Gaussian quantities from above become
numbers at the floor corner, and two of them, the mean and the
dipole, improve on their trivial forms once the concentration bound
is used; the palinstrophy bound is dominated by the second-derivative
constant and is far above the trivial one.  The extremal defect row,
which the third cycle had left conditional on a positive threshold,
is structurally void: the enstrophy constraint forbids the smallness
of the strain that the row would need.  This joins the palinstrophy
row (void for \(\kappa<1\)) in the list of rows whose extremal evaluation
cannot close under the class bounds: the damping of \(\|Y\|^2\) is \(\frac12\), of
the enstrophy \(1\), and the stretching must exceed the latter.
\end{assessment}

% =============================================================================
\section{Integrated V13 conclusion and next acquisition order}
\label{sec:c4v13-conclusion}
% =============================================================================

\begin{theorem}[What V13 proves]
\label{thm:c4v13-conclusion}
Relative to the three legacy cycles and V1--V12: the vorticity rows
with values (Theorem~\ref{thm:c4v13-rows}), the (E5-small) reading
(Proposition~\ref{prop:c4v13-E5}) and the voidness of the extremal defect
row (Theorem~\ref{thm:c4v13-void}).  The full Navier--Stokes
stability/global-regularity theorem remains unproved.
\end{theorem}

\begin{proof}
The cited results.
\end{proof}

\begin{programme}[V14 acquisition order]
\label{prog:c4v14}
\begin{enumerate}[label=\textup{(AC\arabic*)},leftmargin=4.8em]
\item The quantitative description of the hypothetical homogeneous
      type-I point in one table: the certified region, the derivative
      constants, the pressure and defect norms, the Gaussian energy,
      dissipation and vorticity quantities, with their values at the
      corner and their scaling in \((\kappa,C_I)\); the list of the rows that
      stay symbolic and why (\(c_\Phi\), \(\underline\delta\), \(K_P\), local data).
\item The list of rows found void by the cycle (palinstrophy, loops,
      extremal defect) with the structural reason.
\item The next compulsory companion audit is V15.
\end{enumerate}
\end{programme}

\begin{assessment}[Position after V13]
\label{ass:c4v13-position}
The vorticity rows have values; the defect row is void.  No full
stability or global-regularity claim is made.
\end{assessment}

% =============================================================================
\section*{Version V13 (fourth cycle) append-only record}
\addcontentsline{toc}{section}{Version V13 (fourth cycle) append-only record}
% =============================================================================

The complete V12 source of the fourth cycle was retained before this
continuation.  Its SHA--256 digest was
\begin{center}\scriptsize\ttfamily
80db168aa90c37aec13fc00b5acebc8db7a2f2fd800d5ec11dc29b0c264504bd.
\end{center}
V13 gave the vorticity rows their values and found the extremal
defect row void.  No full regularity claim is made.

% =============================================================================
\section{V14: the quantitative description of the hypothetical singular point}
\label{sec:c4v14-entry}
% =============================================================================

Items (AC1) and (AC2) of Programme~\ref{prog:c4v14} are carried out:
the certified quantities at a homogeneous type-I point of a smooth
periodic solution are collected in one table with their values at
the floor corner and their scaling in \((\kappa,C_I)\); the rows that stay
symbolic are listed with the reason; the rows found void are listed
with the structural reason.  No global regularity claim is made.  The
next compulsory companion audit is V15.

% =============================================================================
\section{The table}
\label{sec:c4v14-table}
% =============================================================================

\begin{theorem}[Quantitative description at a homogeneous type-I point]
\label{thm:c4v14-table}
Every homogeneous type-I point of a smooth periodic solution with
velocity constant \(\kappa\) and type-I enstrophy constant \(C_I\) satisfies
the following, all certified; the column ``corner'' is the point
\((\kappa,C_I)=(0.3438,1.93)\) of the region's boundary, with \(C_\infty'\le1.152\),
\(c''=\nu^{1/2}C_\infty''\le6.62\); the column ``scaling'' gives the dependence on
\((\kappa,C_I)\) of the bound, the derivative constants entering through
\(\gamma_1^{\rm anc}(\kappa)\), \(\gamma_2^{\rm anc}(\kappa)\) of V3--V4.

\medskip
\begin{center}\small
\begin{tabular}{@{}p{0.34\textwidth}p{0.2\textwidth}p{0.24\textwidth}p{0.12\textwidth}@{}}
\toprule
\textbf{Quantity} & \textbf{Corner} & \textbf{Scaling / formula} & \textbf{Source}\\
\midrule
velocity constant \(\kappa\) & \(\ge0.3438\) & \(\ge F(C_I)\), \(F\) of V7 & V3, V7\\
enstrophy constant \(C_I\) & \(\ge1.93\) at \(\kappa=0.3438\) & \(\ge G(\kappa)\) for \(\kappa<\sqrt2\) & V7, V8\\
gradient constant \(C_\infty'\) & \(\le1.152\) & \(\gamma_1^{\rm anc}(\kappa)\kappa\) & V3\\
strain constant (lower) & \(\ge1.152\)\,\((=\)\,upper at floor\()\) & \(\ge(3/2)^{1/2}(1-\kappa^2/2)\) & C3-V93\\
\(\nu^{1/2}C_\infty''\) & \(\le6.62\) & \(\gamma_2^{\rm anc}(\kappa)\kappa\) & V4\\
pressure \(|P_V|\) & \(\le2.53\,\nu\) & \eqref{eq:c4v12-P0} & V12\\
pressure gradient \(|\nabla P_V|\) & \(\le31.95\,\nu^{1/2}\) & \eqref{eq:c4v12-P} & V12\\
time derivative \(|V_t|\), defect \(|Y(z)|\) & \(\le38.96\,\nu^{1/2}\), \(\le39.14\nu^{1/2}+0.576|z|\) & \eqref{eq:c4v12-t} & V12\\
\(\|N\|_G\), \(\|Y\|_G\) & \(\le215.9\,\nu^{5/4}\), \(\le270.6\,\nu^{5/4}\) & \eqref{eq:c4v12-YN} & V12\\
Gaussian energy \(\varphi\) & \(\le5.27\,\nu^{5/2}\) & \(\le\min\{44.55\kappa^2,40\kappa^2C_I^2\}\nu^{5/2}\) & V11\\
pinching constant \(c_\Phi\) & \(\le17.6\,\nu^{5/2}\) & \(\le40\kappa^2C_I^2\nu^{5/2}\) & V11\\
defect rate \(\underline\delta\) & \(\le4.40\,\nu^{5/2}\) & \(\le10\kappa^2C_I^2\nu^{5/2}\) & V11\\
Gaussian dissipation \(d\) & \(\le14.9\,\nu^{5/2}\) & \(\le4C_I^2\nu^{5/2}\) & V11\\
Dirichlet form \(\mathcal L\) & \(\le3.18\,\nu^{5/2}\) & \(\le(\frac12C_I^2+11.14\kappa^2)\nu^{5/2}\) & V11\\
Gaussian enstrophy \(\mathcal E\) & \(\le7.45\,\nu^{3/2}\) & \(\le2C_I^2\nu^{3/2}\) & V13\\
second vorticity moment \(\mathcal E_1\) & \(\le44.7\,\nu^{5/2}\) & \(\le12C_I^2\nu^{5/2}\) at the floor & V4, V13\\
palinstrophy \(\mathcal P\) & \(\le3904\,\nu^{1/2}\) & \(\le89.1\,c''^2\nu^{1/2}\) & V13\\
mean vorticity \(|G||\bar\Omega|^2\) & \(\le1.32\,\nu^{3/2}\) & \(\le3\kappa^2C_I^2\nu^{3/2}\) & V13\\
vorticity dipole \(|D|^2\) & \(\le383\,\nu^4\) & \(\le871\,\kappa^2C_I^2\nu^4\) & V13\\
\bottomrule
\end{tabular}
\end{center}
\medskip

\noindent
(\(\varphi\le44.55\kappa^2\nu^{5/2}=5.27\nu^{5/2}\) at the corner, the pinching bound \(17.6\) being
larger there; \(\mathcal L\le\frac12\cdot3.725+11.14\cdot0.118=3.18\); \(871=44.55\cdot19.56\).)
\end{theorem}

\begin{proof}
The cited results; the arithmetic at the corner.
\end{proof}

\begin{proposition}[Rows that stay symbolic]
\label{prop:c4v14-symbolic}
The following remain in terms of constants the class does not fix
numerically: every lower bound involving \(c_\Phi\) (the Dirichlet lower
bound \(\mathcal L_-\), the variance cap \(\mathrm{Var}_{\max}\), the push threshold of C3-V24,
the window floor \(v_{\rm win}\)); every statement involving \(\underline\delta\) from below
(N-C3-3); the mean constraint of C2-V36 through \(K_P\); the fourth
defect bound of C3-V59 through the anchor radius \(\rho_0\) and the halo
volume \(|U''|\); the explicit variance floor \(v_{\rm exp}\), which is explicit
but involves \(c_\Phi\), \(C_\Phi\), \(K_6\), \(C\) through \(K_r\) and the growth constants
and is exponentially small.
\end{proposition}

\begin{proof}
Inspection of the cited statements.
\end{proof}

\begin{proposition}[Item (AC2): rows found void]
\label{prop:c4v14-void}
At every homogeneous type-I point the following extremal rows give
no bound, for the structural reason that their damping is below the
stretching forced by the enstrophy constraint: the palinstrophy row
(damping \(\frac32\), stretching \(2(2/3)^{1/2}C_\infty'\), void for \(\kappa<1\); C3-V93); the
loop constraints and floors of C3-V74, V76, V77 (they need the
palinstrophy row); the extremal defect row (damping \(\frac12\), stretching
\((2/3)^{1/2}C_\infty'\); Theorem~\ref{thm:c4v13-void}).  The rows that do close are
those whose damping exceeds the forced stretching: the enstrophy
row itself (damping \(1\), at the threshold), the weighted-enstrophy
row (damping \(\frac32\) against \((2/3)^{1/2}C_\infty'\), applies with margin \(\frac12\)), and
the mean and dipole rows (pure damping with sources bounded by the
concentration).
\end{proposition}

\begin{proof}
The cited statements.
\end{proof}

\begin{assessment}[Reading of the description]
\label{ass:c4v14-reading}
A homogeneous type-I singular point of a smooth periodic solution,
if one exists, is described at the level of the class by the table:
its rescaled velocity is at least a third of \(\nu^{1/2}\), its strain is of
order one, its Gaussian energy is at most a few \(\nu^{5/2}\), its Gaussian
vorticity is concentrated (second moment at most twelve times the
enstrophy scale), its mean vorticity and dipole are small compared
with its enstrophy, and its pressure gradient and defect are at most
a few tens of \(\nu^{1/2}\).  Nothing in the table is contradictory, and the
class fixes neither \(\kappa\) nor \(C_I\); the description is a certified region,
not an exclusion.  The one structural fact the block added is the
classification of the extremal rows by damping against forced
stretching, which says which rows can and cannot close in the class.
\end{assessment}

% =============================================================================
\section{Integrated V14 conclusion and next acquisition order}
\label{sec:c4v14-conclusion}
% =============================================================================

\begin{theorem}[What V14 establishes]
\label{thm:c4v14-conclusion}
Relative to the three legacy cycles and V1--V13: the quantitative
description (Theorem~\ref{thm:c4v14-table}), the symbolic rows
(Proposition~\ref{prop:c4v14-symbolic}) and the void rows
(Proposition~\ref{prop:c4v14-void}).  The full Navier--Stokes
stability/global-regularity theorem remains unproved.
\end{theorem}

\begin{proof}
The cited results.
\end{proof}

\begin{programme}[V15 acquisition order]
\label{prog:c4v15}
\begin{enumerate}[label=\textup{(AC\arabic*)},leftmargin=4.8em]
\item The compulsory companion audit: read the current SM and YM
      notes and record their digests and decision.
\item Record the position after fifteen versions and the direction
      for V16--V20 among (D3)--(D5) of the third cycle's V99; delivery.
\end{enumerate}
\end{programme}

\begin{assessment}[Position after V14]
\label{ass:c4v14-position}
The quantitative description is written.  No full stability or
global-regularity claim is made.
\end{assessment}

% =============================================================================
\section*{Version V14 (fourth cycle) append-only record}
\addcontentsline{toc}{section}{Version V14 (fourth cycle) append-only record}
% =============================================================================

The complete V13 source of the fourth cycle was retained before this
continuation.  Its SHA--256 digest was
\begin{center}\scriptsize\ttfamily
cdf4e06ea9f32e67d7055276ca50a6ffb902fa13d31b5426d411e27ecee13ef7.
\end{center}
V14 wrote the quantitative description.  No full regularity claim
is made.

% =============================================================================
\section{V15: compulsory companion audit and the position after fifteen versions}
\label{sec:c4v15-entry}
% =============================================================================

This is the fifteenth version of the fourth cycle.  The compulsory
review of the two companion programmes is performed first, and the
position after fifteen versions is recorded with the direction for
V16--V20.  No global regularity claim is made.  The next compulsory
companion audit is V20.

% =============================================================================
\section{Compulsory V15 review of the current companion notes}
\label{sec:c4v15-companion-audit}
% =============================================================================

The current companion files in \texttt{latex\_notes} were inspected
read-only.  Since the V10 audit the SM programme has moved from V57
to V58 of its seventh series, the YM programme from V32 to V33 of its
sixth series.  The frozen checkpoints are
\begin{align}
 &\texttt{Renewal\_SM\_Einstein\_Continuum\_Research\_Notes\_V58.tex},
 \notag\\[-2pt]
 &\qquad
 \texttt{eab9a2107808e75e863046422c5680b9acfb9d0c046079668da6a6dc73489112},
 \label{eq:c4v15-SM-checkpoint}\\
 &\texttt{Renewal\_Yang\_Mills\_Mass\_Gap\_Research\_Notes\_V33.tex},
 \notag\\[-2pt]
 &\qquad
 \texttt{fdd4ffcaf4cdb9fe321e0f854b7872a6e6449bdc8c9afe55a35d5002a49b2d94}.
 \label{eq:c4v15-YM-checkpoint}
\end{align}
(SM V58 has 23311 lines; YM V33 has 15029 lines.  The folder also
held YM V32 with digest \texttt{4db29068\dots} and 14580 lines, differing from
the state recorded at V10 (\texttt{1d61f4bc\dots}, 14579 lines): the YM programme
edited its V32 after that audit; recorded, nothing depends on it.)

\emph{SM V58.}  The SM programme treated the common renormalized
Einstein heat stress: a metric-gradient convention, the complete
logarithmic heat gradient, lower owners and renormalization
conditions, an off-shell conservation identity, the natural
finite-energy extension, and an exact no-escape card.

\emph{YM V33.}  The YM programme treated an optimal polynomial scale
corridor and a generated-interaction compiler.

\begin{theorem}[V15 companion-audit decision]
\label{thm:c4v15-companion-decision}
Nothing is imported from SM V58 or YM V33.  Neither bears on the
fourth cycle's blocks.
\end{theorem}

\begin{proof}
The SM results concern heat-trace stresses of lattice Dirac
interfaces; the YM results concern scale corridors of a lattice
gauge theory.  None is a Navier--Stokes statement, and none concerns
the mild formulation, the kernels, the joint region or the ledger
with values.
\end{proof}

% =============================================================================
\section{Position after fifteen versions, and the direction for V16--V20}
\label{sec:c4v15-position}
% =============================================================================

\begin{assessment}[Position after fifteen versions of the fourth cycle]
\label{ass:c4v15-position-fifteen}
Three blocks are written: the ancient-form block (V2--V4), the
far-field block (V6--V8) and the ledger with values (V11--V14).  The
certified description of a homogeneous type-I point is
Theorem~\ref{thm:c4v14-table}; the structural results are the
classification of the extremal rows by damping against forced
stretching (Proposition~\ref{prop:c4v14-void}) and the joint region
(Theorem~\ref{thm:c4v8-region}).  The block V16--V20 takes direction (D3)
of the third cycle's V99: the gate G1 on the spread anchor
(Theorem C3-V68-gate) with the values of the ledger, i.e.\ the
explicit bounds on the Lamb pairing and the pressure work on the
anchor in terms of \((\kappa,C_I)\) and the local data of the configuration,
and the question whether any part of the region forces the sign;
the negative is to be recorded if the class does not fix the
orientation, as the third cycle found symbolically.  V16 the gate's
inequality with values; V17 the sign question by cases; V18 the
nearly constant and spread anchors with the concentration bound
(the second vorticity moment at the anchor); V19 consolidation; V20
audit.  No global regularity claim is made.
\end{assessment}

% =============================================================================
\section{Integrated V15 conclusion and next acquisition order}
\label{sec:c4v15-conclusion}
% =============================================================================

\begin{theorem}[What V15 establishes]
\label{thm:c4v15-conclusion}
Relative to the three legacy cycles and V1--V14: the companion-audit
decision (Theorem~\ref{thm:c4v15-companion-decision}) and the position.
The full Navier--Stokes stability/global-regularity theorem remains
unproved.
\end{theorem}

\begin{proof}
The cited records.
\end{proof}

\begin{programme}[V16 acquisition order]
\label{prog:c4v16}
\begin{enumerate}[label=\textup{(AC\arabic*)},leftmargin=4.8em]
\item Restate the gate on the spread anchor (Theorem C3-V68-gate and
      Theorem C3-V58-gate) with its two order-one quantities, and
      bound each with the values of V11--V13 (\(C_Y\), \(C_{\mathcal L}\), \(C_P\), the
      concentration bound, \(K_6=S_3C_I\nu^{3/4}\)) in terms of \((\kappa,C_I)\) and the
      anchor's local data (radius, halo volume, mass on the anchor
      event).
\item Record the resulting inequality at the corner and along the
      region's boundary.
\item The next compulsory companion audit is V20.
\end{enumerate}
\end{programme}

\begin{assessment}[Position after V15]
\label{ass:c4v15-position}
The companion audit is recorded; the gate block begins.  No full
stability or global-regularity claim is made.
\end{assessment}

% =============================================================================
\section*{Version V15 (fourth cycle) append-only record}
\addcontentsline{toc}{section}{Version V15 (fourth cycle) append-only record}
% =============================================================================

The complete V14 source of the fourth cycle was retained before this
continuation.  Its SHA--256 digest was
\begin{center}\scriptsize\ttfamily
5645ba9574708c81664e927e7b5e42f19dd1877fe1fca1791a1effab92e457df.
\end{center}
V15 performed the compulsory companion audit and recorded the
position.  No full regularity claim is made.

% =============================================================================
\section{V16: the gate on the anchor with values}
\label{sec:c4v16-entry}
% =============================================================================

Items (AC1) and (AC2) of Programme~\ref{prog:c4v16} are carried out.
The gate on the anchor-alone event (Theorems C3-V58-gate and
C3-V68-gate) bounds the defect at an extremal-Dirichlet configuration
by the Lamb pairing's constant \(\kappa(2\mathcal L)^{1/2}\) and the pressure work's
constant \(\mathcal M_{U_0}\); with the values of V11--V13 the first is at most
\(0.867\,\nu^{5/4}\) at the corner of the region and the second is
\(0.163\,\nu^{5/4}\) times the product of the anchor radius and the halo
diameter in profile coordinates; hence at such a configuration
\(\|Y\|_G\le(0.867+0.163\,r_0u)\nu^{5/4}\), against the class bound \(270.6\,\nu^{5/4}\).  The
inequality is recorded along the region's boundary.  No global
regularity claim is made.  The next compulsory companion audit is
V20.

% =============================================================================
\section{The two constants of the gate with values}
\label{sec:c4v16-constants}
% =============================================================================

Notation of the third cycle's V56--V68: \(\mathcal E_0\) the anchor-alone event,
\(U_0\) the anchor's neighbourhood of radius \(\rho_0\), \(U''\) the halo, \(|U''|\) its
volume, \(\mathcal A_{U_0}\) the Lamb pairing and \(\mathcal B_{U_0}\) the pressure work on \(U_0\),
\(\mathcal M_{U_0}\) of (C3-V58-bound); \(r_0:=\rho_0\nu^{-1/2}\) and \(u:=|U''|^{1/3}\nu^{-1/2}\) the anchor radius
and the halo diameter in units of \(\nu^{1/2}\) (the profile coordinates).

\begin{proposition}[The two constants at the corner]
\label{prop:c4v16-constants}
At every homogeneous type-I point with \((\kappa,C_I)\) at the corner
\((0.3438,1.93)\), on the trapped sub-windows of \(\mathcal E_0\):
\begin{equation}
 \kappa(2\mathcal L)^{1/2}\le\kappa(2C_{\mathcal L})^{1/2}\le0.867\,\nu^{5/4},\qquad
 \mathcal M_{U_0}\le\frac12\Bigl(\frac13+\Bigl(\frac23\Bigr)^{1/2}\Bigr)\kappa S_3C_I\,r_0u\,\nu^{5/4}\|Y\|_{L^2(G;U_0)}\le0.163\,r_0u\,\nu^{5/4}\|Y\|_{L^2(G;U_0)} ,
 \label{eq:c4v16-constants}
\end{equation}
and in general \(\kappa(2C_{\mathcal L})^{1/2}=\kappa(C_I^2+22.28\kappa^2)^{1/2}\nu^{5/4}\), \(\mathcal M_{U_0}\le0.2457\,\kappa C_I\,r_0u\,\nu^{5/4}\|Y\|_{L^2(G;U_0)}\).
\end{proposition}

\begin{proof}
\(C_{\mathcal L}\le(\frac12C_I^2+11.14\kappa^2)\nu^{5/2}\) (Lemma~\ref{lem:c4v11-trivial}), so
\(\kappa(2C_{\mathcal L})^{1/2}=\kappa(C_I^2+22.28\kappa^2)^{1/2}\nu^{5/4}=0.3438(3.725+2.634)^{1/2}=0.867\).  In
(C3-V58-bound), \(\|V\|_{L^4(U'')}^2\le C_\infty K_6|U''|^{1/3}\) with \(C_\infty=\kappa\nu^{1/2}\), \(K_6\le S_3C_I\nu^{3/4}\)
(Lemma~\ref{lem:c4v6-holder}), so
\(\mathcal M_{U_0}\le\frac1{2\nu}(\frac13+(\frac23)^{1/2})\kappa S_3C_I\nu^{5/4}|U''|^{1/3}\rho_0\|Y\|_{L^2(G;U_0)}\), and \(|U''|^{1/3}\rho_0=r_0u\,\nu\);
\(\frac12(\frac13+0.8165)\cdot0.4273=0.2457\); \(0.2457\cdot0.3438\cdot1.93=0.163\).
\end{proof}

% =============================================================================
\section{The gate with values}
\label{sec:c4v16-gate}
% =============================================================================

\begin{theorem}[The gate on the anchor at an extremal-Dirichlet configuration, with values]
\label{thm:c4v16-gate}
At a homogeneous type-I point with \((\kappa,C_I)\), if a maximizer or
minimizer \(W^*\) of \(\mathcal L\) over the family lies in \(\mathcal E_0\) on a trapped
sub-window, then at \(W^*\)
\begin{equation}
 \boxed{\quad
 \|Y\|_G\ \le\ \Bigl[\kappa(C_I^2+22.28\kappa^2)^{1/2}+0.2457\,\kappa C_I\,r_0u\Bigr]\nu^{5/4}+C_{11}\varepsilon_{\rm th}+o(1),
 \quad}
 \label{eq:c4v16-gate}
\end{equation}
at the corner \(\|Y\|_G\le(0.867+0.163\,r_0u)\nu^{5/4}+C_{11}\varepsilon_{\rm th}+o(1)\), the \(o(1)\) being the
far-cell and Gaussian-tail terms of Theorem C3-V58-gate divided by
\(\|Y\|_G\).  Along the region's boundary the Lamb constant
\(\kappa(C_I^2+22.28\kappa^2)^{1/2}\) is \(0.867\), \(1.01\), \(1.38\), \(2.44\), \(4.78\), \(6.82\), \(9.25\) at
\(\kappa=0.3438,0.4,0.5,0.7,1.0,1.2,1.4\), and the pressure-work constant
\(0.2457\kappa C_I\) is \(0.163\), \(0.166\), \(0.177\), \(0.189\), \(0.179\), \(0.147\), \(0.038\).
\end{theorem}

\begin{proof}
Theorem C3-V58-gate(i) with Proposition~\ref{prop:c4v16-constants} and
\(\|Y\|_{L^2(G;U_0)}\le\|Y\|_G\); divide by \(\|Y\|_G\).  The boundary values from
Theorem~\ref{thm:c4v8-region} by substitution, e.g.\ at \((0.5,1.44)\):
\(0.5\,(2.074+5.57)^{1/2}=1.38\) and \(0.2457\cdot0.5\cdot1.44=0.177\).
\end{proof}

\begin{remark}[Comparison with the class bound]
\label{rem:c4v16-comparison}
The class bounds \(\|Y\|_G\le C_Y=270.6\,\nu^{5/4}\) at the corner (Theorem~\ref{thm:c4v12-YN}).
The gate says that at an extremal-Dirichlet anchor configuration the
defect is smaller by two orders of magnitude, unless the product
\(r_0u\) of the anchor radius and the halo diameter is of order \(10^3\); the
third cycle bounded the halo diameter by \(2(r_2+R_{c_\Phi'}+3)+1\) in cell units,
where \(R_{c_\Phi'}\) is the Gaussian radius that captures the pinching
energy, a quantity in \(c_\Phi\), so \(u\) is explicit only in the symbolic
pinching constant; \(r_0\) is the anchor radius of the configuration.
The gate is therefore a bound on the defect at a special
configuration, in terms of two local numbers, and it is the only
place in the series where \(\|Y\|_G\) is bounded by something of order
one.
\end{remark}

\begin{assessment}[Item (AC2): reading]
\label{ass:c4v16-reading}
With values, the gate on the anchor reads: at an extremal-Dirichlet
spread anchor, \(\|Y\|_G\lesssim(0.9+0.16\,r_0u)\nu^{5/4}\).  The two constants are of
order one, as the third cycle found symbolically, and the Lamb
constant grows along the region's boundary with \(\kappa\), while the
pressure-work constant peaks near \(\kappa=0.7\).  Whether the sign of
\(\mathcal A_{U_0}+\mathcal B_{U_0}\) can be forced in any part of the region is the question
of V17; the magnitudes alone do not force it, both terms being
bounded in absolute value only.
\end{assessment}

% =============================================================================
\section{Integrated V16 conclusion and next acquisition order}
\label{sec:c4v16-conclusion}
% =============================================================================

\begin{theorem}[What V16 proves]
\label{thm:c4v16-conclusion}
Relative to the three legacy cycles and V1--V15: the two constants
with values (Proposition~\ref{prop:c4v16-constants}) and the gate with
values (Theorem~\ref{thm:c4v16-gate}).  The full Navier--Stokes
stability/global-regularity theorem remains unproved.
\end{theorem}

\begin{proof}
The cited results.
\end{proof}

\begin{programme}[V17 acquisition order]
\label{prog:c4v17}
\begin{enumerate}[label=\textup{(AC\arabic*)},leftmargin=4.8em]
\item The sign question by cases: at an extremal-Dirichlet anchor
      configuration, \(-\|Y\|_G^2=\mathcal A_{U_0}+\mathcal B_{U_0}+E\); (a) if \(\|Y\|_G\) exceeds the
      right side of \eqref{eq:c4v16-gate}, no such configuration exists
      (the extremum of \(\mathcal L\) is not on \(\mathcal E_0\)); (b) otherwise the sign is
      carried by the Lamb pairing or by the pressure work; use the
      sign structure of C3-V42 and C3-V57 (one and two cells) with the
      values, and record whether a sub-case forces a contradiction
      or the negative.
\item The consequences for the blow-up measure: which extremal
      configurations can lie on \(\mathcal E_0\).
\item The next compulsory companion audit is V20.
\end{enumerate}
\end{programme}

\begin{assessment}[Position after V16]
\label{ass:c4v16-position}
The gate has values.  No full stability or global-regularity claim
is made.
\end{assessment}

% =============================================================================
\section*{Version V16 (fourth cycle) append-only record}
\addcontentsline{toc}{section}{Version V16 (fourth cycle) append-only record}
% =============================================================================

The complete V15 source of the fourth cycle was retained before this
continuation.  Its SHA--256 digest was
\begin{center}\scriptsize\ttfamily
cf0fd8c0cd9553970a89cf12d57e65ade666f836988825ef643a8433533061c2.
\end{center}
V16 gave the gate on the anchor its values.  No full regularity
claim is made.

% =============================================================================
\section{V17: the sign question by cases, and the dichotomy at extremal configurations}
\label{sec:c4v17-entry}
% =============================================================================

Items (AC1) and (AC2) of Programme~\ref{prog:c4v17} are carried out.
At an extremal-Dirichlet configuration on the anchor-alone event the
anti-alignment \(-\|Y\|_G^2\) is realized by the Lamb pairing and the
pressure work; the Lamb pairing is bounded through the component of
the defect normal to the velocity--vorticity plane and vanishes on
Beltrami-type and irrotational configurations, the pressure work has
the one-cell and two-cell sign structures of the third cycle's V57.
With the values of V16 neither sign is forced by the class in any
part of the region: the negative record of C3-V68 stands with
numbers.  What the values give is a dichotomy for the blow-up
measure: either no extremum of the Dirichlet form lies on the
anchor-alone event, or the defect at such an extremum is at most
\((0.867+0.163\,r_0u)\nu^{5/4}\) at the corner.  No global regularity claim is
made.  The next compulsory companion audit is V20.

% =============================================================================
\section{Item (AC1): the cases}
\label{sec:c4v17-cases}
% =============================================================================

At an extremal-Dirichlet configuration \(W^*\in\mathcal E_0\) on a trapped
sub-window, Theorem C3-V68-gate(ii) reads
\(-\|Y\|_G^2=\mathcal A_{U_0}+\mathcal B_{U_0}+E\), \(|E|\le\mathcal S'+C_{11}\varepsilon_{\rm th}\|Y\|_G\).

\begin{proposition}[The Lamb pairing]
\label{prop:c4v17-lamb}
\(\mathcal A_{U_0}=\int_{U_0}\ell\cdot Y^\perp G\) with \(\ell=\Omega\times V\) and \(Y^\perp\) the component of \(Y\)
normal to \(\operatorname{span}(V,\Omega)\); \(|\mathcal A_{U_0}|\le\kappa\nu^{1/2}\mathcal E^{1/2}\|Y^\perp\|_{L^2(G;U_0)}\le0.938\,\nu^{5/4}\|Y^\perp\|_{L^2(G;U_0)}\)
at the corner (and \(\le\kappa(2C_{\mathcal L})^{1/2}\|Y\|_G=0.867\nu^{5/4}\|Y\|_G\) by the route of V16).
It vanishes identically wherever \(\Omega\parallel V\) (Beltrami type),
\(\Omega=0\) (irrotational), or \(Y\in\operatorname{span}(V,\Omega)\); its sign is that of the defect's
normal component against the Lamb vector, which the class does not
fix (Proposition C3-V42-pointwise, Proposition C3-V43-beltrami).
\end{proposition}

\begin{proof}
Proposition C3-V42-pointwise with \(\|\Omega\|_G=\mathcal E^{1/2}\le(2C)^{1/2}=(2C_I^2)^{1/2}\nu^{3/4}\)
(Theorem~\ref{thm:c4v13-rows}): \(0.3438\cdot(7.45)^{1/2}=0.938\).
\end{proof}

\begin{proposition}[The pressure work]
\label{prop:c4v17-pressure}
\(\mathcal B_{U_0}=-\frac1{6\nu}\int\chi_{U'}|V|^2f_{U_0}+\frac1{2\nu}\int\chi_{U'}(V\otimes V)^\circ:H_{U_0}+E_{\rm far}\) (Theorem C3-V58-group),
with \(f_{U_0}=(z\cdot Y)G\) on \(U_0\).  For one cell with outward radial defect
(\(z\cdot Y\ge0\)) the trace part is \(\le0\) and the pressure work is negative
whenever \(\frac13\int\chi_{U'}|V|^2f_\alpha>(\frac23)^{1/2}\|V\|_{L^4(U'')}^2\|f_\alpha\|_{L^2}\) (Theorem C3-V57-one-cell);
for two cells of opposite mass the cross part is
\(\frac m{8\pi\nu\rho^3}(Q_\alpha-Q_\beta)+E_{\rm mult}\), with the axial anisotropies \(Q_\gamma\) of the
velocity tensor (Theorem C3-V57-two-cell).  With the values, the
whole of \(\mathcal B_{U_0}\) is at most \(0.163\,r_0u\,\nu^{5/4}\|Y\|_{L^2(G;U_0)}\) in absolute value at
the corner (Proposition~\ref{prop:c4v16-constants}); the sign is that of
the radial defect's direction against the velocity tensor's
anisotropy, which the class does not fix.
\end{proposition}

\begin{proof}
The cited theorems and Proposition~\ref{prop:c4v16-constants}.
\end{proof}

\begin{firewall}[Item (AC1): the sign is not forced in any part of the region]
\label{fw:c4v17-sign}
For every \((\kappa,C_I)\in\mathcal R\) and every value of the local data \((r_0,u)\), the
class admits at an extremal-Dirichlet anchor configuration both
signs of \(\mathcal A_{U_0}\) (the normal component of the defect is free) and both
signs of \(\mathcal B_{U_0}\) (the radial defect's direction and the velocity
tensor's anisotropy are free), and the magnitudes of V16 are upper
bounds only.  No sub-case forces a contradiction: the Beltrami-type
and irrotational sub-cases kill the Lamb pairing and leave the
pressure work, whose sign is free; the one-cell and two-cell
sub-cases give the sign of the pressure work in terms of quantities
the class does not bound below.  The negative record of
Assessment C3-V68-order-one stands with the values of V16.  What
would lift it: a sign for the normal component of the defect against
the Lamb vector from the equation (the angle statement of Theorem
C3-V39-angle), or a floor on the radial defect's mass on the anchor.
\end{firewall}

% =============================================================================
\section{Item (AC2): the dichotomy at extremal configurations}
\label{sec:c4v17-dichotomy}
% =============================================================================

\begin{theorem}[Dichotomy]
\label{thm:c4v17-dichotomy}
At a homogeneous type-I point with \((\kappa,C_I)\), let \(W^*\) be a maximizer or
minimizer of \(\mathcal L\) over \(\mathcal F(x_0)\) (it exists, Lemma C3-V34-extremal).  Then
either \(W^*\notin\mathcal E_0\) on every trapped sub-window (the extremum of the
Dirichlet form is not an anchor-alone configuration), or
\begin{equation}
 \|Y(W^*)\|_G\ \le\ \Bigl[\kappa(C_I^2+22.28\kappa^2)^{1/2}+0.2457\,\kappa C_I\,r_0u\Bigr]\nu^{5/4}+C_{11}\varepsilon_{\rm th}+o(1),
 \label{eq:c4v17-dichotomy}
\end{equation}
i.e.\ at the corner \(\|Y(W^*)\|_G\le(0.867+0.163\,r_0u)\nu^{5/4}\) up to the small terms.
In the second case the defect at the extremal configuration is
smaller than the class bound \(C_Y\) by the factor
\(270.6/(0.867+0.163\,r_0u)\).
\end{theorem}

\begin{proof}
Theorem~\ref{thm:c4v16-gate}.
\end{proof}

\begin{assessment}[Reading]
\label{ass:c4v17-reading}
The gate with values does not decide the sign, but it turns the
anchor-alone event into a quantitative dichotomy for the extremal
configurations of the Dirichlet form: they avoid the anchor-alone
event, or they carry a defect of order one in \(\nu^{5/4}\) (times the
local factor \(r_0u\)).  The first alternative says that the maximal and
minimal Dirichlet forms are attained on configurations with several
groups or on the far field, the second that the defect is small
where the Dirichlet form is extremal, which is the direction of
Gate G1 (a small defect at extremal configurations) without being
its statement (the gate needs the defect rate, a mean over the
family, not a value at one configuration).  The consequence for the
blow-up measure is recorded in V19.
\end{assessment}

% =============================================================================
\section{Integrated V17 conclusion and next acquisition order}
\label{sec:c4v17-conclusion}
% =============================================================================

\begin{theorem}[What V17 establishes]
\label{thm:c4v17-conclusion}
Relative to the three legacy cycles and V1--V16: the Lamb pairing and
pressure work with values (Propositions~\ref{prop:c4v17-lamb},
\ref{prop:c4v17-pressure}), the negative record (Firewall~\ref{fw:c4v17-sign})
and the dichotomy (Theorem~\ref{thm:c4v17-dichotomy}).  The full
Navier--Stokes stability/global-regularity theorem remains
unproved.
\end{theorem}

\begin{proof}
The cited results.
\end{proof}

\begin{programme}[V18 acquisition order]
\label{prog:c4v18}
\begin{enumerate}[label=\textup{(AC\arabic*)},leftmargin=4.8em]
\item Concentration at the anchor with values: from \(\sup\mathcal E_1\le12C_I^2\nu^{5/2}\),
      the Gaussian enstrophy outside radius \(\tilde r\nu^{1/2}\) is at most
      \(12C_I^2\nu^{3/2}/\tilde r^2\), against the total \(2C_I^2\nu^{3/2}\): the vorticity is
      concentrated within \(\tilde r=\sqrt6\); the mean vorticity and dipole
      with values; what this says for the anchor (within the Gaussian
      radius) and for the far field.
\item The nearly constant anchor: with the values, the exclusion of
      C3-V62 restated quantitatively.
\item The next compulsory companion audit is V20.
\end{enumerate}
\end{programme}

\begin{assessment}[Position after V17]
\label{ass:c4v17-position}
The sign is not forced; the dichotomy is recorded.  No full
stability or global-regularity claim is made.
\end{assessment}

% =============================================================================
\section*{Version V17 (fourth cycle) append-only record}
\addcontentsline{toc}{section}{Version V17 (fourth cycle) append-only record}
% =============================================================================

The complete V16 source of the fourth cycle was retained before this
continuation.  Its SHA--256 digest was
\begin{center}\scriptsize\ttfamily
81e4078c8b367f55cdeccd663e00a29cf52c55670347cc1a797205b5a5e744cb.
\end{center}
V17 examined the sign question and recorded the dichotomy.  No full
regularity claim is made.

% =============================================================================
\section{V18: concentration at the anchor, and the nearly constant anchor with values}
\label{sec:c4v18-entry}
% =============================================================================

Items (AC1) and (AC2) of Programme~\ref{prog:c4v18} are carried out.
The second vorticity moment bound of V4 and V13 says that the
Gaussian enstrophy outside the radius \(\tilde r\nu^{1/2}\) is at most \(6/\tilde r^2\)
times the class bound on the total, so that in Gaussian units the
vorticity of a homogeneous type-I point lies, up to a tenth of the
class bound, within radius \(8\); the mean vorticity and the dipole
are small against the pointwise bound.  The nearly constant anchor
exclusion of the third cycle's V62 is restated with the values: at
an extremal-Dirichlet nearly constant anchor the defect is at most
an explicit multiple of \(\eta^{1/2}\nu^{5/4}\) plus the anchor-tail term
\(0.0493\,\widehat T_2(\hat r)^{1/2}\nu^{5/4}\), which is \(0.002\,\nu^{5/4}\) at \(\hat r=8\); the
\(\eta\)-terms stay symbolic through the level bound \(C_r\), which involves
the pinching constant.  No global regularity claim is made.  The
next compulsory companion audit is V20.

% =============================================================================
\section{Item (AC1): concentration at the anchor with values}
\label{sec:c4v18-concentration}
% =============================================================================

\begin{proposition}[Enstrophy outside a Gaussian radius]
\label{prop:c4v18-outside}
At every homogeneous type-I point with \((\kappa,C_I)\) at the floor
\(\kappa=0.3438\), for every profile of the family and every \(\tilde r>0\),
\begin{equation}
 \int_{|z|>\tilde r\nu^{1/2}}|\Omega|^2G\ \le\ \frac{\sup\mathcal E_1}{\tilde r^2\nu}\ \le\ \frac{12C_I^2}{\tilde r^2}\,\nu^{3/2}
 \ =\ \frac6{\tilde r^2}\cdot\bigl(2C_I^2\nu^{3/2}\bigr),
 \label{eq:c4v18-outside}
\end{equation}
the last factor being the class bound on the total \(\mathcal E\).  At the
corner, the enstrophy outside radius \(\tilde r\) (Gaussian units) is at
most \(7.45\), \(4.97\), \(2.79\), \(1.79\), \(1.24\), \(0.70\) times \(\nu^{3/2}\) for
\(\tilde r=\sqrt6,3,4,5,6,8\), i.e.\ \(100\%\), \(67\%\), \(38\%\), \(24\%\), \(17\%\), \(9\%\) of the class
bound on \(\mathcal E\).  If a profile has \(\mathcal E\ge\vartheta\cdot2C_I^2\nu^{3/2}\), the fraction of its
enstrophy outside radius \(\tilde r\) is at most \(6/(\vartheta\tilde r^2)\).
\end{proposition}

\begin{proof}
Chebyshev: \(|z|^2\ge\tilde r^2\nu\) on the domain, so the left side is at most
\(\mathcal E_1/(\tilde r^2\nu)\), and \(\sup\mathcal E_1\le12C_I^2\nu^{5/2}\) (Theorem~\ref{thm:c4v13-rows}).
\end{proof}

\begin{proposition}[Mean vorticity and dipole against the pointwise bound]
\label{prop:c4v18-mean}
At the corner, \(|\bar\Omega|\le(1.32/44.55)^{1/2}=0.172\) while \(|\Omega|\le\sqrt2\,C_\infty'\le1.63\) pointwise;
\(|D|\le19.6\,\nu^2\), while the trivial bound is \(25.8\,\nu^2\).  In profile units
the mean vorticity is at most a tenth of the pointwise bound.
\end{proposition}

\begin{proof}
Theorem~\ref{thm:c4v13-rows}, \(|G|=44.55\nu^{3/2}\), \(|\Omega|^2\le2|\nabla V|^2\), \(383^{1/2}=19.6\), \(664^{1/2}=25.8\).
\end{proof}

\begin{assessment}[What concentration says for the anchor and the far field]
\label{ass:c4v18-reading-concentration}
For the anchor \(U_0\) of radius \(r_0\nu^{1/2}\): the enstrophy outside it is at
most \(44.7\nu^{3/2}/r_0^2\), so an anchor of radius \(r_0\ge8\) holds all but a
tenth of the class bound; an anchor of radius \(r_0\le\sqrt6\) is not
constrained by concentration at all.  For the far field beyond the
Gaussian radius \(R_G=\hat R\nu^{1/2}\): the enstrophy there is at most \(44.7\nu^{3/2}/\hat R^2\),
a polynomial tail, against the exponential tail \(e^{-\hat R^2/8}\) of the
third cycle's tail terms, which the Gaussian weight gives for the
energy but which the concentration bound does not give for the
enstrophy.  Concentration is a statement about where the enstrophy
is, not about how much there is: \(\mathcal E\) has no lower bound in the class
(Firewall C3-V66-enstrophy), and the fraction statement of
Proposition~\ref{prop:c4v18-outside} needs the hypothesis \(\mathcal E\ge\vartheta\cdot2C\).
\end{assessment}

% =============================================================================
\section{Item (AC2): the nearly constant anchor with values}
\label{sec:c4v18-nearly-constant}
% =============================================================================

Notation of the third cycle's V56--V62: on the trapped sub-windows of
\(\mathcal E_0\) with dominant level \(0\) and \(\eta=16\,\mathrm{Var}\le\eta_1<\frac12\), the cross-term
bound \(|\langle N(V),Y\rangle_G|\le\mathcal S(\eta)\|Y\|_G+\mathcal S'\) of Theorem C3-V62-cross, with
\(\mathcal S(\eta)\) of (C3-V62-S), \(c_r=(\frac18+2C_r^2)^{1/2}\), \(C_r=(4\nu C+2C_\Phi)/(4c_\Phi)\) (C2-V33),
\(T_2(R)\le4\pi e^{-R^2/(4\nu)}(2\nu R^3+12\nu^2R+24\nu^3/R)\) (C3-V56), \(R_0=\hat r\nu^{1/2}\) the anchor tail
radius, \(\rho_2=r_2\nu^{1/2}\) the halo radius.  Put
\(\widehat T_2(\hat r):=4\pi e^{-\hat r^2/4}(2\hat r^3+12\hat r+24/\hat r)\), so \(T_2(R_0)=\widehat T_2(\hat r)\nu^{5/2}\).

\begin{proposition}[The cross-term constant on a nearly constant anchor, with values]
\label{prop:c4v18-S}
At the corner \((0.3438,1.93)\), with \(\varphi\le44.55\kappa^2\nu^{5/2}=5.27\nu^{5/2}\),
\begin{equation}
 \mathcal S(\eta)\ \le\ \Bigl[1.116\,(c_r\eta)^{1/2}+1.579\,\bigl((3+4c_r)\eta\bigr)^{1/2}+1.579\,r_0e^{r_2^2/8}\eta^{1/2}+0.0493\,\widehat T_2(\hat r)^{1/2}\Bigr]\nu^{5/4}+C_{11}\varepsilon_{\rm th},
 \label{eq:c4v18-S}
\end{equation}
where \(0.0493\,\widehat T_2(\hat r)^{1/2}=0.561,\ 0.319,\ 0.136,\ 0.0437,\ 0.0020,\ 3\cdot10^{-5}\) for
\(\hat r=3,4,5,6,8,10\); and \(C_r\le6.36\,\nu^{5/2}/c_\Phi\), so \(c_r\) is explicit only through
the pinching constant \(c_\Phi\) (\(c_r\ge0.79\) in any case, from \(r\ge\frac12\)).  In general
the three \(\eta\)-coefficients are \(\kappa(2\varphi)^{1/2}\), \(2\kappa\varphi^{1/2}\), \(2\kappa\varphi^{1/2}\) and the tail
coefficient is \((\frac14+\frac16)\varphi/(\nu|G|)\), with \(\varphi\le\min\{44.55\kappa^2,40\kappa^2C_I^2\}\nu^{5/2}\).
\end{proposition}

\begin{proof}
In (C3-V62-S): \(\kappa(2\varphi c_r\eta)^{1/2}=0.3438\cdot(10.53)^{1/2}(c_r\eta)^{1/2}=1.116(c_r\eta)^{1/2}\);
\(\varphi|G|^{-1/2}\nu^{-1/2}=5.27/6.675=0.789\) and \(\kappa\varphi^{1/2}=0.3438\cdot2.295=0.789\) (equal because
\(\varphi\le|G|\kappa^2\nu\) is used for \(\varphi\)), giving \(1.579((3+4c_r)\eta)^{1/2}\); the core term
(C3-V61-small): \(\frac{\rho_0}{2\nu}[2(\varphi/|G|)^{1/2}+2C_\infty]e^{\rho_2^2/(8\nu)}(\varphi\eta)^{1/2}\le\frac{r_0}2\cdot4\kappa\nu^{1/2}\cdot e^{r_2^2/8}\cdot2.295\eta^{1/2}\nu^{5/4}\nu^{-1/2}=1.579r_0e^{r_2^2/8}\eta^{1/2}\nu^{5/4}\);
the two constant-level terms \((\frac14+\frac16)\frac{\varphi}{\nu|G|}T_2(R_0)^{1/2}\) with \(\varphi/(\nu|G|)=0.1182\);
\(\frac5{12}\cdot0.1182=0.0493\).  \(C_r\): \(4\nu C=4C_I^2\nu^{5/2}=14.9\nu^{5/2}\), \(C_\Phi\le5.27\nu^{5/2}\),
\((14.9+10.53)/4=6.36\).  The tail values by evaluation of \(\widehat T_2\).
\end{proof}

\begin{theorem}[The exclusion of C3-V62 with values]
\label{thm:c4v18-nearly-constant}
At a homogeneous type-I point at the corner, if a maximizer or
minimizer of \(\mathcal L\) over the family is a nearly constant anchor-alone
configuration (\(\eta\le\eta_1\)), then at it
\begin{equation}
 \|Y\|_G\ \le\ \Bigl[1.116\,(c_r\eta)^{1/2}+1.579\,\bigl((3+4c_r)\eta\bigr)^{1/2}+1.579\,r_0e^{r_2^2/8}\eta^{1/2}+0.0493\,\widehat T_2(\hat r)^{1/2}\Bigr]\nu^{5/4}+C_{11}\varepsilon_{\rm th}+\mathcal S'^{1/2},
 \label{eq:c4v18-nearly-constant}
\end{equation}
with \(\mathcal S'=C_7\varepsilon_{\rm th}^{1/4}\cdot270.6\nu^{5/4}+C_{\rm class}(1+R_G^4)e^{-R_G^2/(8\nu)}\).  As \(\eta\to0\) the bound
tends to \(0.0493\,\widehat T_2(\hat r)^{1/2}\nu^{5/4}+C_{11}\varepsilon_{\rm th}+\mathcal S'^{1/2}\): an extremal-Dirichlet
nearly constant anchor with tail radius \(\hat r=8\) has defect at most
\(0.002\,\nu^{5/4}\) up to the threshold terms, against the class bound
\(270.6\,\nu^{5/4}\).  The blow-up measure statement of Theorem C3-V62-measure
reads \(\mathfrak m(A^c)\ge(\underline\delta-270.6\,\mathcal S(\eta_1)\nu^{5/4}-\mathcal S')/(270.6\,\overline M\nu^{5/4})\) whenever the
numerator is positive, and stays conditional: \(\underline\delta\) has no lower
bound (N-C3-3).
\end{theorem}

\begin{proof}
Theorems C3-V62-extremal and C3-V62-measure with
Proposition~\ref{prop:c4v18-S} and \(C_Y=270.6\nu^{5/4}\) (Theorem~\ref{thm:c4v12-YN}).
\end{proof}

\begin{assessment}[Reading]
\label{ass:c4v18-reading}
With values the nearly constant anchor is the configuration on which
the gate is sharpest: its defect at an extremal-Dirichlet
configuration is of order \(\eta^{1/2}\) plus an anchor tail that is
exponentially small in the tail radius, five orders of magnitude
below the class bound at \(\hat r=8\).  The \(\eta\)-terms carry the level bound
\(C_r\), which the class fixes only through \(c_\Phi\); this is the same
symbolic dependence as in Proposition~\ref{prop:c4v14-symbolic}, and it
is the reason the nearly constant row is not fully numerical.  The
spread anchor (\(\eta>\eta_1\)) is the complement: there the gate of V16
applies with both constants of order one, and V17 recorded that the
sign is not forced.  Together V16--V18 say: the anchor-alone
configurations of extremal Dirichlet form carry a defect at most of
order one (spread) or explicitly small (nearly constant), and the
blow-up measure must charge the complement of the nearly constant
anchors whenever the defect rate exceeds the explicit small
quantities.
\end{assessment}

% =============================================================================
\section{Integrated V18 conclusion and next acquisition order}
\label{sec:c4v18-conclusion}
% =============================================================================

\begin{theorem}[What V18 establishes]
\label{thm:c4v18-conclusion}
Relative to the three legacy cycles and V1--V17: concentration with
values (Propositions~\ref{prop:c4v18-outside}, \ref{prop:c4v18-mean}), the
cross-term constant with values (Proposition~\ref{prop:c4v18-S}) and the
exclusion of C3-V62 with values (Theorem~\ref{thm:c4v18-nearly-constant}).
The full Navier--Stokes stability/global-regularity theorem remains
unproved.
\end{theorem}

\begin{proof}
The cited results.
\end{proof}

\begin{programme}[V19 acquisition order]
\label{prog:c4v19}
\begin{enumerate}[label=\textup{(AC\arabic*)},leftmargin=4.8em]
\item Consolidation of the gate block V16--V18: one statement
      collecting the gate with values, the dichotomy, the nearly
      constant bound, and the consequences for the blow-up measure
      (which extremal configurations can lie on \(\mathcal E_0\)); the negative
      records of the block in one place.
\item The direction for V21--V25 among (D4) (the loop method above
      \(\kappa=1\)) and (D5) (lifts of N-C3-1 and N-C3-3), with the reason.
\item The next compulsory companion audit is V20.
\end{enumerate}
\end{programme}

\begin{assessment}[Position after V18]
\label{ass:c4v18-position}
The nearly constant anchor has values.  No full stability or
global-regularity claim is made.
\end{assessment}

% =============================================================================
\section*{Version V18 (fourth cycle) append-only record}
\addcontentsline{toc}{section}{Version V18 (fourth cycle) append-only record}
% =============================================================================

The complete V17 source of the fourth cycle was retained before this
continuation.  Its SHA--256 digest was
\begin{center}\scriptsize\ttfamily
1159b50c93b5bf24b0170f0188cb3c1463d13c4541b9617eb786e6b6d7e60f25.
\end{center}
V18 gave concentration and the nearly constant anchor their values.
No full regularity claim is made.

% =============================================================================
\section{V19: consolidation of the gate block, and the direction for V21--V25}
\label{sec:c4v19-entry}
% =============================================================================

Items (AC1) and (AC2) of Programme~\ref{prog:c4v19} are carried out.
The gate block V16--V18 is collected in one statement: the gate with
values, the dichotomy, the nearly constant bound, and what they say
about which extremal configurations can lie on the anchor-alone
event; the block's negative records are listed in one place.  The
direction for V21--V25 is (D5) of the third cycle's V99, the two
lifts, with (D4) closed in V21 by a short record, for the reason
given below.  No global regularity claim is made.  The next
compulsory companion audit is V20.

% =============================================================================
\section{Item (AC1): the gate block in one statement}
\label{sec:c4v19-block}
% =============================================================================

\begin{theorem}[The gate block with values]
\label{thm:c4v19-block}
At a homogeneous type-I point of a smooth periodic solution with
\((\kappa,C_I)\in\mathcal R\), let \(W^*\) be a maximizer or minimizer of the Dirichlet
form \(\mathcal L\) over the family \(\mathcal F(x_0)\).  Then exactly one of the following
holds on each trapped sub-window.
\begin{enumerate}[label=\textup{(\roman*)},leftmargin=3.2em]
\item \(W^*\) is not an anchor-alone configuration: it is multi-group
      within the Gaussian radius, or its energy sits at the Gaussian
      radius (the band of C3-V6).  The gate says nothing about it.
\item \(W^*\) is a spread anchor (\(\eta=16\,\mathrm{Var}>\eta_1\)): then
      \(\|Y(W^*)\|_G\le[\kappa(C_I^2+22.28\kappa^2)^{1/2}+0.2457\,\kappa C_I\,r_0u]\nu^{5/4}+C_{11}\varepsilon_{\rm th}+o(1)\),
      at the corner \((0.867+0.163\,r_0u)\nu^{5/4}\) (Theorem~\ref{thm:c4v16-gate}).
\item \(W^*\) is a nearly constant anchor (\(\eta\le\eta_1\)): then
      \(\|Y(W^*)\|_G\le\mathcal S(\eta)+\mathcal S'^{1/2}\) with \(\mathcal S(\eta)\) of \eqref{eq:c4v18-S}, which at
      the corner tends to \(0.0493\,\widehat T_2(\hat r)^{1/2}\nu^{5/4}\) plus threshold terms
      as \(\eta\to0\) (Theorem~\ref{thm:c4v18-nearly-constant}).
\end{enumerate}
In cases (ii) and (iii) the defect at \(W^*\) is below the class bound
\(C_Y=270.6\,\nu^{5/4}\) by the factor \(270.6/(0.867+0.163\,r_0u)\), respectively by
five orders of magnitude at \(\hat r=8\) and small \(\eta\).  In no case is the
sign of \(\mathcal A_{U_0}+\mathcal B_{U_0}\) forced (Firewall~\ref{fw:c4v17-sign}).  For the
blow-up measure \(\mathfrak m\): the set \(A\) of nearly constant anchor-alone
configurations has \(\mathfrak m(A^c)\ge(\underline\delta-270.6\,\mathcal S(\eta_1)\nu^{5/4}-\mathcal S')/(270.6\,\overline M\nu^{5/4})\)
whenever the numerator is positive.
\end{theorem}

\begin{proof}
The trichotomy is the case split of Assessment C3-V64-direction
(anchor-alone or not; on the anchor-alone event, \(\eta\le\eta_1\) or not).
The bounds are Theorems~\ref{thm:c4v16-gate}, \ref{thm:c4v17-dichotomy},
\ref{thm:c4v18-nearly-constant}; the measure statement is Theorem
C3-V62-measure with the values.
\end{proof}

\begin{proposition}[Which extremal configurations can lie on the anchor-alone event]
\label{prop:c4v19-which}
An extremal-Dirichlet configuration can lie on \(\mathcal E_0\) only with a
defect of the size in (ii) or (iii).  Since the class does not bound
\(\|Y\|_G\) below at any single configuration (the defect rate \(\underline\delta\) is a
mean, N-C3-3), neither case is excluded: an extremal configuration
with small defect is admissible, and the block is a description, not
an exclusion.  What it excludes is the combination ``extremal
Dirichlet form, anchor-alone, defect of the order of the class
bound'': at an anchor-alone extremum the defect is at most of order
one in \(\nu^{5/4}\).
\end{proposition}

\begin{proof}
Theorem~\ref{thm:c4v19-block}.
\end{proof}

\begin{theorem}[Negative records of the gate block]
\label{thm:c4v19-negatives}
\begin{enumerate}[label=\textup{(N\arabic*)},leftmargin=3.6em]
\item The sign of the Lamb pairing and of the pressure work is not
      forced by the class in any part of \(\mathcal R\) (Firewall~\ref{fw:c4v17-sign});
      lifts: a sign for the defect's normal component against the
      Lamb vector from the equation, or a floor on the radial
      defect's mass on the anchor.
\item The nearly constant bound is symbolic in the level bound \(C_r\)
      through \(c_\Phi\) (Proposition~\ref{prop:c4v18-S}); lift: a numerical
      lower bound on the Gaussian energy \(\varphi\) (the pinching constant).
\item The measure statement is conditional on \(\underline\delta\) exceeding the
      explicit small quantities (N-C3-3).
\item The gate involves the local data \(r_0\), \(u\), \(\hat r\), \(r_2\) of the
      configuration, which the class bounds only through \(c_\Phi\)
      (Remark~\ref{rem:c4v16-comparison}).
\end{enumerate}
\end{theorem}

\begin{proof}
The cited records.
\end{proof}

% =============================================================================
\section{Item (AC2): the direction for V21--V25}
\label{sec:c4v19-direction}
% =============================================================================

\begin{assessment}[Why (D4) is closed and (D5) is taken]
\label{ass:c4v19-direction}
Direction (D4) of the third cycle's V99 is the loop method above
\(\kappa=1\).  Its loop constraints (C3-V74, V76, V77) have the form
\(3\kappa^2+(1+2\kappa^2)q_P\ge1\) and \(4\delta_{67}+3\kappa^2+2\kappa^2q_P\ge1\), which are satisfied
identically for \(\kappa\ge1/\sqrt3=0.577\) whatever \(q_P\), \(\delta_{67}\ge0\); above \(\kappa=1\)
they carry no information.  Their other content, a derivative bound
from the class constants directly, is what the ancient-form block
already provides for every \(\kappa\) (the bound \(\gamma_1^{\rm anc}(\kappa)\) of V3 exists for
every \(\kappa\) by taking \(\theta\) small, the denominator
\(1-2(8/3)^{1/2}c_O^{s0}\kappa(\theta/(1+\theta))^{1/2}\) being positive for \(\theta<\theta_\kappa\)).  So (D4)
reduces to the record that the loop method is void on the whole of
\(\mathcal R\); V21 writes that record with the value of \(\gamma_1^{\rm anc}\) along the
region's boundary, and the block V21--V25 takes (D5), the two lifts:
N-C3-1 (concentration in a ball gives no variance floor; lift: a
pointwise use of the equation on the ball) and N-C3-3 (no
quantitative defect-rate inequality; lift: an explicit Liouville
constant for (E9)).  The reason: the mild formulation of the fourth
cycle is exactly a pointwise use of the equation, the ancient form
writing \(W\) at a point as a kernel integral of \(W\otimes W\); whether it
gives a floor on the variance of a profile concentrated in a ball,
or an explicit distance of every profile of the family from the
self-similar profiles, is the question the fourth cycle is equipped
to ask and the third was not.  The expected outcome is recorded
honestly: the third cycle's negative records were structural, and a
lift needs a new quantitative input, not a reformulation.
\end{assessment}

% =============================================================================
\section{Integrated V19 conclusion and next acquisition order}
\label{sec:c4v19-conclusion}
% =============================================================================

\begin{theorem}[What V19 establishes]
\label{thm:c4v19-conclusion}
Relative to the three legacy cycles and V1--V18: the gate block in
one statement (Theorem~\ref{thm:c4v19-block}), the admissible extremal
configurations (Proposition~\ref{prop:c4v19-which}), the negative records
(Theorem~\ref{thm:c4v19-negatives}) and the direction.  The full
Navier--Stokes stability/global-regularity theorem remains
unproved.
\end{theorem}

\begin{proof}
The cited results.
\end{proof}

\begin{programme}[V20 acquisition order]
\label{prog:c4v20}
\begin{enumerate}[label=\textup{(AC\arabic*)},leftmargin=4.8em]
\item The compulsory companion audit: read the current SM and YM
      notes and record their digests and decision.
\item Record the position after twenty versions and the plan for
      V21--V25 under (D5); delivery.
\end{enumerate}
\end{programme}

\begin{assessment}[Position after V19]
\label{ass:c4v19-position}
The gate block is consolidated; the direction is (D5).  No full
stability or global-regularity claim is made.
\end{assessment}

% =============================================================================
\section*{Version V19 (fourth cycle) append-only record}
\addcontentsline{toc}{section}{Version V19 (fourth cycle) append-only record}
% =============================================================================

The complete V18 source of the fourth cycle was retained before this
continuation.  Its SHA--256 digest was
\begin{center}\scriptsize\ttfamily
f1e26879f4d29b840e2547bf064cdd5bafe4ab2c23914c69cfb1ac97b4b91f11.
\end{center}
V19 consolidated the gate block and fixed the direction.  No full
regularity claim is made.

% =============================================================================
\section{V20: compulsory companion audit and the position after twenty versions}
\label{sec:c4v20-entry}
% =============================================================================

This is the twentieth version of the fourth cycle.  The compulsory
review of the two companion programmes is performed first, and the
position after twenty versions is recorded with the plan for
V21--V25 under direction (D5).  No global regularity claim is made.
The next compulsory companion audit is V25.

% =============================================================================
\section{Compulsory V20 review of the current companion notes}
\label{sec:c4v20-companion-audit}
% =============================================================================

The current companion files in \texttt{latex\_notes} were inspected
read-only.  Since the V15 audit the SM programme has not moved (its
V58 is unchanged, digest \eqref{eq:c4v15-SM-checkpoint}); the YM
programme has moved from V33 to V34 of its sixth series.  The frozen
checkpoints are
\begin{align}
 &\texttt{Renewal\_SM\_Einstein\_Continuum\_Research\_Notes\_V58.tex},
 \notag\\[-2pt]
 &\qquad
 \texttt{eab9a2107808e75e863046422c5680b9acfb9d0c046079668da6a6dc73489112},
 \label{eq:c4v20-SM-checkpoint}\\
 &\texttt{Renewal\_Yang\_Mills\_Mass\_Gap\_Research\_Notes\_V34.tex},
 \notag\\[-2pt]
 &\qquad
 \texttt{9e3cd3090f7b839a7f62f545e3ebf78bd5bf9d1a5e1638b5f056375453510098}.
 \label{eq:c4v20-YM-checkpoint}
\end{align}
(SM V58 has 23311 lines, as at V15; YM V34 has 15389 lines.  The
folder also holds YM V32 and V33 with the digests recorded at V15,
unchanged.)

\emph{SM V58.}  Unchanged since the V15 audit; nothing new to
review.

\emph{YM V34.}  The YM programme treated the upstream response
interface of its generated action: an exact marginal coefficient
from two response derivatives at a selected root, the seven-stock
contribution of a plaquette marginal, sufficient scale conditions
for a bounded marginal compiler, and a regular/ramified domain
fork; its claim boundary records that the mass gap is not proved.

\begin{theorem}[V20 companion-audit decision]
\label{thm:c4v20-companion-decision}
Nothing is imported from SM V58 or YM V34.  Neither bears on the
fourth cycle's blocks.
\end{theorem}

\begin{proof}
The SM file is unchanged.  The YM results concern the response
matching of a lattice gauge action; none is a Navier--Stokes
statement, and none concerns the mild formulation, the kernels, the
joint region, the ledger with values or the gate block.
\end{proof}

% =============================================================================
\section{Position after twenty versions, and the plan for V21--V25}
\label{sec:c4v20-position}
% =============================================================================

\begin{assessment}[Position after twenty versions of the fourth cycle]
\label{ass:c4v20-position-twenty}
Four blocks are written: the ancient-form block (V2--V4, floor
\(\kappa\ge0.3438\)), the far-field block (V6--V8, the joint region \(\mathcal R\)), the
ledger with values (V11--V14, Theorem~\ref{thm:c4v14-table}) and the gate
block (V16--V18, Theorem~\ref{thm:c4v19-block}).  The certified
description of a homogeneous type-I point is a region in \((\kappa,C_I)\)
with a table of bounds, and the gate says that an anchor-alone
extremum of the Dirichlet form carries a defect of order one at
most; no sign is forced and no exclusion is obtained.  The block
V21--V25 takes (D5): V21 the record closing (D4) with \(\gamma_1^{\rm anc}\) along
the boundary; V22 the pointwise use of the equation on a ball
through the ancient form, against N-C3-1 (does the kernel integral
of a profile concentrated in a ball bound its variance below?);
V23 the same question for the mean over the family; V24 the
distance from the self-similar profiles through the ancient form,
against N-C3-3 (a self-similar profile satisfies the ancient form
with \(Y=0\); an explicit Liouville constant is a floor on \(\|Y\|_G\) in the
mean); V25 audit and delivery.  No global regularity claim is made.
\end{assessment}

% =============================================================================
\section{Integrated V20 conclusion and next acquisition order}
\label{sec:c4v20-conclusion}
% =============================================================================

\begin{theorem}[What V20 establishes]
\label{thm:c4v20-conclusion}
Relative to the three legacy cycles and V1--V19: the companion-audit
decision (Theorem~\ref{thm:c4v20-companion-decision}) and the position.
The full Navier--Stokes stability/global-regularity theorem remains
unproved.
\end{theorem}

\begin{proof}
The cited records.
\end{proof}

\begin{programme}[V21 acquisition order]
\label{prog:c4v21}
\begin{enumerate}[label=\textup{(AC\arabic*)},leftmargin=4.8em]
\item Close (D4): state that the loop constraints of C3-V74, V76, V77
      hold identically on \(\mathcal R\) (\(\kappa\ge0.3438\) is not enough, \(\kappa\ge1/\sqrt3\) is;
      record where on the boundary they become identities) and that
      above \(\kappa=1\), where the palinstrophy row applies, the derivative
      bound is \(\gamma_1^{\rm anc}(\kappa)\kappa\) with \(\theta\) chosen below \(\theta_\kappa\); tabulate
      \(\gamma_1^{\rm anc}\) and \(\nu^{1/2}C_\infty''\) at \(\kappa=0.5,0.7,1.0,1.2,1.4\).
\item State what the palinstrophy row gives above \(\kappa=1\) with these
      values (Theorem C3-V93-rows with \(\theta_P^{ST}=-\frac12+\kappa^2/2\)).
\item The next compulsory companion audit is V25.
\end{enumerate}
\end{programme}

\begin{assessment}[Position after V20]
\label{ass:c4v20-position}
The companion audit is recorded; the (D5) block begins.  No full
stability or global-regularity claim is made.
\end{assessment}

% =============================================================================
\section*{Version V20 (fourth cycle) append-only record}
\addcontentsline{toc}{section}{Version V20 (fourth cycle) append-only record}
% =============================================================================

The complete V19 source of the fourth cycle was retained before this
continuation.  Its SHA--256 digest was
\begin{center}\scriptsize\ttfamily
1ad841869cf2a4a6634adc1f5e73127d13d57a066d404d7b213ec6d669be183a.
\end{center}
V20 performed the compulsory companion audit and recorded the
position.  No full regularity claim is made.

% =============================================================================
\section{V21: closing direction (D4), and the derivative constants along the boundary}
\label{sec:c4v21-entry}
% =============================================================================

Items (AC1) and (AC2) of Programme~\ref{prog:c4v21} are carried out.
The loop constraints of the third cycle's V74--V77 are identities for
\(\kappa\ge1/\sqrt3\) and need the palinstrophy row below it, where that row
is void; hence the loop method gives nothing on the whole of the
region \(\mathcal R\), and direction (D4) is closed.  The derivative constants
of the ancient-form block are tabulated along the region's boundary:
they grow rapidly with \(\kappa\) because the absorption forces the window
fraction \(\theta\) to zero.  Above \(\kappa=1\) the palinstrophy row applies only if
the gradient constant lies in an explicit window that the class does
not place it in; where it applies it gives a conditional bound on
the palinstrophy, recorded with its value.  No global regularity
claim is made.  The next compulsory companion audit is V25.

% =============================================================================
\section{Item (AC1): the loop method on the region, and the derivative table}
\label{sec:c4v21-loops}
% =============================================================================

\begin{proposition}[The loop constraints are void on the region]
\label{prop:c4v21-loops}
The loop constraints of the third cycle, \(3\kappa^2+(1+2\kappa^2)q_{\mathcal P}\ge1\)
(C3-V74) and \(4\delta_{67}+3\kappa^2+2\kappa^2q_{\mathcal P}\ge1\) (C3-V77), with \(q_{\mathcal P},\delta_{67}\ge0\), hold
identically for \(\kappa\ge1/\sqrt3=0.5774\).  For \(0.3438\le\kappa<0.5774\) they read
\(q_{\mathcal P}\ge(1-3\kappa^2)/(1+2\kappa^2)\), i.e.\ \(q_{\mathcal P}\ge0.522\) at the floor and \(q_{\mathcal P}\ge0.167\) at
\(\kappa=0.5\); but \(q_{\mathcal P}\) is defined through the palinstrophy row, which
is void for \(\kappa<1\) (Proposition C3-V93-consequences), so the constraints
carry no information there either.  On the boundary of \(\mathcal R\) the
identities begin at \(\kappa=0.5774\), where \(C_I\ge G(0.5774)\approx1.3\).  Direction (D4)
of the third cycle's V99 is closed: the loop method is void on the
whole of \(\mathcal R\).
\end{proposition}

\begin{proof}
\(3\kappa^2\ge1\) iff \(\kappa\ge1/\sqrt3\); the rest is Theorem~\ref{thm:c4v13-void} and
Proposition~\ref{prop:c4v14-void}.
\end{proof}

\begin{theorem}[The derivative constants along the boundary]
\label{thm:c4v21-table}
With the certified kernel constants \(c_O^{s0}=2.766\), \(c_{O,2}=1.788\), \(c_{O,3}=1.897\), the
bounds of Theorems~\ref{thm:c4v3-gradient} and \ref{thm:c4v4-gamma2},
minimized over \(\theta\) on a geometric grid, give at the boundary
points of \(\mathcal R\):

\medskip
\begin{center}\small
\begin{tabular}{@{}lcccccc@{}}
\toprule
\(\kappa\) & \(0.3438\) & \(0.5\) & \(0.7\) & \(1.0\) & \(1.2\) & \(1.4\)\\
\midrule
\(\theta_\kappa\) (absorption limit) & \(0.116\) & \(0.0515\) & \(0.0257\) & \(0.0124\) & \(0.00858\) & \(0.00629\)\\
\(\theta^*\) (optimal, first derivative) & \(0.0093\) & \(0.0033\) & \(0.0013\) & \(0.00054\) & \(0.00034\) & \(0.00023\)\\
\(\gamma_1^{\rm anc}(\kappa)\) & \(3.35\) & \(5.63\) & \(8.80\) & \(13.9\) & \(17.5\) & \(21.2\)\\
\(C_\infty'\le\gamma_1^{\rm anc}\kappa\) & \(1.152\) & \(2.82\) & \(6.16\) & \(13.9\) & \(21.0\) & \(29.7\)\\
\(C_\infty'\ge(3/2)^{1/2}(1-\kappa^2/2)\) & \(1.152\) & \(1.072\) & \(0.925\) & \(0.612\) & \(0.343\) & \(0.024\)\\
\(\nu^{1/2}C_\infty''\le\) & \(6.62\) & \(23.2\) & \(70.0\) & \(222\) & \(398\) & \(651\)\\
\bottomrule
\end{tabular}
\end{center}
\medskip

\noindent
\(\theta_\kappa\) is the value at which \(2(8/3)^{1/2}c_O^{s0}\kappa(\theta/(1+\theta))^{1/2}=1\), i.e.\ \(\theta_\kappa=x^2/(1-x^2)\)
with \(x=(2(8/3)^{1/2}c_O^{s0}\kappa)^{-1}=0.1107/\kappa\); the bounds exist for every \(\kappa\).
\end{theorem}

\begin{proof}
Evaluation of \eqref{eq:c4v3-gamma1} and \eqref{eq:c4v4-gamma2} with
\(J_2(\theta)=\theta^{1/2}/(1+\theta)+\arctan\theta^{1/2}\), \(J_3(\theta)=2\theta^{-1/2}-\pi+2\arctan\theta^{1/2}\), on the grid
\(\theta=10^{-6+j/100}\), \(0\le j<600\); the floor row is Proposition
C3-V93-consequences(1).
\end{proof}

\begin{assessment}[Reading of the table]
\label{ass:c4v21-reading-table}
The upper bound on the gradient constant is tight only at the floor,
where it meets the lower bound; it grows like \(\kappa^2\log(1/\theta^*)\) with \(\theta^*\)
forced to zero as \(1/\kappa^2\) by the absorption, so that at \(\kappa=1\) the
class allows \(C_\infty'\) anywhere in \([0.61,13.9]\) and at \(\kappa=1.4\) anywhere in
\([0.02,29.7]\).  The second-derivative bound grows faster (it carries
\(\gamma_1^2\)).  The ancient-form block is therefore a floor argument: it
is sharp where the floor is and loose elsewhere, and the region \(\mathcal R\)
above \(\kappa=0.5\) is described by the far-field block (the curve \(F(C_I)\)),
not by the derivative constants.
\end{assessment}

% =============================================================================
\section{Item (AC2): the palinstrophy row above \texorpdfstring{\(\kappa=1\)}{kappa=1}}
\label{sec:c4v21-palinstrophy}
% =============================================================================

\begin{proposition}[The window of applicability]
\label{prop:c4v21-window}
The palinstrophy row of Theorem C3-V93-rows(3) applies iff
\(\theta_{\mathcal P}^{ST}=\frac32-2(2/3)^{1/2}C_\infty'-\frac{\kappa^2}2>0\), i.e.\ iff
\(C_\infty'<0.6124\,(\frac32-\frac{\kappa^2}2)\).  With the floor \(C_\infty'\ge1.2247\,(1-\frac{\kappa^2}2)\) the
window is nonempty iff \(\kappa>1\): it is \([0.343,0.478)\) at \(\kappa=1.2\) and
\([0.024,0.318)\) at \(\kappa=1.4\), empty at \(\kappa=1\).  The class does not place \(C_\infty'\)
in the window: its upper bound is \(21.0\) and \(29.7\) there
(Theorem~\ref{thm:c4v21-table}).  Where the row applies,
\begin{equation}
 \sup_{\mathcal F(x_0)}\mathcal P\ \le\ 2\Bigl(\frac{C_\infty''}{\theta_{\mathcal P}^{ST}}\Bigr)^2C
 \ =\ \frac{2c''^2C_I^2}{(\theta_{\mathcal P}^{ST})^2}\,\nu^{1/2},\qquad\theta_{\mathcal P}^{ST}\le-\frac12+\frac{\kappa^2}2 ,
 \label{eq:c4v21-palin}
\end{equation}
against the pointwise bound \(\mathcal P\le89.1\,c''^2\nu^{1/2}\); the ratio is
\(2C_I^2/(89.1(\theta_{\mathcal P}^{ST})^2)\), at least \(0.115\) at \(\kappa=1.2\) (\(\theta_{\mathcal P}^{ST}\le0.22\), \(C_I\ge0.498\)):
the row, where it applies, improves the pointwise bound by at most
one order of magnitude.
\end{proposition}

\begin{proof}
Theorems C3-V93-rows(3), C3-V72-extremal with \(\theta_{\mathcal P}^{ST}\) in place of
\(\theta_{\mathcal P}\), \(\mathcal E\le2C\), \(C=C_I^2\nu^{3/2}\); \((2/3)^{1/2}\cdot2=1.633\), \(1/1.633=0.6124\), \((3/2)^{1/2}=1.2247\);
the pointwise bound is \eqref{eq:c4v13-EP}; \(G(1.2)=0.498\) from
Theorem~\ref{thm:c4v8-region}.
\end{proof}

\begin{firewall}[Item (AC2): the palinstrophy row above \(\kappa=1\) is conditional and weak]
\label{fw:c4v21-palinstrophy}
Above \(\kappa=1\) the palinstrophy row applies only under the hypothesis
that \(C_\infty'\) lies in the window of Proposition~\ref{prop:c4v21-window}, a
hypothesis the class neither implies nor excludes; and under it the
bound \eqref{eq:c4v21-palin} is within an order of magnitude of the
pointwise bound.  No loop closes on it (Proposition~\ref{prop:c4v21-loops}).
Direction (D4) is closed with this record.  What would lift it: an
upper bound on \(C_\infty'\) above \(\kappa=1\) that does not go through the
absorption (the sup-norm route), i.e.\ a derivative bound from the
class constants \(C_I\), \(K_6\) directly, which is the far-field block's
bound \(C_\infty'\le[\dots]/(1-(8/3)^{1/2}\kappa A)\) of V7, itself subject to the same
absorption.
\end{firewall}

% =============================================================================
\section{Integrated V21 conclusion and next acquisition order}
\label{sec:c4v21-conclusion}
% =============================================================================

\begin{theorem}[What V21 establishes]
\label{thm:c4v21-conclusion}
Relative to the three legacy cycles and V1--V20: the closure of (D4)
(Proposition~\ref{prop:c4v21-loops}, Firewall~\ref{fw:c4v21-palinstrophy}),
the derivative table (Theorem~\ref{thm:c4v21-table}) and the window
(Proposition~\ref{prop:c4v21-window}).  The full Navier--Stokes
stability/global-regularity theorem remains unproved.
\end{theorem}

\begin{proof}
The cited results.
\end{proof}

\begin{programme}[V22 acquisition order]
\label{prog:c4v22}
\begin{enumerate}[label=\textup{(AC\arabic*)},leftmargin=4.8em]
\item The pointwise use of the equation on a ball through the ancient
      form, against N-C3-1: restate the negative record (Firewalls
      C3-V22-no-floor-from-concentration and C3-V23-no-equation)
      and its lift; write the ancient form at a point of the ball for a
      profile concentrated in the ball \(B_R\) (energy fraction \(1-\epsilon\)
      inside), split the kernel integral into the ball's contribution
      and the exterior's, and bound the exterior by \(\epsilon\) and the
      kernel tails of V2, V4.
\item Determine whether the ball's contribution forces a nonzero
      variance (a profile nearly constant in the ball is nearly a
      constant vector \(\bar V\), and \(\bar V\otimes\bar V\) is annihilated by the
      Leray projection of the divergence: the ancient form of a
      constant is zero, so a nearly constant profile must be nearly
      zero, against \(\kappa\ge0.3438\)); make this quantitative or record the
      negative.
\item The next compulsory companion audit is V25.
\end{enumerate}
\end{programme}

\begin{assessment}[Position after V21]
\label{ass:c4v21-position}
(D4) is closed; the (D5) block begins with N-C3-1.  No full stability
or global-regularity claim is made.
\end{assessment}

% =============================================================================
\section*{Version V21 (fourth cycle) append-only record}
\addcontentsline{toc}{section}{Version V21 (fourth cycle) append-only record}
% =============================================================================

The complete V20 source of the fourth cycle was retained before this
continuation.  Its SHA--256 digest was
\begin{center}\scriptsize\ttfamily
c57f860d202d2e8ca4de280197a5f36974dab6dcdba3882d2395087f469a56bd.
\end{center}
V21 closed direction (D4) and tabulated the derivative constants.  No
full regularity claim is made.

% =============================================================================
\section{V22: the pointwise use of the equation on a ball through the ancient form}
\label{sec:c4v22-entry}
% =============================================================================

Items (AC1) and (AC2) of Programme~\ref{prog:c4v22} are carried out.
The ancient form annihilates constant fields exactly, so a profile
that is nearly a constant in the Gaussian sense must reconstruct its
sup norm from its fluctuation; splitting the kernel integral into a
far past, a ball around the point where the sup is attained, its
exterior, and a short recent interval gives an explicit inequality
between the sup norm and the Gaussian fluctuation on the window.
Its consequence is a floor on the level variance of a nearly
constant profile, conditional on the window's length and on the
sup being nearly attained at its end, and of size \(10^{-406}\) at the
corner: an explicit constant, but exponentially small in the square
of the ball radius, which the polynomial tail of the Leray
projection forces to be about sixty Gaussian units.  This is the
structural reason for N-C3-1, now with a number, and the record is
negative in substance.  No global regularity claim is made.  The
next compulsory companion audit is V25.

% =============================================================================
\section{Item (AC1): the ancient form on a ball}
\label{sec:c4v22-ball}
% =============================================================================

Notation: \(W\) an ancient solution of the family at a homogeneous
type-I point, \(s<0\), \(L:=-s\), \(\omega(s):=(L/\nu)^{1/2}\|W(s)\|_\infty\le\kappa\); for \(\tau<0\),
\(\overline W(\tau)\) the mean of \(W(\tau)\) against the Gaussian of scale \(-\tau\) (the
Gaussian mean \(\bar V\) of the profile at \(\sigma=-\log(-\tau)\)) and
\(\widetilde W=W-\overline W\); \(\varphi\eta\ge\|\widetilde V\|_G^2\) on the window (the third cycle's V61);
\(x^*\) a point with \(|W(x^*,s)|=\|W(s)\|_\infty\) and \(\zeta:=|x^*|(\nu L)^{-1/2}\) its position in
profile units; \(m_2:=\|\,\|\nabla K^O_1\|_{s0}\|_{L^2(\mathbb R^3)}\).

\begin{lemma}[The ancient form annihilates constants; the \(L^2\) and tail norms of the kernel]
\label{lem:c4v22-constants}
\begin{enumerate}[label=\textup{(\roman*)},leftmargin=3.2em]
\item For every constant vector \(a\) and \(t>0\), \(e^{\nu t\Delta}\PP\operatorname{div}(a\otimes a)=0\); hence for
      every \(\tau\), \(\int\nabla K^O_{\nu(s-\tau)}(x-y):\bigl(\overline W\otimes\overline W\bigr)(\tau)\dd y=0\), and the ancient form
      reads
      \begin{equation}
       W(x,s)=-\int_{-\infty}^s\!\!\int\nabla K^O_{\nu(s-\tau)}(x-y):\bigl(W\otimes W-\overline W\otimes\overline W\bigr)_0(\tau,y)\dd y\dd\tau .
       \label{eq:c4v22-ancient}
      \end{equation}
\item \(m_2\le0.0676\) (certified: midpoint sum of \(4\pi r^2\|M\|_{s0}^2\) on \([0,8]\) with
      \(4\cdot10^6\) cells and the sup bound \(0.0358\) below \(r=0.01\), plus the tail
      \(54(1+\epsilon_8)^2/(20\pi\cdot8^5)\)).
\item For \(\rho\ge8(\nu t)^{1/2}\), \(\int_{|y|>\rho}\|\nabla K^O_{\nu t}(y)\|_{s0}\dd y\le\sqrt{54}(1+\epsilon_8)/\rho\), independent
      of \(t\).
\end{enumerate}
\end{lemma}

\begin{proof}
(i) \(\operatorname{div}(a\otimes a)=0\).  Then \(\nabla K^O_t=t^{-2}M(\cdot/t^{1/2})\) with \(M=\nabla K^O_1\): (ii) by the
scaling \(\|\nabla K^O_{\nu t}\|_{L^2}=(\nu t)^{-5/4}m_2\) and the tail bound \(\|M(w)\|_{s0}\le\sqrt{54}(1+\epsilon_8)/(4\pi|w|^4)\)
for \(|w|\ge8\) (Theorem~\ref{thm:c4v2-certified}); (iii)
\((\nu t)^{-2}\int_{|y|>\rho}\|M(y/(\nu t)^{1/2})\|\dd y=(\nu t)^{-1/2}\int_{|w|>\rho/(\nu t)^{1/2}}\|M\|\dd w\le(\nu t)^{-1/2}\sqrt{54}(1+\epsilon_8)(\nu t)^{1/2}/\rho\).
\end{proof}

\begin{theorem}[The sup norm against the Gaussian fluctuation on a window]
\label{thm:c4v22-window}
For every \(s<0\), \(\theta>0\), \(0<\delta<\theta L\) and \(r\ge8\theta^{1/2}\), with \(s_0:=s-\theta L\) and
\(\rho:=r(\nu L)^{1/2}\),
\begin{equation}
 \boxed{\quad
 \omega(s)\ \le\ \Bigl(\frac23\Bigr)^{1/2}c_O^{s0}\kappa^2\Bigl[(\pi-2\arctan\theta^{1/2})+2\Bigl(\frac\delta L\Bigr)^{1/2}\Bigr]
 +2\Bigl(\frac23\Bigr)^{1/2}\sqrt{54}\,\kappa^2\frac{\log(1+\theta)}r
 +8\kappa m_2e^{(\zeta+r)^2/8}\Bigl(\frac\delta L\Bigr)^{-1/4}\sup_{[s_0,s]}\Bigl(\frac{\varphi\eta}{\nu^{5/2}}\Bigr)^{1/2} ,
 \quad}
 \label{eq:c4v22-window}
\end{equation}
the four terms being the far past \((-\infty,s_0)\), the recent interval
\([s-\delta,s]\), the exterior of the ball \(B_\rho(x^*)\) on \([s_0,s-\delta]\), and the ball
on \([s_0,s-\delta]\).
\end{theorem}

\begin{proof}
Evaluate \eqref{eq:c4v22-ancient} at \(x^*\).  Since the subtracted term
integrates to zero for each \(\tau\) separately, it may be omitted on any
\(\tau\)-interval; on the far past and the recent interval it is omitted,
and the sup-norm bound of Theorem C3-V81-ancient with the traceless
reduction \(|(W\otimes W)_0|\le(2/3)^{1/2}C_\infty^2(-\tau)^{-1}\) gives
\((2/3)^{1/2}c_O^{s0}\kappa^2\nu^{1/2}\int(s-\tau)^{-1/2}(-\tau)^{-1}\dd\tau\), which
is \(L^{-1/2}(\pi-2\arctan\theta^{1/2})\) on \((-\infty,s_0)\) and at most \(2\delta^{1/2}L^{-1}\) on \([s-\delta,s]\).
On \([s_0,s-\delta]\) subtract \(\overline W\otimes\overline W\) (Lemma~\ref{lem:c4v22-constants}(i)); on the
exterior \(|(W\otimes W-\overline W\otimes\overline W)_0|\le2(2/3)^{1/2}C_\infty^2(-\tau)^{-1}\) (\(|\overline W|\le\|W\|_\infty\)) and
Lemma~\ref{lem:c4v22-constants}(iii) with \(\rho\ge8(\nu(s-\tau))^{1/2}\) for \(s-\tau\le\theta L\), then
\(\int_{s_0}^s(-\tau)^{-1}\dd\tau=\log(1+\theta)\).  On the ball, Cauchy--Schwarz with
\(\|\nabla K^O_{\nu(s-\tau)}\|_{L^2}=(\nu(s-\tau))^{-5/4}m_2\) and
\(\|W\otimes W-\overline W\otimes\overline W\|_{L^2(B_\rho(x^*))}\le2C_\infty(-\tau)^{-1/2}\|\widetilde W(\tau)\|_{L^2(B_\rho(x^*))}\); in profile
variables \(\widetilde W(y,\tau)=(-\tau)^{-1/2}\widetilde V(y(-\tau)^{-1/2})\), so
\(\|\widetilde W(\tau)\|_{L^2(B_\rho(x^*))}^2=(-\tau)^{1/2}\int_{B}|\widetilde V|^2\dd z\le(-\tau)^{1/2}e^{(\zeta+r)^2/4}\|\widetilde V\|_G^2\), since on the
ball \(|z|\le(|x^*|+\rho)(-\tau)^{-1/2}\le(\zeta+r)\nu^{1/2}\) for \(-\tau\ge L\), and \(G^{-1}\le e^{(\zeta+r)^2/4}\) there;
then \(\int_{s_0}^{s-\delta}(s-\tau)^{-5/4}(-\tau)^{-1/4}\dd\tau\le L^{-1/4}\cdot4\delta^{-1/4}\).  Collect, multiply by
\((L/\nu)^{1/2}\), and use \(C_\infty=\kappa\nu^{1/2}\).
\end{proof}

% =============================================================================
\section{Item (AC2): the conditional floor and the record}
\label{sec:c4v22-floor}
% =============================================================================

\begin{theorem}[A conditional explicit floor on the level variance of a nearly constant profile]
\label{thm:c4v22-floor}
At the corner \(\kappa=0.3438\), let \(s\) be a time at which \(\omega(s)\ge(1-\epsilon)\kappa\) and
suppose the profile is nearly constant (dominant level \(0\), level
variance \(\eta\)) throughout the window \([s_0,s]\) with \(\theta=38.5\), i.e.\ of
\(\sigma\)-length \(\Lambda=\log(1+\theta)=3.68\).  Then, with \(r=61\), \(\delta/L=0.026\),
\begin{equation}
 \sup_{[s_0,s]}\eta\ \ge\ \frac{\nu^{5/2}}{\varphi}\Bigl(\frac14-\epsilon\Bigr)^2\frac{(\delta/L)^{1/2}}{64m_2^2}e^{-(\zeta+r)^2/4}
 \ \ge\ 0.104\Bigl(\frac14-\epsilon\Bigr)^2e^{-(\zeta+61)^2/4}
 \ \ge\ 10^{-406}\quad(\zeta=0,\ \epsilon\to0).
 \label{eq:c4v22-floor}
\end{equation}
\end{theorem}

\begin{proof}
In \eqref{eq:c4v22-window} the three fluctuation-free terms are
\(0.0853\), \(0.0861\), \(0.0855\) (\((2/3)^{1/2}\cdot2.766\cdot0.1182=0.267\), \(\pi-2\arctan\sqrt{38.5}=0.319\);
\(0.267\cdot2\cdot0.161=0.0861\); \(2\cdot0.8165\cdot0.1182\cdot7.348\cdot3.676/61=0.0855\)), each at most
\(\kappa/4=0.0860\), and \(r=61\ge8\sqrt{38.5}=49.6\); so the ball term is at least
\((\frac14-\epsilon)\kappa\), i.e.\ \((\varphi\eta/\nu^{5/2})^{1/2}\ge(\frac14-\epsilon)(\delta/L)^{1/4}/(8m_2e^{(\zeta+r)^2/8})\); square, with
\(m_2\le0.0676\), \((0.026)^{1/2}=0.161\), \(0.161/(64\cdot0.00457)=0.55\), and \(\varphi\le5.27\nu^{5/2}\):
\(0.55/5.27=0.104\), and at \(\epsilon=0\), \(0.104/16=6.5\cdot10^{-3}\); \(e^{-61^2/4}=e^{-930}=10^{-404}\).
\end{proof}

\begin{firewall}[Item (AC2): N-C3-1 with the equation, the record]
\label{fw:c4v22-record}
The pointwise use of the equation on a ball lifts N-C3-1 only in the
following sense: a nearly constant profile has a level variance
bounded below by an explicit constant, \eqref{eq:c4v22-floor}, but (a)
only on windows of \(\sigma\)-length at least \(3.7\) that end where the sup
of \((-s)^{1/2}\|W\|_\infty\) is nearly attained (at other times the class bounds
\(\omega(s)\) below only through \((\varphi/|G|\nu)^{1/2}\), hence through \(c_\Phi\)); (b) with
a floor of size \(e^{-(\zeta+61)^2/4}\), which is \(10^{-406}\) at best; (c) with the
position \(\zeta\) of the sup entering, which the class does not bound.
The size is structural, not an artefact of the constants: the
exterior of the ball is controlled only by the polynomial tail
\(\sqrt{54}/\rho\) of the Oseen kernel's gradient, the nonlocality of the
Leray projection, while the fluctuation is controlled only inside
the Gaussian, so the ball must be large (\(r\ge60\)) and the Gaussian
factor \(e^{r^2/4}\) follows.  The same obstruction appears if the ancient
form is paired with a Gaussian test function instead of evaluated at
a point (the Leray projection of a Gaussian has a polynomial tail).
The record of the third cycle stands: concentration, even with the
equation, gives no usable floor on the variance.  What would lift it:
a Gaussian-weighted control of the exterior, i.e.\ a bound on
\(\int_{|y|>\rho}|\widetilde W|^2\) without the weight, which is a decay statement on the
profile that the class does not supply (the profile is bounded, not
decaying).
\end{firewall}

\begin{assessment}[Reading]
\label{ass:c4v22-reading}
The fourth cycle asked the question the third could not: does the
equation, used pointwise, force a nearly constant profile to
fluctuate?  The answer is yes in principle (a constant is not a
solution, and the ancient form says so exactly) and no in practice
(the fluctuation forced is \(10^{-406}\) of the energy at the corner, under
conditions).  The obstruction is the one the whole series has met in
different forms: the pressure is nonlocal and the Gaussian control
is local.  The (D5) block turns to N-C3-3 in V23, where the same
tool, the ancient form, is applied to the defect rather than to the
variance: a self-similar profile satisfies the ancient form with
\(Y=0\), and the question is whether the ancient form measures the
distance from self-similarity.
\end{assessment}

% =============================================================================
\section{Integrated V22 conclusion and next acquisition order}
\label{sec:c4v22-conclusion}
% =============================================================================

\begin{theorem}[What V22 establishes]
\label{thm:c4v22-conclusion}
Relative to the three legacy cycles and V1--V21: the ancient form on
a ball (Lemma~\ref{lem:c4v22-constants}, Theorem~\ref{thm:c4v22-window}),
the conditional floor (Theorem~\ref{thm:c4v22-floor}) and the record
(Firewall~\ref{fw:c4v22-record}).  The full Navier--Stokes
stability/global-regularity theorem remains unproved.
\end{theorem}

\begin{proof}
The cited results.
\end{proof}

\begin{programme}[V23 acquisition order]
\label{prog:c4v23}
\begin{enumerate}[label=\textup{(AC\arabic*)},leftmargin=4.8em]
\item The ancient form and the defect: write the defect \(Y=\partial_\sigma V\) of
      the profile in terms of the ancient form (the time derivative of
      \eqref{eq:c4v22-ancient} in the profile variables), and identify
      the self-similar profiles as the fixed points; state the
      identity for \(\|Y\|_G\) that results.
\item Determine whether the identity bounds \(\|Y\|_G\) below by a
      distance of the profile from the self-similar profiles that the
      class controls (a Liouville constant for (E9)), or record the
      negative with the structural reason; the plan of V20 is
      amended: V23 replaces the mean version of the ball argument,
      which V22 shows to be void for the same structural reason.
\item The next compulsory companion audit is V25.
\end{enumerate}
\end{programme}

\begin{assessment}[Position after V22]
\label{ass:c4v22-position}
N-C3-1 is recorded with the equation; N-C3-3 is next.  No full
stability or global-regularity claim is made.
\end{assessment}

% =============================================================================
\section*{Version V22 (fourth cycle) append-only record}
\addcontentsline{toc}{section}{Version V22 (fourth cycle) append-only record}
% =============================================================================

The complete V21 source of the fourth cycle was retained before this
continuation.  Its SHA--256 digest was
\begin{center}\scriptsize\ttfamily
9159a5cf188d664d233382a3b1ba0b64b74bf6a93eb91c756943d1cefd87fe08.
\end{center}
V22 used the equation pointwise on a ball and recorded the result.
No full regularity claim is made.

% =============================================================================
\section{V23: the ancient form in profile variables, and the time derivative as a neutral mode}
\label{sec:c4v23-entry}
% =============================================================================

Items (AC1) and (AC2) of Programme~\ref{prog:c4v23} are carried out.
In profile variables the ancient form is the variation-of-constants
formula for the profile equation with the semigroup of \(L_-\), which
is a heat flow composed with a dilation; the self-similar profiles
are its fixed points.  Differentiating in time gives a linear
homogeneous equation for the time derivative along every trajectory:
it is a neutral mode of the linearized ancient form, and the defect
is that mode minus the dilation generator.  The
equation reproduces a floor on \(\kappa\) (weaker than the certified one)
and gives no floor on the defect, for the structural reason that a
homogeneous linear equation bounds its solution only by itself; the
quantitative Liouville statement that would lift N-C3-3 needs the
decay of approximate profiles at infinity, which the class does not
supply.  N-C3-3 stands with this reason.  V22 is corrected in place
(a double-counted factor in the middle expression of its floor).  No
global regularity claim is made.  The next compulsory companion audit
is V25.

\begin{remark}[Correction to V22, made in place]
\label{rem:c4v23-correction}
The middle expression of (\ref{eq:c4v22-floor}) carried the factor
\((\frac14-\epsilon)^2\) and also the numerical factor \(\frac1{16}\) that is its value at
\(\epsilon=0\); it now reads \(0.104(\frac14-\epsilon)^2e^{-(\zeta+61)^2/4}\), with the final value \(10^{-406}\)
unchanged.  The draft and the live file were patched before this
version was appended; the digest recorded below is that of the
patched V22.
\end{remark}

% =============================================================================
\section{Item (AC1): the profile ancient form and the linearized equation}
\label{sec:c4v23-profile}
% =============================================================================

\begin{lemma}[The semigroup of \(L_-\)]
\label{lem:c4v23-semigroup}
\(L_-=\nu\Delta-\frac12z\cdot\nabla-\frac12\) generates on bounded fields
\begin{equation}
 e^{uL_-}f(z)=e^{-u/2}\bigl(e^{\nu(1-e^{-u})\Delta}f\bigr)(e^{-u/2}z),\qquad u\ge0 ,
 \label{eq:c4v23-semigroup}
\end{equation}
and \(e^{uL_-}\PP\operatorname{div}=e^{u/2}\PP\operatorname{div}\,e^{uL_-}\) (the semigroup acting componentwise on
matrices).  Consequently \(\|e^{uL_-}\PP\operatorname{div}F\|_\infty\le c_O^{s0}\nu^{-1/2}e^{-u/2}(1-e^{-u})^{-1/2}\|F_0\|_\infty\) for
symmetric \(F\), and \(\int_0^\infty e^{-u/2}(1-e^{-u})^{-1/2}\dd u=\pi\).
\end{lemma}

\begin{proof}
Differentiate the right side of \eqref{eq:c4v23-semigroup} in \(u\): with
\(g=e^{\nu(1-e^{-u})\Delta}f\) and \(a=e^{-u}\), \(\partial_u[a^{1/2}g(a^{1/2}z)]=a^{1/2}[-\frac12g+\nu a\Delta g-\frac12(z\cdot\nabla g)](a^{1/2}z)\)
and \(\nu\Delta_z[a^{1/2}g(a^{1/2}z)]=a^{3/2}\nu(\Delta g)(a^{1/2}z)\), \(z\cdot\nabla_z[a^{1/2}g(a^{1/2}z)]=a^{1/2}(z\cdot\nabla g)(a^{1/2}z)\); so
the right side satisfies \(\partial_u=L_-\), with value \(f\) at \(u=0\).  The
commutation: \(\PP\) and \(\operatorname{div}\) commute with the heat flow; under the
dilation \(f\mapsto f(a^{1/2}\cdot)\), \(\operatorname{div}\) picks up the factor \(a^{1/2}=e^{-u/2}\) and \(\PP\)
is invariant (a Fourier multiplier of degree \(0\)).  The sup bound is
Theorem C3-V86-absorptions with \(\|e^{\nu t\Delta}\PP\operatorname{div}F\|_\infty\le c_O^{s0}(\nu t)^{-1/2}\|F_0\|_\infty\) at
\(t=1-e^{-u}\).  The integral: \(a=e^{-u}\), \(\int_0^1a^{-1/2}(1-a)^{-1/2}\dd a=B(\frac12,\frac12)=\pi\).
\end{proof}

\begin{theorem}[The ancient form in profile variables]
\label{thm:c4v23-profile-ancient}
For every element of the family and every \(\sigma\),
\begin{equation}
 \boxed{\quad
 V(\sigma)=-\int_0^\infty e^{uL_-}\,N(V)(\sigma-u)\,\dd u,\qquad N(V)=\PP\operatorname{div}(V\otimes V),
 \quad}
 \label{eq:c4v23-profile-ancient}
\end{equation}
the integral converging absolutely in \(L^\infty\); this is the
variation-of-constants formula for \(\partial_\sigma V=L_-V-N(V)\) on \((-\infty,\sigma]\), the
initial term vanishing.  A profile \(V\) independent of \(\sigma\) satisfies it
iff \(V=-(-L_-)^{-1}N(V)\), i.e.\ iff \(L_-V=N(V)\): the self-similar profiles are
the fixed points of \(\mathcal A(V):=-\int_0^\infty e^{uL_-}N(V)\dd u\).
\end{theorem}

\begin{proof}
Rescale Theorem C3-V81-ancient: with \(W(\tau)=(-\tau)^{-1/2}V(\cdot(-\tau)^{-1/2},\sigma')\), \(\sigma'=-\log(-\tau)\),
\(e^{\nu(s-\tau)\Delta}\PP\operatorname{div}(W\otimes W)(\tau)(x)=(-\tau)^{-3/2}\bigl(e^{\nu(1-a)\Delta}N(V)(\sigma')\bigr)(x(-\tau)^{-1/2})\) with \(a=(-s)/(-\tau)=e^{\sigma'-\sigma}\);
evaluate at \(x=(-s)^{1/2}z\), multiply by \((-s)^{1/2}\), and use \(\dd\tau=(-\tau)\dd\sigma'\):
\(V(z,\sigma)=-\int_{-\infty}^\sigma a^{1/2}\bigl(e^{\nu(1-a)\Delta}N(V)(\sigma')\bigr)(a^{1/2}z)\dd\sigma'\), which is
\eqref{eq:c4v23-profile-ancient} by \eqref{eq:c4v23-semigroup} with \(u=\sigma-\sigma'\).
Convergence: \(\|N(V)\|\) enters through the kernel bound of Lemma~\ref{lem:c4v23-semigroup},
integrable.  The fixed-point statement: \(\int_0^\infty e^{uL_-}\dd u=(-L_-)^{-1}\) on the
range of the bounded functional calculus (spectrum of \(L_-\) in
\(\{\Re\le-\frac12\}\) on \(L^2(G)\), and \(\|e^{uL_-}\|_{\infty\to\infty}\le e^{-u/2}\)).
\end{proof}

\begin{theorem}[The time derivative is a neutral mode of the linearized ancient form]
\label{thm:c4v23-neutral}
Let \(\dot W:=\partial_sW\) and \(M_t:=\sup_s(-s)^{3/2}\nu^{-1/2}\|\dot W(s)\|_\infty\) (finite: \(M_t\le38.96\) at the
corner, Theorem~\ref{thm:c4v12-YN}).  Along every ancient solution of
the family,
\begin{equation}
 \boxed{\quad
 \dot W(s)=-\int_{-\infty}^se^{\nu(s-\tau)\Delta}\PP\operatorname{div}\bigl(W\otimes\dot W+\dot W\otimes W\bigr)(\tau)\dd\tau
 \ =:\ -\mathcal K_W\dot W(s),
 \qquad M_t\ \le\ \pi c_O^{s0}\kappa\,M_t ,
 \quad}
 \label{eq:c4v23-neutral}
\end{equation}
the integral converging absolutely in \(L^\infty\): \(\dot W\in\ker(I+\mathcal K_W)\).  Consequently:
\begin{enumerate}[label=\textup{(\roman*)},leftmargin=3.2em]
\item if \(\kappa<1/(\pi c_O^{s0})=0.1151\), then \(\dot W\equiv0\), \(W\) is a steady solution with
      \(|W|\le C_\infty(-s)^{-1/2}\to0\) as \(s\to-\infty\), so \(W\equiv0\), against \(\kappa>0\): the floor
      \(\kappa\ge0.1151\), which is the \(\varepsilon\)-regularity floor \(1/(\pi c_O)=0.1122\) of C3-V81
      recovered through the time derivative with the traceless
      constant, weaker than the certified \(0.3438\);
\item in profile variables \(\dot W=(-s)^{-3/2}\bigl[Y+\frac12(V+z\cdot\nabla V)\bigr]\): the defect is
      the neutral mode minus the dilation generator.  The defect
      itself satisfies formally the profile-linearized equation
      \(Y(\sigma)=-\int_0^\infty e^{uL_-}\PP\operatorname{div}(V\otimes Y+Y\otimes V)(\sigma-u)\dd u\), but the class does not make
      that integral converge: \(Y\) grows linearly (\(|Y(z)|\le39.14\nu^{1/2}+0.576|z|\) at
      the corner) and the gradient kernel's tail \(|\cdot|^{-4}\) against linear
      growth diverges logarithmically;
\item no lower bound on \(M_t\) or on \(\|Y\|\) in any norm follows from
      \eqref{eq:c4v23-neutral}: the equation is linear and homogeneous in
      \(\dot W\), so it bounds \(\dot W\) by \(\dot W\), and the alternative it yields is
      \(\dot W\equiv0\) or \(\pi c_O^{s0}\kappa\ge1\), a floor on \(V\), not on the mode.
\end{enumerate}
\end{theorem}

\begin{proof}
The time-translated solution \(W(\cdot+h)\) satisfies the Duhamel formula
of C3-V81 on \([s_0,s]\) with data \(W(s_0+h)\); differentiate in \(h\) at \(h=0\):
\(\dot W(s)=e^{\nu(s-s_0)\Delta}\dot W(s_0)-\int_{s_0}^se^{\nu(s-\tau)\Delta}\PP\operatorname{div}(W\otimes\dot W+\dot W\otimes W)(\tau)\dd\tau\), and
\(\|\dot W(s_0)\|_\infty\le M_t\nu^{1/2}(-s_0)^{-3/2}\to0\) as \(s_0\to-\infty\).  The bound:
\(|(W\otimes\dot W+\dot W\otimes W)_0|\le2C_\infty M_t\nu^{1/2}(-\tau)^{-2}\), the kernel \(c_O^{s0}(\nu(s-\tau))^{-1/2}\), and
\(\int_{-\infty}^s(s-\tau)^{-1/2}(-\tau)^{-2}\dd\tau=L^{-3/2}B(\frac12,\frac32)=\frac\pi2L^{-3/2}\); so
\(\|\dot W(s)\|_\infty\le2c_O^{s0}\nu^{-1/2}\kappa\nu^{1/2}M_t\nu^{1/2}\frac\pi2L^{-3/2}\), i.e.\ \(M_t\le\pi c_O^{s0}\kappa M_t\); \(\pi\cdot2.766=8.69\),
\(1/8.69=0.1151\).  (i): a steady bounded solution of the Navier--Stokes
system on \(\mathbb R^3\) that tends to zero uniformly along a sequence of
times is zero.  (ii): with \(W(x,s)=(-s)^{-1/2}V(x(-s)^{-1/2},\sigma)\), \(\sigma=-\log(-s)\), the chain
rule gives \(\partial_sW=(-s)^{-3/2}[\frac12V+\frac12z\cdot\nabla V+Y]\); the growth of \(Y\) is
Theorem~\ref{thm:c4v12-YN}; \(\int_{|y|>R}|y|^{-4}|y|\dd y=\infty\).  (iii) is the statement.
\end{proof}

% =============================================================================
\section{Item (AC2): the Liouville constant and the record}
\label{sec:c4v23-liouville}
% =============================================================================

\begin{proposition}[What a Liouville constant for (E9) in the class would be]
\label{prop:c4v23-liouville}
A quantitative (E9) in the class is a statement of the form: for every
profile \(V\) of the class with \(\|L_-V-N(V)\|_X\le\delta_0\) in some norm \(X\),
\(\|V\|\le\epsilon(\delta_0)\) with \(\epsilon(\delta_0)\to0\); its contrapositive with \(\|V\|_\infty\ge(c_\Phi/|G|)^{1/2}\) is
the floor \(\|Y\|_X\ge\delta_0\).  The proof of (E9) (Ne\v cas--R\accent23u\v zi\v cka--\v Sver\'ak; Tsai)
runs: the profile equation with \(V\in L^q\) gives decay of \(V\), \(\nabla V\) and
of the head \(\Pi=P+\frac12|V|^2+\frac12z\cdot V\) at infinity; the identity
\(\nu\Delta\Pi-(V+\frac12z)\cdot\nabla\Pi=\nu|\Omega|^2\ge0\) and the maximum principle give \(\Pi\le0\)
with equality forcing \(\Omega\equiv0\); then \(V\) harmonic and decaying is zero.
With the defect the identity is \(\nu\Delta\Pi-(V+\frac12z)\cdot\nabla\Pi=\nu|\Omega|^2+(V+\frac12z)\cdot Y\)
(Theorem C3-V31-NRS), whose source has no sign and grows linearly
in \(|z|\) (the class bounds \(|Y(z)|\le39.14\nu^{1/2}+0.576|z|\) at the corner,
Theorem~\ref{thm:c4v12-YN}), and the decay step needs \(Y\) small in an
unweighted norm at infinity, which the class does not give: the
defect is bounded below only in the mean of \(L^2(G)\) (the defect
rate), and above pointwise with linear growth.  In the ancient form
the same obstruction reads: the tail of \(e^{uL_-}\PP\operatorname{div}\) is polynomial,
so \eqref{eq:c4v23-profile-ancient} transmits decay only if \(V\otimes V\)
decays, and the class gives \(V\in L^6\) without a rate.
\end{proposition}

\begin{proof}
The cited theorems and the description of the classical proof.
\end{proof}

\begin{firewall}[Item (AC2): N-C3-3 with the ancient form, the record]
\label{fw:c4v23-record}
The ancient form does not give an explicit Liouville constant for
(E9) in the class.  Two structural reasons, one new: (a) the
time derivative is a neutral mode of the linearized ancient form and
the defect is that mode minus the dilation generator
(Theorem~\ref{thm:c4v23-neutral}), so every bound derived from the
equation is homogeneous in the mode and yields a floor on \(V\), never
on the mode or on \(Y\);
the defect rate \(\underline\delta\) is the compactness constant of the
alternative ``\(Y\equiv0\) is excluded by (E9)'' and remains so; (b) the
quantitative (E9) needs decay of approximate profiles, which the
class does not supply, as in the third cycle's record.  The Gaussian
\(L^2\) version of \(\mathcal K_V\) meets in addition the Riesz obstruction (the Leray
projection is not bounded on \(L^2(G)\); the split constants \(c_1\)--\(c_6\)
of V12 bound it at the price of the unweighted \(L^6\) norm), and gives
again an upper bound on \(Y\) through \(Y\).  What would lift it: a decay
rate for the profile at infinity from the class (an \(L^6\) modulus, or a
bound on the far-field enstrophy that the concentration of V18 gives
only inside the Gaussian), or a second, inhomogeneous equation for
\(Y\) with a source the class bounds below.  The (D5) block is closed
with this record and V22's; V24 consolidates it.
\end{firewall}

\begin{assessment}[Reading]
\label{ass:c4v23-reading}
The two lifts of (D5) fail for the same reason, seen twice: the
equation is nonlocal through the pressure, the class controls the
profile locally (Gaussian) or without rate (\(L^6\)), and every use of
the equation at infinity needs a rate.  The neutral-mode identity is
the fourth cycle's statement of N-C3-3: the time derivative is not a
source the equation feeds, it is a solution of the linearized
equation along the trajectory, the defect differs from it by the
dilation generator, and the class fixes the size of either only
through compactness.  This is worth recording as the reason the series has
no explicit defect rate in any of its four cycles.
\end{assessment}

% =============================================================================
\section{Integrated V23 conclusion and next acquisition order}
\label{sec:c4v23-conclusion}
% =============================================================================

\begin{theorem}[What V23 establishes]
\label{thm:c4v23-conclusion}
Relative to the three legacy cycles and V1--V22: the profile ancient
form (Lemma~\ref{lem:c4v23-semigroup}, Theorem~\ref{thm:c4v23-profile-ancient}),
the neutral-mode equation for the time derivative (Theorem~\ref{thm:c4v23-neutral}),
the description of the Liouville constant (Proposition~\ref{prop:c4v23-liouville})
and the record (Firewall~\ref{fw:c4v23-record}).  The full Navier--Stokes
stability/global-regularity theorem remains unproved.
\end{theorem}

\begin{proof}
The cited results.
\end{proof}

\begin{programme}[V24 acquisition order]
\label{prog:c4v24}
\begin{enumerate}[label=\textup{(AC\arabic*)},leftmargin=4.8em]
\item Consolidation of the (D5) block V21--V23: the closure of (D4),
      the two records with their structural reason, and the standing
      negative records of the fourth cycle in one table (N-C4-1: sign
      not forced; N-C4-2: no variance floor from the equation on a
      ball; N-C4-3: no Liouville constant from the ancient form).
\item The direction for V26--V30: the remaining quantitative content
      of the fourth cycle is the region \(\mathcal R\); state what would move
      its boundary (the far-field block's constants, the strain
      constant, the traceless reduction) and choose among them with
      the reason.
\item The next compulsory companion audit is V25.
\end{enumerate}
\end{programme}

\begin{assessment}[Position after V23]
\label{ass:c4v23-position}
The (D5) block is closed with two records.  No full stability or
global-regularity claim is made.
\end{assessment}

% =============================================================================
\section*{Version V23 (fourth cycle) append-only record}
\addcontentsline{toc}{section}{Version V23 (fourth cycle) append-only record}
% =============================================================================

The complete V22 source of the fourth cycle (patched as in
Remark~\ref{rem:c4v23-correction}) was retained before this
continuation.  Its SHA--256 digest was
\begin{center}\scriptsize\ttfamily
a979217d41580e3692851e46c8c674b69c45ea77d1334199123f53a7ad23c978.
\end{center}
V23 wrote the profile ancient form and the neutral-mode equation for
the time derivative and recorded N-C3-3.  No full regularity claim is made.

% =============================================================================
\section{V24: consolidation of the (D5) block, the standing records, and the direction for V26--V30}
\label{sec:c4v24-entry}
% =============================================================================

Items (AC1) and (AC2) of Programme~\ref{prog:c4v24} are carried out.
The block V21--V23 closed directions (D4) and (D5) of the third
cycle's V99 with three records, each with a structural reason; the
standing negative records of the fourth cycle are collected in one
table.  The direction for V26--V30 is the rank-one refinement of the
kernel constants: the ancient form acts on \(W\otimes W\), a rank-one
field, while every certified constant of the cycle is a supremum
over all symmetric traceless matrices, so the constants entering
the far-past terms can be lowered, and with them the floor and the
region.  V23 was corrected in place (the neutral mode is the time
derivative, not the defect, whose linearized equation the class does
not make converge).  No global regularity claim is made.  The next
compulsory companion audit is V25.

\begin{remark}[Correction to V23, made in place]
\label{rem:c4v24-correction}
As first appended, Theorem~\ref{thm:c4v23-neutral} stated the
linearized ancient form for the defect \(Y=\partial_\sigma V\) in \(L^\infty\) and drew a
floor \(0.0575\) from it; but \(Y\) is not bounded (it grows linearly, by the
dilation term \(\frac12z\cdot\nabla V\)), and the integral does not converge for
linearly growing inputs.  The theorem now concerns the time
derivative \(\dot W=\partial_sW\), which is bounded by the class, gives the floor
\(\kappa\ge1/(\pi c_O^{s0})=0.1151\), and relates to the defect by
\(\dot W=(-s)^{-3/2}[Y+\frac12(V+z\cdot\nabla V)]\); the entry paragraph, the record and
the reading were adjusted accordingly.  The draft and the live file
were patched before this version was appended; the digest recorded
below is that of the patched V23.
\end{remark}

% =============================================================================
\section{Item (AC1): the (D5) block in one statement, and the standing records}
\label{sec:c4v24-block}
% =============================================================================

\begin{theorem}[The block V21--V23]
\label{thm:c4v24-block}
At every homogeneous type-I point with \((\kappa,C_I)\in\mathcal R\):
\begin{enumerate}[label=\textup{(\roman*)},leftmargin=3.2em]
\item the loop constraints of the third cycle are identities for
      \(\kappa\ge0.5774\) and void below (Proposition~\ref{prop:c4v21-loops}); the
      derivative constants are \(C_\infty'\le\gamma_1^{\rm anc}\kappa\) and \(\nu^{1/2}C_\infty''\le\gamma_2^{\rm anc}\) with the
      values of Theorem~\ref{thm:c4v21-table}, sharp only at the floor;
      the palinstrophy row applies above \(\kappa=1\) only if \(C_\infty'\) lies in the
      window \([1.2247(1-\kappa^2/2),\,0.6124(\frac32-\kappa^2/2))\), which the class does
      not decide (Proposition~\ref{prop:c4v21-window});
\item a nearly constant profile on a window of \(\sigma\)-length \(3.7\) ending
      where the sup is nearly attained has level variance at least
      \(0.104(\frac14-\epsilon)^2e^{-(\zeta+61)^2/4}\nu^{5/2}/\varphi\), i.e.\ \(10^{-406}\) at best
      (Theorem~\ref{thm:c4v22-floor});
\item the time derivative satisfies the linearized ancient form
      \eqref{eq:c4v23-neutral}, whence \(\kappa\ge0.1151\), and no floor on it or on the
      defect (Theorem~\ref{thm:c4v23-neutral}); the profile ancient form
      \eqref{eq:c4v23-profile-ancient} has the self-similar profiles as
      fixed points.
\end{enumerate}
\end{theorem}

\begin{proof}
The cited results.
\end{proof}

\begin{theorem}[Standing negative records of the fourth cycle]
\label{thm:c4v24-records}
\medskip
\begin{center}\small
\begin{tabular}{@{}p{0.1\textwidth}p{0.3\textwidth}p{0.28\textwidth}p{0.26\textwidth}@{}}
\toprule
\textbf{Record} & \textbf{Statement} & \textbf{Reason} & \textbf{What would lift it}\\
\midrule
N-C4-0 (Firewalls~\ref{fw:c4v4-no-second-split}, \ref{fw:c4v8-no-split}) &
 the second derivative admits no second split of the window &
 the third derivative on the kernel is not integrable at the window's end &
 a bound on \(\nabla^2(W\otimes W)\) in a weaker norm\\[3pt]
N-C4-1 (Firewall~\ref{fw:c4v17-sign}) &
 the sign of the Lamb pairing and of the pressure work at an anchor-alone extremum is not forced anywhere in \(\mathcal R\) &
 the normal component of the defect and the radial defect's direction are free in the class &
 a sign from the equation for the defect's normal component; a floor on the radial defect's mass\\[3pt]
N-C4-2 (Firewall~\ref{fw:c4v21-palinstrophy}) &
 the palinstrophy row above \(\kappa=1\) is conditional and within an order of magnitude of the pointwise bound &
 the class does not place \(C_\infty'\) in the window of applicability &
 a derivative bound above \(\kappa=1\) not through the absorption\\[3pt]
N-C4-3 (Firewall~\ref{fw:c4v22-record}) &
 the equation used pointwise on a ball gives no usable variance floor &
 the Leray projection's polynomial tail against the Gaussian locality of the fluctuation forces a ball of radius \(60\) &
 unweighted control of the exterior fluctuation (a decay rate)\\[3pt]
N-C4-4 (Firewall~\ref{fw:c4v23-record}) &
 the ancient form gives no Liouville constant for (E9) in the class &
 the time derivative is a neutral mode (homogeneous equation); decay of approximate profiles is not supplied &
 a decay rate from the class, or an inhomogeneous equation for the defect with a signed source\\
\bottomrule
\end{tabular}
\end{center}
\medskip
Records N-C4-3 and N-C4-4 are the fourth cycle's forms of N-C3-1 and
N-C3-3, each now with a structural reason; N-C3-2 was closed by the
third cycle.
\end{theorem}

\begin{proof}
The cited firewalls.
\end{proof}

% =============================================================================
\section{Item (AC2): what would move the boundary of the region, and the direction}
\label{sec:c4v24-direction}
% =============================================================================

\begin{proposition}[The constants that fix the region]
\label{prop:c4v24-constants}
The floor \(\kappa\ge0.3438\) and the curve \(F(C_I)\) are determined by: the enstrophy
constraint in strain--transport form \((2/3)^{1/2}C_\infty'+\kappa^2/2\ge1\) (C3-V93), whose
constant \((2/3)^{1/2}\) is the sharp bound \(\lambda_{\max}(S)\le(2/3)^{1/2}|S|_F\) for traceless symmetric
\(S\); the gradient bound of V3 and the split bound of V7, whose
constants are \(c_{O,2}\) (far past), \(c_O^{s0}\) and \(c_O^{s0}(R)\) (window, absorption),
\(m_3(R)\) (far field), \(S_3\) (Sobolev, sharp); and the traceless reduction
\(|(W\otimes W)_0|\le(2/3)^{1/2}|W|^2\), sharp for rank one.  Of these, only the kernel
constants are not sharp for the input they receive: they are
suprema of the pointwise norms \(\|\nabla K^O_1(w)\|_{s0}\), \(\|\nabla^2K^O_1(w)\|_{s0}\) over all
symmetric traceless matrices, while the far-past terms receive
\((W\otimes W)_0\), the traceless part of a rank-one matrix, at each point.
Replacing the suprema by the suprema over rank-one inputs,
\(\|M(w)\|_{r1}:=\sup_{|a|=1}|M(w):(a\otimes a)_0|/|(a\otimes a)_0|_F\le\|M(w)\|_{s0}\), lowers \(c_{O,2}\) in the numerator of
\(\gamma_1^{\rm anc}\) and \(c_{O,3}\) in \(\gamma_2^{\rm anc}\), and \(c_O^{s0}\) in the \(\varepsilon\)-regularity floor of
C3-V81 and in Theorem~\ref{thm:c4v23-neutral}; it does not lower the
absorption constants, whose inputs \(\nabla W\otimes W+W\otimes\nabla W\) are not rank one.
For a fixed traceless symmetric bilinear form \(B\) the two suprema are
the spectral radius and the Frobenius norm, with ratio between \(2^{-1/2}\)
and \((2/3)^{1/2}\); for the vector-valued kernel the ratio is to be
computed.
\end{proposition}

\begin{proof}
Inspection of the cited proofs; for the ratio, \(\rho(B)\le|B|_F\) with equality
impossible for traceless \(B\ne0\), the extreme cases being eigenvalues
\((1,-1,0)\) and \((2,-1,-1)\).
\end{proof}

\begin{assessment}[Direction for V26--V30: the rank-one refinement]
\label{ass:c4v24-direction}
The block V26--V30 computes the rank-one constants \(c_{O,2}^{r1}\), \(c_{O,3}^{r1}\), \(c_O^{r1}\)
by the method of V2 and V4 (radial form of the kernel derivatives,
exact pointwise norm over rank-one inputs, tail, Lipschitz constant,
midpoint sum), re-derives the gradient and second-derivative bounds
with them, and recomputes the floor, the curve \(F(C_I)\), the region, the
ledger rows that depend on \(C_\infty'\) and \(C_\infty''\), and the derivative table.
V26: the pointwise rank-one norms of \(\nabla K^O_1\) and \(\nabla^2K^O_1\), with the
channel structure of V2; V27: the certified constants and the new
floor; V28: the region and the ledger with the new values; V29:
consolidation; V30: audit.  The reason for this choice over the
alternatives: the far-field constants \(m_3(R)\) are already sharp for
their input (a Hölder bound), the strain constant is sharp, and the
sup-norm route's structure is fixed by the absorption, so the
rank-one refinement is the only certified constant of the cycle
with slack, and it enters the floor quadratically in \(\kappa\) through
\(c_{O,2}\).  The expected gain is a few percent to some tens of percent on
\(c_{O,2}\), and correspondingly on the floor; it is recorded honestly
that this moves the region, not the conclusion.
\end{assessment}

% =============================================================================
\section{Integrated V24 conclusion and next acquisition order}
\label{sec:c4v24-conclusion}
% =============================================================================

\begin{theorem}[What V24 establishes]
\label{thm:c4v24-conclusion}
Relative to the three legacy cycles and V1--V23: the block in one
statement (Theorem~\ref{thm:c4v24-block}), the standing records
(Theorem~\ref{thm:c4v24-records}), the constants that fix the region
(Proposition~\ref{prop:c4v24-constants}) and the direction.  The full
Navier--Stokes stability/global-regularity theorem remains
unproved.
\end{theorem}

\begin{proof}
The cited results.
\end{proof}

\begin{programme}[V25 acquisition order]
\label{prog:c4v25}
\begin{enumerate}[label=\textup{(AC\arabic*)},leftmargin=4.8em]
\item The compulsory companion audit: read the current SM and YM
      notes and record their digests and decision.
\item Record the position after twenty-five versions; delivery.
\end{enumerate}
\end{programme}

\begin{assessment}[Position after V24]
\label{ass:c4v24-position}
Directions (D1)--(D5) of the third cycle are all treated; the fourth
cycle's own direction begins at V26.  No full stability or
global-regularity claim is made.
\end{assessment}

% =============================================================================
\section*{Version V24 (fourth cycle) append-only record}
\addcontentsline{toc}{section}{Version V24 (fourth cycle) append-only record}
% =============================================================================

The complete V23 source of the fourth cycle (patched as in
Remark~\ref{rem:c4v24-correction}) was retained before this
continuation.  Its SHA--256 digest was
\begin{center}\scriptsize\ttfamily
bee7c2e5915753ec5dcaec973e49a7016fb49cca0a2298a488c6648c2a18525c.
\end{center}
V24 consolidated the (D5) block and fixed the direction.  No full
regularity claim is made.

% =============================================================================
\section{V25: compulsory companion audit and the position after twenty-five versions}
\label{sec:c4v25-entry}
% =============================================================================

This is the twenty-fifth version of the fourth cycle.  The compulsory
review of the two companion programmes is performed first, and the
position after twenty-five versions is recorded.  No global
regularity claim is made.  The next compulsory companion audit is
V30.

% =============================================================================
\section{Compulsory V25 review of the current companion notes}
\label{sec:c4v25-companion-audit}
% =============================================================================

The current companion files in \texttt{latex\_notes} were inspected
read-only.  Since the V20 audit the SM programme has moved from V58
to V61 of its seventh series (the V58 file is no longer in the folder;
its digest stands recorded at \eqref{eq:c4v15-SM-checkpoint}), and the
YM programme from V34 to V35 of its sixth series.  The frozen
checkpoints are
\begin{align}
 &\texttt{Renewal\_SM\_Einstein\_Continuum\_Research\_Notes\_V61.tex},
 \notag\\[-2pt]
 &\qquad
 \texttt{103e39c5b66373db7d0efeebe6d478b403af9b1b634c50f426e618f5263a854a},
 \label{eq:c4v25-SM-checkpoint}\\
 &\texttt{Renewal\_Yang\_Mills\_Mass\_Gap\_Research\_Notes\_V35.tex},
 \notag\\[-2pt]
 &\qquad
 \texttt{12e332699e7819595ab4550e9457796d0a8700d2fb9e4d2df63c9ef6c7551909}.
 \label{eq:c4v25-YM-checkpoint}
\end{align}
(SM V61 has 24431 lines; YM V35 has 15944 lines.  The folder also
holds YM V32, V33, V34 with the digests recorded at V15 and V20,
unchanged.)

\emph{SM V59--V61.}  The SM programme computed the flat conformal
Hessian of an \(R^2\)-completed Einstein action, gave its quadratic
positivity condition, proved that degree-two inverse polynomial
covariances fail scalar reflection positivity, and formulated four
admissible branches for its continuum construction; its claim
boundary records that the continuum is not constructed.

\emph{YM V35.}  The YM programme performed its own compulsory audit,
derived a root-growth requirement for its global chart, an
oscillation bound for a generated four-loop response numerator, a
response-bracket slope floor, a variable-clock spectral theorem with
an ultralocality obstruction, and a conditional closure of its
continuum mass handoff; its claim boundary records that the mass gap
is not proved.

\begin{theorem}[V25 companion-audit decision]
\label{thm:c4v25-companion-decision}
Nothing is imported from SM V61 or YM V35.  Neither bears on the
fourth cycle's blocks.
\end{theorem}

\begin{proof}
The SM results concern positivity of covariances of a gravitational
action and the YM results the response matching and spectral
structure of a lattice gauge theory; none is a Navier--Stokes
statement, and none concerns the ancient form, the kernel constants,
the region, the ledger, the gate block or the records.  The YM
programme's own audit of these notes is a reading, not an input.
\end{proof}

% =============================================================================
\section{Position after twenty-five versions}
\label{sec:c4v25-position}
% =============================================================================

\begin{assessment}[Position after twenty-five versions of the fourth cycle]
\label{ass:c4v25-position-twentyfive}
Five blocks are written: the ancient-form block (V2--V4), the
far-field block (V6--V8), the ledger with values (V11--V14), the gate
block (V16--V19) and the (D4)/(D5) block (V21--V24).  The certified
description of a homogeneous type-I point is the region \(\mathcal R\) with the
table of Theorem~\ref{thm:c4v14-table}; the gate block bounds the
defect at anchor-alone extrema; the standing negative records are
Theorem~\ref{thm:c4v24-records}.  All five directions of the third
cycle's V99 are treated.  The block V26--V30 is the rank-one
refinement of the kernel constants (Assessment~\ref{ass:c4v24-direction}):
V26 the pointwise rank-one norms of \(\nabla K^O_1\) and \(\nabla^2K^O_1\); V27 the
certified constants and the new floor; V28 the region and the
ledger with the new values; V29 consolidation; V30 audit and
delivery.  No global regularity claim is made.
\end{assessment}

% =============================================================================
\section{Integrated V25 conclusion and next acquisition order}
\label{sec:c4v25-conclusion}
% =============================================================================

\begin{theorem}[What V25 establishes]
\label{thm:c4v25-conclusion}
Relative to the three legacy cycles and V1--V24: the companion-audit
decision (Theorem~\ref{thm:c4v25-companion-decision}) and the position.
The full Navier--Stokes stability/global-regularity theorem remains
unproved.
\end{theorem}

\begin{proof}
The cited records.
\end{proof}

\begin{programme}[V26 acquisition order]
\label{prog:c4v26}
\begin{enumerate}[label=\textup{(AC\arabic*)},leftmargin=4.8em]
\item The pointwise rank-one norm of \(\nabla K^O_1\): with the radial tensor of
      C3-V82/V89 (axis \(n\), coefficients \(c\), \(h_K\), \(\varphi'''\)), compute
      \(\|M(w)\|_{r1}=\sup_{|a|=1}|M(w):(a\otimes a)|/(2/3)^{1/2}\) by reduction to the polar angle of
      \(a\) against \(n\), in closed form if possible, else certified
      numerically; compare with \(\|M(w)\|_{s0}\) and give the ratio as a
      function of \(|w|\).
\item The same for \(\nabla^2K^O_1\) with the three-channel form of V2
      (\(A\,n^4+e\,\Sigma_6+(c/r)\Sigma_3\)); the rank-one norm of the 4-tensor over
      inputs \(a\otimes a\) with the output measured in Frobenius norm over
      the two free indices.
\item The next compulsory companion audit is V30.
\end{enumerate}
\end{programme}

\begin{assessment}[Position after V25]
\label{ass:c4v25-position}
The companion audit is recorded; the rank-one block begins.  No full
stability or global-regularity claim is made.
\end{assessment}

% =============================================================================
\section*{Version V25 (fourth cycle) append-only record}
\addcontentsline{toc}{section}{Version V25 (fourth cycle) append-only record}
% =============================================================================

The complete V24 source of the fourth cycle was retained before this
continuation.  Its SHA--256 digest was
\begin{center}\scriptsize\ttfamily
db877c141d81260bd8a70a436b5b77e0aecd42df5c2c57bbaa68bf622a059213.
\end{center}
V25 performed the compulsory companion audit and recorded the
position.  No full regularity claim is made.

% =============================================================================
\section{V26: the pointwise rank-one norms of the kernel derivatives}
\label{sec:c4v26-entry}
% =============================================================================

Items (AC1) and (AC2) of Programme~\ref{prog:c4v26} are carried out.
For the first derivative of the Oseen kernel the rank-one norm
coincides with the norm over all symmetric traceless matrices at
every radius: the axial channel dominates, and it is realized by a
rank-one input.  For the second derivative the rank-one norm is
strictly smaller on the range \(0.3<|w|<3.6\), by up to eleven percent
near \(|w|=2\), and equal elsewhere; both norms have closed forms, the
squared rank-one norm being a quadratic polynomial in \(\mu^2\), \(\mu\) the
cosine of the angle between the input direction and the axis, so
the supremum is exact.  No global regularity claim is made.  The
next compulsory companion audit is V30.

% =============================================================================
\section{Item (AC1): the first derivative}
\label{sec:c4v26-first}
% =============================================================================

Notation of C3-V82, V89 and V2: \(w\in\mathbb R^3\), \(r=|w|\), \(n=w/r\); \(K_1=(4\pi)^{-3/2}e^{-r^2/4}\),
\(h_K=(r/2)K_1\); \(\Phi_1=\operatorname{erf}(r/2)/(4\pi r)\) with radial derivatives \(\varphi',\varphi'',\varphi''',\varphi''''\);
\(c=(\varphi''-\varphi'/r)/r\), \(e=(\varphi'''-3c)/r\), \(A=\varphi''''-6e-3c/r\); \(M=\nabla K^O_1\) the 3-tensor
\(M_{ijk}=\partial_kK^O_{1,ij}\), \(M_2=\nabla^2K^O_1\) the 4-tensor.  For a unit vector \(a\) put \(\mu=n\cdot a\)
and \(x=\mu^2\in[0,1]\).  The rank-one norms are
\(\|M\|_{r1}:=(3/2)^{1/2}\sup_{|a|=1}|M:(a\otimes a)|\), \(\|M_2\|_{r1}:=(3/2)^{1/2}\sup_{|a|=1}|M_2:(a\otimes a)|_F\), so that
\(|M:(W\otimes W)_0|\le\|M\|_{r1}(2/3)^{1/2}|W|^2\) pointwise, as in the far-past terms of V3
and V4 with \(\|\cdot\|_{s0}\).

\begin{proposition}[The first derivative on a rank-one input]
\label{prop:c4v26-first}
With \(u:=2c-h_K\), for \(|a|=1\),
\begin{equation}
 M:(a\otimes a)=\bigl[c-(u+3c)\mu^2\bigr]n+u\mu\,a,\qquad
 |M:(a\otimes a)|^2=c^2+(u^2-6c^2)x+(9c^2-u^2)x^2 .
 \label{eq:c4v26-first}
\end{equation}
Hence \(\|M\|_{r1}=(3/2)^{1/2}\max\{2|c|,\ (c^2+(u^2-6c^2)^2/(4(u^2-9c^2)))^{1/2}\ \text{if }u^2>12c^2\}\), and on the
whole half-line \(\|M\|_{r1}=\sqrt6|c|=\|M\|_{s0}\): the second alternative never
wins, and the axial channel \(\sqrt6|c|\) of C3-V89, realized by \(a=n\), is the
norm over all symmetric traceless inputs as well.  Consequently
\(c_O^{r1}=c_O^{s0}\) and nothing is gained at the first derivative.
\end{proposition}

\begin{proof}
\(M_{ijk}=(\varphi'''-3c)n_in_jn_k+c(\delta_{ik}n_j+\delta_{jk}n_i+\delta_{ij}n_k)-h_Kn_k\delta_{ij}\) (C3-V82) contracted
with \(a_ja_k\): \((\varphi'''-3c)\mu^2n+c(\mu a+n+\mu a)-h_K\mu a\), and \(\varphi'''-h_K=-2c\) (the exact
identity of C3-V89) gives \(\varphi'''-3c=-(u+3c)\).  Expand the square with
\(n\cdot a=\mu\); the quadratic \(f(x)=c^2+px+qx^2\), \(p=u^2-6c^2\), \(q=9c^2-u^2\), has
\(f(0)=c^2\), \(f(1)=4c^2\), and an interior maximum iff \(q<0\) and \(0<-p/(2q)<1\),
i.e.\ \(u^2>12c^2\), of value \(c^2+p^2/(4|q|)\).  The comparison with \(\|M\|_{s0}=\max\{\sqrt6|c|,|u|/\sqrt2\}\)
is the inequality \(|u|\le2\sqrt3|c|\) on the half-line, verified on the
certification grid of V27 (\(2\cdot10^7\) cells on \([0,8]\), where both channels are
explicit; below \(r=0.01\) both norms are replaced by the common sup bound
\(0.0358\) of C3-V89, so the certified integrals coincide) and by the
tail forms (\(c\sim3/(4\pi r^4)\), \(u\sim6/(4\pi r^4)\) as \(r\to\infty\), since \(h_K\) is exponentially
small, ratio \(2<2\sqrt3\)); the closed forms
were also checked against the direct maximization over \(a\) on a grid
of \(2\cdot10^4\) angles at \(r=0.3,1,2,3,5,8\), agreeing to seven digits.
\end{proof}

% =============================================================================
\section{Item (AC2): the second derivative}
\label{sec:c4v26-second}
% =============================================================================

\begin{proposition}[The second derivative on a rank-one input]
\label{prop:c4v26-second}
In the orthonormal frame \((n,t,b)\) with \(a=\mu n+st\), \(s=(1-x)^{1/2}\), the matrix
\(B_{im}=\sum_{l,j}(M_2)_{ilmj}a_la_j\) is
\begin{align}
 B_{nn}&=(e+c/r)+\Bigl(A+5e+2\frac cr+K_1\bigl(\tfrac{r^2}4-\tfrac12\bigr)\Bigr)x,\qquad
 B_{nt}=\Bigl(2e+2\frac cr-\frac{K_1}2\Bigr)\mu s,\qquad
 B_{tn}=\Bigl(2e+2\frac cr+K_1\frac{r^2}4-\frac{K_1}2\Bigr)\mu s,\notag\\
 B_{tt}&=\Bigl(3\frac cr-\frac{K_1}2\Bigr)+\Bigl(e-2\frac cr+\frac{K_1}2\Bigr)x,\qquad
 B_{bb}=\frac cr+e\,x,\qquad\text{other entries }0,
 \label{eq:c4v26-entries}
\end{align}
so that \(|B|_F^2=F_0+F_1x+F_2x^2\) is a quadratic polynomial in \(x\) with explicit
radial coefficients, and \(\|M_2\|_{r1}=(3/2)^{1/2}\max_{[0,1]}(F_0+F_1x+F_2x^2)^{1/2}\) is attained at
\(x=0\), \(x=1\) or the vertex.  The ratio \(\|M_2\|_{r1}/\|M_2\|_{s0}\) is \(1\) for \(r\le0.1\) and for
\(r\ge3.6\), and \(0.985\), \(0.961\), \(0.947\), \(0.890\), \(0.972\) at \(r=0.5,1,1.5,2,3\).
\end{proposition}

\begin{proof}
\((M_2)_{ilmj}=A\,n_in_ln_mn_j+e\,\Sigma_6+(c/r)\Sigma_3+K_1(\frac{r^2}4n_mn_j-\frac12\delta_{mj})\delta_{il}\) (V2, with
\((l,j)\) the input and \((i,m)\) the output indices); contract with \(a_la_j\):
\(B=A\mu^2\,nn^{\!\top}+e\bigl[2\mu(an^{\!\top}+na^{\!\top})+nn^{\!\top}+\mu^2I\bigr]+\frac cr\bigl[I+2aa^{\!\top}\bigr]+K_1\bigl[\frac{r^2}4\mu\,an^{\!\top}-\frac12aa^{\!\top}\bigr]\),
and in the frame \(an^{\!\top}=\mu E_{nn}+sE_{tn}\), \(aa^{\!\top}=\mu^2E_{nn}+\mu s(E_{nt}+E_{tn})+s^2E_{tt}\).  The
entries \eqref{eq:c4v26-entries} follow; the odd entries carry \(\mu s\) and the
even ones are affine in \(x\), so \(|B|_F^2\) is quadratic in \(x\) (\((\mu s)^2=x(1-x)\)).
The formula was checked against the analytic 4-tensor of V2 by
direct contraction and maximization over \(2\cdot10^3\) angles at
\(r=0.3,1,2,3,5,8,12\) (agreement to six digits) and against the
three-channel norm of V2 for \(\|M_2\|_{s0}\); the ratios by evaluation.
\end{proof}

\begin{assessment}[Reading]
\label{ass:c4v26-reading}
The rank-one refinement is empty for the first derivative and worth
about two percent in the integral for the second (V27), the gain
being confined to \(1<r<3\), where the mixed channel of V2 dominates the
norm over all inputs while a rank-one input cannot realize it.  This
is smaller than hoped in Assessment~\ref{ass:c4v24-direction} and is
recorded as such; the block is completed to keep the certified
constants sharp for their inputs, and V27 gives the floor it yields.
\end{assessment}

% =============================================================================
\section{Integrated V26 conclusion and next acquisition order}
\label{sec:c4v26-conclusion}
% =============================================================================

\begin{theorem}[What V26 establishes]
\label{thm:c4v26-conclusion}
Relative to the three legacy cycles and V1--V25: the rank-one norms
in closed form (Propositions~\ref{prop:c4v26-first}, \ref{prop:c4v26-second}).
The full Navier--Stokes stability/global-regularity theorem remains
unproved.
\end{theorem}

\begin{proof}
The cited results.
\end{proof}

\begin{programme}[V27 acquisition order]
\label{prog:c4v27}
\begin{enumerate}[label=\textup{(AC\arabic*)},leftmargin=4.8em]
\item The certified constant \(c_{O,2}^{r1}=\int\|M_2\|_{r1}\) by the midpoint sum of V2
      (\(2\cdot10^6\) cells on \([0,12]\), the Lipschitz error, the tail \(36/(4\pi r^5)\), which
      bounds the rank-one norm too), and the rank-one constant
      \(c_{O,3}^{r1}\) of the third derivative by the method of V4 with a
      certified angular maximization.
\item The gradient bound and the floor with \(c_{O,2}^{r1}\) in the far-past
      term; the second-derivative bound with \(c_{O,3}^{r1}\); the values.
\item The next compulsory companion audit is V30.
\end{enumerate}
\end{programme}

\begin{assessment}[Position after V26]
\label{ass:c4v26-position}
The pointwise rank-one norms are in closed form.  No full stability
or global-regularity claim is made.
\end{assessment}

% =============================================================================
\section*{Version V26 (fourth cycle) append-only record}
\addcontentsline{toc}{section}{Version V26 (fourth cycle) append-only record}
% =============================================================================

The complete V25 source of the fourth cycle was retained before this
continuation.  Its SHA--256 digest was
\begin{center}\scriptsize\ttfamily
9271e0bd31f9a78b4f466d1bee92c6c672ccfda1fe0e818aac02628f4197a9d3.
\end{center}
V26 computed the pointwise rank-one norms.  No full regularity claim
is made.

\end{document}
